% Official author native source lines 139--528
\section{Introduction}

The main theme of this paper and its companion~\cite{oim_copain} is the Landau levels of compact manifolds. For a physicist, the Landau quantization refers to a charged particle confined to two dimensions and exposed to a magnetic field. It has discrete energy levels connected by ladder operators. Besides the planar geometry considered by Landau~\cite{La30}, the case of Riemann surfaces has been investigated in the context of the quantum Hall effect, see~\cite{IeLi94} and references therein.

On a mathematical point of view, a natural generalisation is the Bochner Laplacian acting on the sections of a Hermitian line bundle $L$ on a compact manifold. This Laplacian is defined from two data: a Riemannian metric of the base and a connection of the line bundle. The idea underlying this work is that when the curvature of the connection is non-degenerate and large with respect to the metric, the spectrum of the Laplacian exhibits a structure similar to the Landau quantization.

More specifically, let us assume that the curvature is related to the Riemannian metric by a complex structure, and consider the spectrum of the Laplacian of $L^{k}$ in the large $k$ limit. In this setting, Faure--Tsujii~\cite{FaTs15} have shown that the eigenvalues are grouped in clusters, each of them representing a generalised Landau level. The first level was previously identified by Guillemin--Uribe~\cite{GuUr88} and studied further by Borthwick--Uribe~\cite{BoUr96} as a generalization of K\"ahler quantization. In particular, its dimension is given by a Riemann--Roch number and it comes with an algebra of Toeplitz operators quantizing the classical Poisson algebra.

Our goal in this paper is to extend these results to the higher Landau levels. Our main results are:
\begin{enumerate}
\item the dimension of the $m^{\rm th}$ Landau level is the Riemann--Roch number of $L^k \otimes F_m$ when $k$ is sufficiently large, where $F_m$ is the symmetric $m^{\rm th}$ power of the complex tangent bundle of the base.
\item there is an algebra of Berezin--Toeplitz operators associated to the $m^{\rm th}$ Landau level, the symbol of these operators being sections of the endomorphism bundle $\op{End} F_m$.
\item the $m^{\rm th}$ Landau level is isomorphic with the first Landau level twisted by $F_m$ through a ladder operator, these isomorphisms are compatible with the Berezin--Toeplitz operators.
\end{enumerate}

The main ingredient to establish these results is an asymptotic expansion of the Schwartz kernel of the spectral projector of each level. For the first level, when the complex structure is integrable, the K\"ahler case, this kernel is the Szeg\"o kernel. Its asymptotic is well-understood since the seminal work by Boutet de Monvel and Sj\"ostrand~\cite{BoSj} and has been used in numerous papers starting from~\cite{BoGu,BoMeSc,Ze}. In the non-K\"ahler case, the asymptotics of the first level projector kernel has been obtained by Borthwick--Uribe~\cite{BoUr07} and Ma--Marinescu~\cite{MaMa08}. For the higher Landau levels, this asymptotic expansion will be proved in our second paper~\cite{oim_copain}.

In the current paper, we will rely on this asymptotic expansion or more generally we will show that the previous results hold for higher Landau levels defined as the image of any projector whose Schwartz kernel has the convenient asymptotics. Here the inspiration is the generalised Toeplitz structure of Boutet de Monvel--Guillemin~\cite{BoGu} and our previous work~\cite{oim}, the idea being that the only important feature of the Landau levels is this asymptotic expansion.

The main tool we will use for the proofs is a particular class of operators containing the Landau level projectors, the associated Toeplitz operators and also the generalised ladder operators. The operators in this class are controlled at first order by their symbols, which are defined as sections of a bundle of non-commutative algebras. Each of these algebras is generated by the spectral projectors and ladder operators of a Landau Hamiltonian. By this mechanism, the basic properties of the Landau quantization are transferred to the Bochner Laplacian.


To finish this general introduction, let us mention the two contemporaneous papers~\cite{Yuri_1,Yuri_2} by Yuri Kordyukov on the same subject,  which contain some results on Berezin--Toeplitz operators common with ours, cf. Remark~\ref{sec:comp_Yuri} for a comparison. However, the computation of the dimension of the $m^{\rm th}$ Landau level, the ladder operators and the general operator algebras we use, do not appear in any other work.  Let us mention as well that in a related but different context, belonging to homogeneous microlocal analysis instead of semi-classical analysis, Boutet de Monvel--Guillemin~\cite[Chapter~15]{BoGu} and Epstein--Melrose~\cite[Chapter~6]{EM} have considered generalised Szeg\"o projections at higher level with associated Toeplitz algebras, which are similar to our constructions.


\subsection{Magnetic Laplacian}

\subsubsection{Constant magnetic intensity}
Consider a Riemannian manifold $(M,g)$ with a Hermitian line bundle $L $ equipped with a connection $\nabla$. Associated to these data is a Laplacian $\frac{1}{2} \nabla^*
\nabla$ acting on $\Ci (M, L)$, which from the physical point of view is a Schr\"odinger operator with a magnetic field $\Om = i \op{curv} (\nabla) \in
\Om^2 (M, \R)$.

We will assume that $\Om$ is non-degenerate at each point and has a constant magnetic intensity with respect to $g$ in the following sense. In the case where $M$ is a surface, the magnetic intensity is the positive function defined by $ |\Om| = B \op{vol}_g$, where $\op{vol}_g $ is the Riemannian volume, and we merely assume that $B$ is constant. In higher dimension, $\Om$ being non-degenerate, the dimension of $M$ is even, say $2n$. At any $p \in M$, there exists a skew-symmetric endomorphism $j_B (p)$ of $(T_pM, g_p)$ such that $\Om_p (X,Y) = g_p (j_B(p) X, Y)$. The eigenvalues of $j_B(p)$ are $ \pm i B_{\ell} (p)$ with $0 < B_1(p) \leqslant \ldots \leqslant B_n (p)$. We assume that these eigenvalues are all equal, $B_1 = \ldots = B_{n}$, and do not depend on $p$. Equivalently $j_B(p) = B j(p)$ with $B$ a positive constant and $j$ an almost complex structure of $M$ compatible with $g$, cf.~\cite[Proposition~2.5.6]{McSa}.

So we have that $\Om = B \om$ where $B >0$ is constant and $\om$ is a symplectic form of $M$ defined by $\om (X,Y) = g (jX, Y)$. We will consider the large $B$ limit. To do this, we will replace $L$ by $L^k$, $k \in \N$, so that the curvature of $\nabla^{L^k}$ is $k B \om$, and let $k$ tend to infinity. We will also normalise the metric so that $B=1$, and our magnetic intensity is simply $k$.

Alternatively, we can introduce our data as follows. Consider a compact symplectic manifold $(M^{2n}, \om)$ with a compatible almost complex structure $j$ and a Hermitian line bundle $L \rightarrow M$ with a connection $\nabla$ having curvature $\frac{1}{i} \om$. Such a bundle is called a {\em prequantum} bundle in the Kostant--Souriau theory, where it is used to define the geometric quantization of $M$. For any positive integer $k$, we consider the Laplacian
\begin{gather}
\Delta_k = \tfrac{1}{2}\left(\nabla^{L^k}\right)^* \nabla^{L^k} : \Ci (M, L^k)\rightarrow \Ci \left(M, L^k\right)
\end{gather}
with $\nabla^{L^k} : \Ci (M, L^k) \rightarrow \Om^1 (M, L^k)$ the covariant derivative induced by $\nabla$, and the Riemannian metric $g(X,Y) = \om (jX, Y)$ independent of $k$.


\subsubsection*{Earlier results}

It is known that the spectrum $\si (\Delta_k)$ of $\Delta_k$ is partitioned into clusters around each point of $k (\frac{n}{2} + \N)$ in the large $k$ limit. More precisely, for any $m \in \N$, define the interval $I_m$
\[
I_0 = \left[0, \tfrac{n}{2} + \tfrac{1}{2}\right], \qquad I_m = \left(\tfrac{n}{2} +m\right) + \left[-\tfrac{1}{2}, \tfrac{1}{2} \right[\text{ if
} m \geqslant 1,
\]
so that we have a partition $[0,\infty [= \bigcup _{m\,\in\,\N} I_m$. Then we set
\begin{gather}\label{eq:def_Landau}
\Si_{m,k} := \bigl(k^{-1} \si (\Delta_k) \bigr) \cap I_m, \qquad \Hilb_{m,k} := \bigoplus_{\la\,\in\,\Si_{m,k}} \ker \left(k^{-1}\Delta_k - \la\right).
\end{gather}
It was proved by Faure--Tsuji~\cite{FaTs15} that
\begin{gather}\label{eq:Si_m:Faure_Tsujii}
\Si_{m,k} \subset
\Bigl(\tfrac{n}{2}+ m + C_m k^{-\frac{1}{4}} [-1,1] \Bigr)
\end{gather}
and by Demailly~\cite{Dem} that
\begin{gather}\label{eq:asymptot_Demailly}
\op{dim} \Hilb_{m,k} = \Bigl(\frac{k}{2 \pi} \Bigr)^n \binom{ {m+n-1}}{{n-1} } \op{vol} (M) + \op{o} (k^n)
\end{gather}
For a surface $(n=1)$ with a constant Gauss curvature $S$, more precise results have been obtained by Iengo--Li~\cite{IeLi94}: if $k + m S >0$, then
\begin{gather}\label{eq:surface_eignvalue_dim}
\begin{split}
\Si_{m,k} = \left\{ \tfrac{1}{2} + m + k^{-1} S \tfrac{m(m+1)}{2} \right\}, \\
\op{dim} \Hilb_{m,k} = \tfrac{k}{2 \pi} \op{vol} (M) + \left(\tfrac{1}{2} + m\right)
\chi (M).
\end{split}
\end{gather}
So in this case, when $k$ is sufficiently large, the $m^{\rm th}$ eigenvalue is degenerate with multiplicity equal to $\op{dim} \Hilb_{m,k}$.

The first cluster has been further studied. In the K\"ahler case, that is when the complex structure $j$ is integrable, $L$ has itself a natural holomorphic structure such that $\con{\partial}_L = \nabla^{0,1}$ and by Kodaira identities, we have when $k$ is sufficiently large that
\begin{gather}
\Si_{0,k} = \left\{ \tfrac{n}{2} \right\}, \qquad \Hilb_{0,k} = H^0 \left(M, L^k\right).
\end{gather}
and the dimension of $\Hilb_{0,k}$ is given by the Riemann--Roch--Hirzebruch theorem
\begin{gather}\label{eq:RRH}
\dim \Hilb_{0,k} = \int_M \exp \left(\frac{k\om}{2\pi} \right) \; \op{Todd} M
\end{gather}
Here $\op{Todd} M$ is the Todd class of $(M,j)$.

More generally, when $j$ is not necessarily integrable, it was proved by Guillemin--Uribe~\cite{GuUr88} that
\begin{gather}\label{eq:si_0-k_Guillemin_Uribe}
\Si_{0,k} \subset \left(\tfrac{n}{2} + C_0 k^{-1} [-1,1] \right)
\end{gather}
and by Borthwick--Uribe~\cite{BoUr96} that the dimension of $\Hilb_{0,k}$ is given by~\eqref{eq:RRH} when $k$ is sufficiently large.

\subsection{Main results}

In the sequel $m \in \N$ is a fixed non negative integer and all the results hold in the large $k$ limit, with estimates, bounds depending on $m$.

\subsubsection{Dimension}
Our first result is the computation of the dimension of $\Hilb_{m,k}$, as the Riemann--Roch number of $L^k \otimes \D_m (TM)$ where $\D_m (TM)$ is the $m^{\rm th}$ symmetric power of $(T^{0,1} M)^*$. Here, $T^{0,1}M = \ker (j + i)$ and $j$ is the almost complex structure introduced previously. The reason why we prefer to work with $(T^{0,1}M)^*$ instead of the isomorphic bundle $T^{1,0}M$ should be clear later.

\begin{theo}\label{theo:dim_landau}
If $k$ is sufficiently large, then
\[
\op{dim} \Hilb_{m,k} = \int_M \exp \left(\frac{k\om}{2\pi} \right) \; \op{ch} (\D_m (TM))) \op{Todd} M
\]
with $\op{ch}$ the Chern character and $\op{Todd} M$ the Todd class of $(M,j)$.
\end{theo}

As far as we know, Theorem~\ref{theo:dim_landau} is a new result, except in the cases already mentioned ($n=1$ with constant curvature or $m=0$).

\begin{rema}\leavevmode
\begin{itemize}
\item[-] When $n=1$, $\D_m (TM)$ is isomorphic with $K^{-m}$, $K$ being the canonical bundle, and it is easy to see that we recover the second equation~\eqref{eq:surface_eignvalue_dim}. However, even for $n=1$, Theorem~\ref{theo:dim_landau} goes further since we don't assume that the Gauss curvature is constant. In this generality, it is not likely that $\Si_{m,k}$ consists of a single degenerate eigenvalue, but the dimension of the $m^{\rm th}$ cluster is given by the same formula.
\item[-] For a general dimension $n$, $\D_m (TM)$ has rank $\binom{ {m+n-1}}{{n-1} }$ and we recover the asymptotic~\eqref{eq:asymptot_Demailly}.
\end{itemize}
\end{rema}



\subsubsection{Symbol spaces}
In the sequel, we will use $\D_m (TM)$ as a bosonic space, with associated creation an annihilation operators defined as follows. For any $x \in M$, let us view $ \D_{m} (T_xM)$ as the space of homogeneous polynomials maps $T_x^{0,1} M \rightarrow \C$ with degree $m$. Set
\begin{gather}\label{eq:dtm_intro}
\D (T_xM) := \bigoplus_{m\,\in\,\N} \D_m (T_xM)
\end{gather}
and let $\pi_m(x)$ be the corresponding projector of $\D (T_x M)$ onto $\D_m (T_x M)$. For any $ Y
\in T_xM \otimes \C$, let $\rho (Y)$ be the endomorphism of $\D (T_xM)$ defined as follows. Write $Y = U + \con V$ with $U, V \in T_x ^{1,0} M$. Then
\begin{gather}\label{eq:rho}
\rho (Y)=
\rho (U) + \rho (\con V) \; \text{ with } \;
\begin{cases} \rho (U) = \text{
multiplication by } i
\om (U, \cdot) \\ \rho (\con {V}) = \text{ derivation with respect to } \con V
\end{cases}
\end{gather}
More concretely, let $(U_i)$ be a basis of $T^{1,0}_xM$ such that $\frac{1}{i} \om (U_i,
\con {U}_j) = \delta_{ij}$. Let $(z_i)$ be the basis of $(T^{1,0}_xM)^*$ dual to $(U_i)$, so $z_i = \frac{1}{i} \om (\cdot, \con{U}_i)$ and $ \D (T_x M) = \C [\con{z}_1, \,\ldots,\,\con{z}_n]$. Then for any polynomial $P$ of the variable $\con{z}_1$, $\ldots$, $\con{z}_n$
\begin{gather}
\rho (U_i) P = - \con{z}_i P, \qquad \rho (\con{U}_i) P = \frac{\partial P}{\partial
\con{z}_i}.
\end{gather}
So $-\rho (U_i)$ and $\rho (\con{U}_i)$ are respectively the creation and annihilation operators.

\subsubsection{Berezin--Toeplitz operators}

Our second result is about Berezin--Toeplitz operators. By~\cite{BoUr96}, the spaces $\Hilb_{0,k}$ can be considered as quantizations of $M$, replacing the standard K\"ahler quantization $H^0 (M,L^k)$, for symplectic manifold not necessarily having an integrable complex structure. An important feature is that there is a natural way to pass from classical to quantum Hamiltonians, provided by the Berezin--Toeplitz quantization. In the semi-classical limit, defined here as the large $k$ limit, the product and commutator of quantum observables correspond to the product and Poisson bracket of classical observables, up to some error terms. More precisely, let $\Pi_{m,k}$ be the orthogonal projector of $\Ci (M, L^k)$ onto $\Hilb_{m,k}$ and for any $f \in \Ci (M)$, let $T_{m,k} (f)$ be the endomorphism of $\Hilb_{m,k}$ defined by
\begin{gather}\label{eq:toeplitz}
T_{m,k} (f) \psi = \Pi_{m,k} (f \psi)
\qquad \forall \; \psi \in \Hilb_{m,k}.
\end{gather}
For the first Landau level, it is known~\cite{BoMeSc,oim_op,MaMa} that for any $N$
\begin{gather}\label{eq:t_0_k_produit_tout_ordre}
T_{0,k} (f) T_{0,k} (g) = \sum_{\ell = 0 } ^N k^{-\ell} T_{0,k} (B_\ell (f,g)) +
\bigo \left(k^{-(N+1)}\right)
\end{gather}
for some bidifferential operators $B_{\ell} : \Ci (M) \times \Ci (M)
\rightarrow \Ci (M)$, where
\begin{gather}\label{eq:b_0_b_1}
B_0 (f,g) = fg, \qquad B_1(f,g) = - \tfrac{1}{2} g(X,Y) + \tfrac{1}{2i} \om (X, Y),
\end{gather}
$X$ and $Y$ being the Hamiltonian vector fields of $f$ and $g$ respectively.


For the generalisation to higher Landau levels, we will use in addition to the $T_{m,k} (f)$'s the following operators: let $p \in \N$ and $X_1,\, \ldots,\, X_{2p}$ be vector fields of $M$. Define $ T_{m,k} (X_1,\, \ldots,\, X_{2p}):
\Hilb_{m,k} \rightarrow \Hilb_{m,k}$ by
\begin{gather}\label{eq:toep_der}
T_{m,k} (X_1, \,\ldots,\, X_{2p}) (\Psi) = k^{ -p} \Pi_{m,k} \left(\nabla^{L^k}_{X_1} \ldots
\nabla_{X_{2p}}^{L^k} \Psi\right)
\end{gather}
Let $\clT_m$ be the vector space of families $ (P_k : \Hilb_{m,k} \rightarrow
\Hilb_{m,k}, \; k \in \N)$ spanned by the $(T_{m,k} (f), k \in \N)$'s and $(T_{m,k} (X_1, \,\ldots,\, X_{2p}), \; k \in \N) 's$. Here the functions $f$ or vector fields $X_1,\,
\ldots X_{2p}$ do not depend on $k$.

 Define the semiclassical completion $\clT_m ^{\op{sc}}$ as the vector space of families $(P_k : \Hilb_{m,k} \rightarrow \Hilb_{m,k}, \; k \in
\N)$ such that for any $N$,
\[
P_k = \sum_{\ell= 0 } ^N k^{-\ell} P_{\ell,k} + \bigo \left(k ^{-N-1}\right)
\]
where the coefficients $(P_{\ell, k })_k$, $\ell \in \N$ all belong to $\clT_m$, and the $\bigo $ is for the operator norm. This expansion is meaningful because as we will see later, for any $\ell$, the operator norm $\| P_{\ell, k}\|$ is bounded independently of $k$.



\begin{theo}\label{theo:Toeplitz_Landau}
For any $m \in \N$, we have:
\begin{enumerate}
\item\label{theo1.3.1} $ \clT_m^{\op{sc}}$ is closed under product.
\item\label{theo1.3.2}  There exists a unique linear map $ \mathcal{\tau} : \clT_m^{\op{sc}} \rightarrow \Ci (M,
\op{End} (\D_m (TM))) $ given on the generators ~\eqref{eq:toeplitz} and~\eqref{eq:toep_der} by
\begin{equation}\label{eq:symbolmap}
\begin{split}
\tau (T_{m,k} (f))_x & = f(x) \op{id}_{\D_m (T_x M)}, \\ \tau \left(T_{m,k} \left(X_1,\, \ldots,\, X_{2p}\right)\right)_x & = \pi_m(x) \rho(X_1(x))
\ldots \rho (X_{2p}(x))
\end{split}
\end{equation}
$\tau$ is onto, its kernel consists of $k^{-1} \clT_{m}^{\op{sc}}$.
\item\label{theo1.3.3} For any $P,Q \in \clT_{m}^{\op{sc}}$
\begin{align}
 &\tau (PQ) = \tau (P) \tau (Q),\label{eq:prod} \\
 &\| P_k \| = \sup \big\{ \| \tau (P)_x \|, \; x \in M \big\} + \op{o} (1),  \label{eq:norm}\\ 
 &P_k(x,x) = \Bigl(\frac{k}{2\pi} \Bigr)^n \left(\op{tr} (\tau (P)_x) + \bigo \left(k^{-1}\right)\right).\label{eq:trace}
\end{align}
\item\label{theo1.3.4} For any $f$, $g$ in $\Ci (M)$, we have
\begin{gather}\label{eq:produit}
T_{m,k} (f) T_{m,k} (g) = T_{m,k} (fg) + k^{-1} T_{m,k} (X,Y) + \bigo \left(k^{-2}\right)
\end{gather}
where $X$, $Y$ are the Hamiltonian vector fields of $f$ and $g$. In particular,
\begin{gather}\label{eq:com_poisson}
ik \left[T_{m,k} (f), T_{m,k} (g)\right] = T_{m,k} (\{ f,g \}) + \bigo \left(k^{-1}\right)
\end{gather}
with $\{ \cdot, \cdot \}$ the Poisson Bracket of $(M, \om)$.
\end{enumerate}
\end{theo}

We call $\tau$ the symbol map. The symbol of the generators~\eqref{eq:symbolmap} is defined in terms of the endomorphisms~\eqref{eq:rho}. The product of symbols in the right-hand side of~\eqref{eq:prod} is the pointwise composition. In the norm estimate~\eqref{eq:norm}, the norm of $\tau (P)_x$ is defined in terms of the hermitian structure of $\D_m (T_xM)$. In~\eqref{eq:trace}, $P_k (x,x)$ is the value of the Schwartz kernel of $P_k$ at $(x,x)$. Integrating~\eqref{eq:trace}, we obtain the following estimate of the trace of $P_k$,
\begin{gather}
\op{tr} P_k = \Bigl(\frac{k}{2\pi} \Bigr)^n \int_M \op{tr} (\tau (P)_x) \mu_M(x) + \bigo \left(k^{n-1}\right)
\end{gather}
where $\mu_M = \om^n /n!$. Since $P_k =T_{m,k} (1)$ is the identity of $\Hilb_{m,k}$, we recover the estimate~\eqref{eq:asymptot_Demailly}.

\break
\begin{rema}\leavevmode \label{rem:Toeplitz}
\begin{enumerate}
\item\label{rema1.4.1} In the surface case, $n=1$, $\D_m(TM)$ is a line bundle, so any endomorphism of $\D_m(T_xM)$ is scalar and by Assertion 2 of Theorem~\ref{theo:Toeplitz_Landau}, $\clT_{m}^{\op{sc}}$ consists of the families
\begin{gather}\label{eq:-p_k-=}
(P_k = T_{m,k} \big(f(\cdot, k)) + \bigo (k^{-\infty}), \; k \in \N\big)
\end{gather}
where the multiplicator $f(\cdot, k)$ depends on $k$ in such a way that it admits an expansion $f(\cdot, k) = f_0 + k^{-1} f_1 + \ldots\ $. It holds as well that the $T_{m,k} (f)$ satisfy~\eqref{eq:t_0_k_produit_tout_ordre} for some bidifferential operators $B_\ell^m$ depending on $m$. Indeed, by~\cite[Section~5.4]{oim}, this property is equivalent to the locality of the product: for any two functions $f$, $g$ with disjoint supports, $T_{m,k} (f) T_{m,k} (g) = \bigo (k^{-\infty})$. This latter property follows from the fact that the Schwartz kernel of $\Pi_{m,k}$ is in $\bigo (k^{-\infty})$ outside the diagonal. 

\eqref{eq:produit} writes in this case
\begin{gather}\label{eq:bm_1-f-g}
B^m_1 (f,g) = - \left(\tfrac{1}{2} +m\right) g(X,Y) + \tfrac{1}{2i} \om (X,Y)
\end{gather}
where $X,Y$ are the Hamiltonian vector fields of $f$ and $g$.

\item\label{rema1.4.2} If $n \geqslant 2$ and $m\geqslant 1$, the Toeplitz algebra $\clT_m^{\op{sc}}$ is strictly larger than the subspace $
\clT_m^{\op{f}}$ consisting of the Toeplitz operators~\eqref{eq:-p_k-=} with scalar multiplicators, because the symbol map $\tau $ is onto. Nevertheless, we may ask if $ \clT_m^{\op{f}}$ is closed under product. This is not the case. Indeed, if follows from~\eqref{eq:produit}, that if $f \in \Ci (M)$ is not locally constant, then there exists no function $h$ such that $T_{m,k} (f)^2 - T_{m,k} (f^2) = k^{-1} T_{m,k} (h) + \bigo (k^{-2})$. So interestingly, by developing the theory of Berezin--Toeplitz operator for higher Landau level, we are naturally led to use matrix valued symbols.
\item\label{rema1.4.3}  The Toeplitz operators defined as in~\eqref{eq:toep_der}, but with an odd number of vector fields, can be incorporated in the theory. They are odd Toeplitz operators, cf. Remark~\ref{rem:odd_toeplitz}, and they belong to $k^{-\frac{1}{2} } \clT_m^{\op{sc}}$. We purposely have avoided any square root of $k$ in Theorem~\ref{theo:Toeplitz_Landau}.
\end{enumerate}
\end{rema}

\begin{rema}\leavevmode \label{sec:comp_Yuri}
\begin{enumerate}
\item\label{rema1.5.1}  The main estimates for the Toeplitz operators associated to functions, that is $T_{m,k} (f) T_{m,k} (g) = T_{m,k} (fg) + \bigo (k^{-1})$ and $ik [T_{m,k} (f), T_{m,k} (g)]= T_{m,k} (\{ f, g \}) + \bigo (k^{-1})$ have been proved independently by Kordyukov~\cite{Yuri_2} by using the techniques of~\cite{MaMa}.
\item \label{rema1.5.2} In~\cite{Yuri_2}, the non degenerate magnetic fields $j_B$ with constant eigenvalues $B_1,\, \ldots,\, B_n$ are considered as well. The corresponding magnetic Laplacian has clusters centered at the points of $k\Si$ where $\Si = \{ \sum B_i (\frac{1}{2} + \al (i)) / \al \in \N^n \}$. In the companion paper~\cite{oim_copain}, we prove that the number of eigenvalues in the cluster at $\La \in \Si$ is given by a Riemann--Roch number as in Theorem~\ref{theo:dim_landau} where $\D_m (TM)$ is replaced by a bundle $F_{\La}$ \footnote{$F_{\Lambda}$ is defined as follows: for any $p \in M$, introduce a basis $(\partial_j, \con{\partial}_j)$ of $T_pM \otimes \C$ such that $j_B
\partial_j = i B_j \partial_j$ and denote by $(z_j, \con{z}_j)$ the dual basis of $T^*_pM \otimes \C$. Then the fiber of $F_{\Lambda}$ at $p$ is spanned by the $\con{z}_1^{\al(1)}
\ldots\, \con{z}_n^{\al (n)}$ where $\al \in \clK_{\Lambda}$. } with rank the cardinal of
\[
\clK_{\La} = \left\{ \al \in \N^n, \sum B_i
\left(\tfrac{1}{2} + \al (i)\right) = \Lambda \right\}.
\]
When $|\clK_{\Lambda} | =1$, it is proved in~\cite{Yuri_2} that the space of Toeplitz operators in the $\Lambda$-cluster associated to functions of $M$ is an algebra. This particular case is very similar to the first Landau level~($m=0$) because the endomorphisms of $K_{\Lambda}$ are scalar, the symbol composition law is commutative and it is not necessary to introduce the operators~\eqref{eq:toep_der}. Observe as well that in dimension $n \geqslant 2$, the condition that $|\clK_{\Lambda}| =1$ for $\Lambda \neq \frac{1}{2} \sum B_i$ implies that the $B_i$ are not all equal, so that the tangent bundle of $M$, endowed with the unique up to isotopy complex structure compatible with the symplectic structure, splits into a non trivial sum of complex subbundles. This seems rather restrictive and is not satisfied for instance by the complex projective spaces $\bbP^n$ when $n$ is even, cf.~\cite{split}.

\item\label{rema1.5.3} We can study as well with our techniques the Toeplitz operators associated to the $\Lambda$-cluster and this will be partly done in Section~\ref{sec:projectors-Toeplitz} where we consider a vector bundle $F$ generalising $\D_m (TM)$ or $F_{\Lambda}$. Assuming that the numbers $|\al| = \al (1) + \ldots + \al (n)$ have all the same parity when $\al$ runs over $\clK_{\La}$, we have the same result as Theorem~\ref{theo:Toeplitz_Landau}. When this parity condition is not satisfied, the square roots of $k$ seem to be unavoidable.

We have chosen to work out the case where all the $B_i$ are equal because first any symplectic manifold can be endowed with such a magnetic field, which is merely a compatible almost complex structure, and second the higher rank vector bundle $\D_m (TM)$ makes it rather different from the already much studied case of first Landau levels.
\end{enumerate}

\end{rema}

\subsubsection{Ladder operators}

The last result we would like to emphasize in this introduction is the construction of some ladder operators for the spaces $\Hilb_{m,k}$. In the surface case with constant Gauss curvature, $\Hilb_{m,k}$ is naturally isomorphic with the space of holomorphic sections of $L^k \otimes K^{-m}$ where $K$ is the canonical bundle~\cite{Te06}, the isomorphism being the \emph{ladder} operator $\con{\partial}_{L^k \otimes K^{-m+1}} \circ \ldots
\circ \con{\partial}_{L^k}$, cf. the Appendix~\ref{sec:appendix}. Here we will show that the family $(\Hilb_{m,k}, \; k \in \N)$ is isomorphic to a quantization of $M$ twisted by the vector bundle $\D_m (TM) $.

Recall that for any Hermitian vector bundle $F \rightarrow M $, we can define a family of finite dimensional subspaces $\Hilb_{F,k} \subset \Ci (M, L^k \otimes F)$, $k
\in \N$, having the following properties:
\begin{enumerate}
\item $\op{dim} \Hilb_{F,k} = \int_M \op{ch} (L^k \otimes F) \; \op{Todd} M$, when $k$ is sufficiently large.
\item the space $\clT_F^{\op{sc}}$, consisting of families of $(T_k \in \op{End} (\Hilb_{F,k}), \; k \in \N)$ having an expansion of the form
\[
T_k = \sum_{\ell = 0 }^N k^{-\ell} T_{F,k} (f_{\ell}) + \bigo \left(k^{-N+1}\right), \qquad
\forall N \in \N
\]
for a sequence $(f_\ell)$ of $\Ci (M, \op{End} F)$, is closed under product. Here, $T_{F,k} (f_{\ell}) (\psi)\linebreak = \Pi_{F,k} (f_{\ell} \psi)$ for any $\psi
\in \Hilb_{F,k} $ where $\Pi_{F,k}$ is the orthogonal projector of $\Ci (M, L^k
\otimes F)$ onto $\Hilb_{F,k}$.
\item At first order, the product is given by the pointwise product, that is 
\[
T_{F,k} (f) T_{F,k} (g) = T_{F,k} (fg) + \bigo \left(k^{-1}\right)
\]
for any $f,g \in \Ci (M,\op{End} F)$.
\end{enumerate}
The algebra $\clT_F^{\op{sc}}$ is called the Toeplitz algebra. In the K\"ahler case, that is when $(M, \om)$ is K\"ahler, $L$ holomorphic with $\nabla$ the Chern connection, and $F$ holomorphic as well, the space $\Hilb_{F,k}$ can be defined as the space $ H^0 (L^k \otimes F)$ of holomorphic sections. In the non-K\"ahler case, various constructions have been developed~\cite{BoUr96, MaMa}: Spin-c quantization, first Landau level of a Laplacian acting on $\Ci (M, L^k
\otimes F)$, or more generally any image of a projector of $\Ci (M, L^k
\otimes F)$ having a specific Schwartz kernel~\cite{oim}. Let us call such a family $(\Hilb_{F,k}, \; k \in \N)$ a {\em quantization of $(M,L)$ twisted by $F$}. These twisted quantizations have sometimes better properties than the non twisted one (corresponding to $F=\C$), typically when $F$ is a half-form bundle~\cite{oim_eq}. The general case where the rank of $F$ is $\geqslant 2$ may be viewed as a free generalization without any application, but interestingly, this is exactly what we need.

Assume $F$ is equipped with a connection $\nabla^F : \Ci (M, F) \rightarrow
\Om^1 (M, F)$. Let $G = (T^{0,1}M)^*$ and $D_{F,k} : \Ci (M, L^k \otimes F) \rightarrow \Ci (M, L^k
\otimes F \otimes G) $ be the $(0,1)$-part of the connection $\nabla^{F \otimes L^k}$ induced by $\nabla^F$ and $\nabla^{L^k}$. Endow $G$ with a connection and define the differential operators
\begin{gather}\label{eq:def_w_k}
\begin{split}
W_k : \Ci \left(M, L^k\right) \rightarrow \Ci \left(M, L^k \otimes \D_m (TM)\right)\\
W_k = R_m D_{G^{\otimes (m-1)},k} \circ D_{G^{\otimes (m-2)},k} \circ \ldots \circ D_{G, k} \circ D_{\C,k}
\end{split}
\end{gather}
where $R_m$ is the projection from $G^{\otimes m }$ onto $\D_m (TM) =
\op{Sym}^m G$.

\begin{theo}\label{theo:ladder}
For any quantization $(\Hilb_{F,k}, \; k \in \N)$ of $(M,L)$
twisted by $F= \D_m (TM)$, the linear maps
\[
V_k = \tfrac{1}{m!} k^{-\frac{m}{2}} \Pi_{F,k} W_k
: \Hilb_{m,k} \rightarrow \Hilb_{F,k}, \qquad k \in \N
\]
satisfy:
\begin{enumerate}
\item $V_kV_k^* = \op{id}_{\Hilb_{F,k}} + \bigo (k^{-1})$ and $V_k^* V_k =\op{id}_{\Hilb_{m,k}} + \bigo (k^{-1})$. In particular, $V_k$ is an isomorphism when $k$ is sufficiently large.
\item the conjugation by $V=(V_k)$ is an isomorphism between the Toeplitz algebra $ \clT^{\op{sc}}_m$ and $\clT_F^{\op{sc}}$ modulo $\bigo (k^{-\infty})$. In particular, for any $(P_k) \in
\clT^{\op{sc}}_m$, $(V_k P_k V_k^*)_k$ belongs to $\clT_F^{\op{sc}}$ and if $f \in \Ci (M, \op{End} F)$ is the symbol $ \tau(P_k)$, then $V_k P_k V_k^* = T_{F,k} (f) + \bigo (k^{-1})$.
\end{enumerate}
\end{theo}

The first assertion of Theorem~\ref{theo:ladder} tells us that $V_k$ is almost unitary. This can be improved by setting $U_k := A_k V_k$ with $A_k$ the endomorphism of $\Hilb_{F,k}$ equal to $(V_kV_k^*)^{-1/2} |_{\Hilb_{m,k}}$ when $k$ is sufficiently large and to $0$ for the first values of $k$. Then $U_k U_k^* = \op{id}_{\Hilb_{F,k}}$ and $U_k^* U_k =
\op{id}_{\Hilb_{m,k}}$ when $k$ is sufficiently large. Furthermore the second assertion of~\ref{theo:ladder} holds with $(U_k)$ instead of $(V_k)$.

 The inspiration for~\eqref{eq:def_w_k} comes from the case of surface $(n =1)$ with constant Gauss curvature. In this case, choosing for $D_{F,k}$ the $
\con\partial$ operator, it holds that $W_k (
\Hilb_{m,k}) \subset \Hilb_{F,k}$, cf. the discussion after Theorem~\ref{theo:appendix}, so the projector $\Pi_{F,k}$ in the definition of $V_k$ is not necessary. We expect that something similar happens in higher dimension under the convenient assumptions. 


\subsection{Generalised Landau level and Schwartz kernel expansion}

\subsubsection{Generalised Landau level}

In the previous results, the $m^{\rm th}$ Landau level $\Hilb_{m,k}$, $k\in \N$ can be replaced by any family $(\Hilb_{m,k} \subset \Ci (M, L^k), \; k \in \N) $ of finite dimensional subspaces, such that the Schwartz kernel of the orthogonal projector $\Pi_{m,k}$ of $\Ci (M, L^k)$ onto $\Hilb_{m,k}$ has a specific behavior in the large $k$ limit. We require first that
\begin{gather}\label{eq:first_order}
\Pi_{m,k} (x,y) = \Bigl(\frac{k}{ 2\pi} \Bigr)^n E^k (x,y) Q^{(n-1)}_m (k
\delta (x,y)) + \bigo \left(k^{n-1}\right)
\end{gather}
where
\begin{itemize}
\item[-] $E$ is a section of $L \boxtimes \con{L}$ such that $|E(x,y)| < 1$ when $x \neq y$, and its second order Taylor expansion along the diagonal has a specific form, cf. Equation~\eqref{eq:hypotheseE}.  In particular, in a coordinate chart at $x$, $\ln |E(x+
\xi, x)|= - \frac{1}{4} |\xi|_x^2 + \bigo (|\xi|^3)$ where $|\cdot |_x$ is the norm of $T_xM$.  \footnote{In the whole paper, when working in a coordinate chart $(U, \chi)$ of $M$, we write $x+ \xi$ for $ \chi^{-1} (
\chi (x) + T_x \chi (\xi))$ where $x \in U$ and $\xi \in T_xM$ sufficiently close to the origin.}
\item[-] $\delta \in \Ci (M^2)$ is any function vanishing to second order along the diagonal and satisfying $\delta (x+ \xi, x) = |\xi|_x^2 + \bigo (|\xi|^3)$, $\xi \in T_xM$. $Q_m^{(p)} $ is the generalised Laguerre polynomial $Q^{(p)}_m (x) = \frac{x^{-p}}{m!} (\frac{d}{dx} -1 ) ^m x ^{m+p}$.
\end{itemize}
The Schwartz kernel of the projector onto the $m^{\rm th}$ Landau level of the magnetic Laplacian of $\C^n$ is given by an expression similar to~\eqref{eq:first_order}, with the convenient section~$E$ and the same Laguerre polynomials, cf.~\eqref{eq:noyau_niveau_m} and~\eqref{eq:2pin-e_cn-u}. 

In addition to~\eqref{eq:first_order}, we require a full expansion of the form
\begin{gather}\label{eq:full}
\Pi_{m,k} (x,y) = \left(\frac{k}{ 2\pi} \right)^n E^k (x,y) \sum _{\ell\,\in\,\Z} k^{-\ell} a_{\ell} (x,y) + \bigo (k^{-\infty})
\end{gather}
with coefficients $a_{\ell} \in \Ci (M^2)$ such that for $\ell<0$, $a_{\ell}$ vanishes to order $m(\ell) \geqslant -2 \ell$ along the diagonal and $m(\ell) + 2\ell \rightarrow
\infty $ as $\ell \rightarrow - \infty$. The meaning of this expansion is not obvious because the negative $\ell$'s give positive powers of $k$. Actually, the condition satisfied by $|E|$ implies that $|E^k (x,y) b(x,y) | = \bigo (k^{-m/2})$ when $b$ vanishes to order $m$ along the diagonal. So the $\ell^{\rm th}$ summand in~\eqref{eq:full} is in $\bigo (k^{n - \frac{1}{2} (m (\ell) + 2\ell)})$, and the expansion is meaningful because of the conditions satisfied by $m (\ell)$.

We will prove that for any family $(\Hilb_{m,k})$ whose associated projector $\Pi_{m,k}$ satisfies~\eqref{eq:first_order} and~\eqref{eq:full}, Theorems~\ref{theo:dim_landau}, \ref{theo:Toeplitz_Landau} and~\ref{theo:ladder} hold. On the other hand, in the second part of this work~\cite{oim_copain}, cf. also~\cite{Yuri_1}, it is proved that the Schwartz kernel of the orthogonal projector $\Pi_{m,k}$ onto the Landau levels $\Hilb_{m,k}$ defined in~\eqref{eq:def_Landau} from the Laplacian $\Delta_k$, satisfies~\eqref{eq:first_order} and~\eqref{eq:full}. The assumption that the magnetic field is constant with respect to the metric can be relaxed. It is actually possible to define some Landau levels and describe the asymptotic expansion of the associated projector as soon as a particular gap condition is satisfied, cf~\cite{oim_copain}.

\break
\subsubsection{The class $\lag (A,B)$}
To establish our results, we will introduce a specific class of operators, containing the projector $\Pi_{m,k}$, the Berezin Toeplitz operators $T_{m,k} (f)$ and $T_{m,k} (X_1,\,\ldots,\,X_{2p})$ and also the projector $\Pi_{F,k}$ of any twisted quantization, the corresponding Toeplitz operators $T_{F,k} (g)$, the isomorphisms $V_k$ of Theorem~\ref{theo:ladder} and their unitarizations $(U_k)$. This operator class has a natural filtration, with associated symbol spaces, which allows to prove most of the results by successive approximations as often in microlocal analysis. Interestingly, in the symbolic calculus appear the eigenprojectors of  the Landau Laplacian of $\C^n$, providing another link between the usual Landau levels and our geometric Landau levels. 

Introduce two auxiliary Hermitian vector bundles $A$, $B$ over $ M$. Then $\lag (A,B)$ consists of families $(P_k : \Ci(M, L^k \otimes A) \rightarrow \Ci(M, L^k \otimes B), \; k \in \N)$ of operators having a smooth Schwartz kernel satisfying
\begin{gather}\label{eq:exp_lag}
P_k (x,y) = \Bigl(\frac{k}{2\pi} \Bigr)^n E^{k}(x,y) \sum_{\ell\, \in\,\Z } k^{-\frac{\ell}{2}} b_{\ell}(x,y) + \bigo \left(k^{- \infty}\right)
\end{gather}
where $E$ is defined as in~\eqref{eq:first_order}; the coefficients $b_{\ell} $ are in $\Ci (M^2, B \boxtimes \con{A})$; for $\ell <0$, $b_{\ell}$ vanishes to order $m (\ell)
\geqslant -\ell $ along the diagonal; $m(\ell) + \ell \rightarrow \infty $ as $\ell
\rightarrow - \infty$, and the meaning of this expansion is the same as in~\eqref{eq:full}. We have a decomposition into even/odd elements: $(P_k) \in \lag^{+} (A,B)$ (resp. $\lag^{-} (A,B)$) if the expansion~\eqref{eq:exp_lag} holds with a sum over the $\ell$'s even (resp. odd).

The main property is that this class of operators is closed under composition:
\begin{gather}\label{eq:comp}
\lag^{\ep'} (B, C) \cdot \lag^{\ep} (A,B) \subset \lag^{\ep \ep'} (A,C),
\qquad \ep, \ep' \in \{ \pm 1\}.
\end{gather}
In particular, $\lag^+ (A) := \lag^{+} (A,A)$ is an algebra. We also have a filtration $ \lag_q(A,B) := \lag (A,B) \cap
\bigo (k^{-q/2})$, $q \in \N$ and the corresponding graduation is described by symbol maps $\si_q$
\[
0 \rightarrow \lag_{q+1} (A,B) \rightarrow \lag_q (A,B) \xrightarrow{\si_q} \Ci \big(M
, \symb (M) \otimes \op{Hom} (A,B)\big) \rightarrow 0
\]
Here, $\symb (M)$ is an infinite rank vector bundle over $M$, each fiber $\symb_x(M)$ is a subalgebra of the algebra of endomorphisms of the space $\D (T_xM)$ defined in~\eqref{eq:dtm_intro}. This is compatible with the composition~\eqref{eq:comp} in the sense that $\lag_p (B,C) \cdot \lag_q (A,B) \subset \lag_{p+q} (A,C)$ and the corresponding product of symbols is the pointwise product of $\symb(M)$ tensored by $\op{Hom} (B,C) \otimes \op{Hom} (A,B) \rightarrow \op{Hom} (A,C)$.


The projector $(\Pi_{m,k})_k$ is an idempotent of $\lag^+ (\C)$ with symbol $\si_0 (\Pi_{m})$ equal at $x$ to the projector $\pi_m (x)$ onto the $m^{\rm th}$ summand in~\eqref{eq:dtm_intro}. The Toeplitz algebra introduced previously is
\begin{gather}\label{eq:toep_land}
\clT_m^{\op{sc}} = \big\{ P
\in \lag^{+} (\C) / \; \Pi_m P \Pi_m = P \big\}.
\end{gather}
The isomorphism $V$ of Theorem~\ref{theo:ladder} belongs to $\lag (\C, F)$ and has the same parity as~$m$.

Interestingly, $\symb_x(M)$ has a representation as operators of $L^2 (\C^n)$, and in this representation, $ \pi_m(x) = \si_0 (\Pi_m)(x)$ is the projector onto the $m^{\rm th}$ Landau level of a magnetic Laplacian of $\C^n$.


\break
\subsubsection{Outline of the paper}
In Section~\ref{sec:class-lag-pres}, we introduce the class $\lag (A,B)$, and state its main properties. In Section~\ref{sec:projectors-Toeplitz}, we prove variations of the theorems stated before, where the $\Hilb_{m,k}$ are subspaces of $\Ci (M, L^k \otimes A)$ such that the corresponding family $(\Pi_{m,k})$ of orthogonal projectors belongs to $\lag^+ (A)$ with a convenient symbol. Sections~\ref{sec:landau-levels-cn}, \ref{sec:schw-kern-oper} and~\ref{sec:derivatives} are devoted to the proof of the properties of $\lag (A,B)$. The proofs of Theorems~\ref{theo:dim_landau}, \ref{theo:Toeplitz_Landau} and~\ref{theo:ladder} is given in the last Subsection~\ref{sec:ledernier}. In the Appendix~\ref{sec:appendix}, we prove formulas~\eqref{eq:surface_eignvalue_dim} on constant curvature surface.



\subsection*{Acknowledgments}
I would like to thank Yuri Kordyukov for his collaboration at an early stage of this work. I thank as well the anonymous referees for their careful reading, their comments and their suggestions.




% Official author native source lines 529--1102
\section{The class \texorpdfstring{$\lag (A,B)$}{L(A,B)}} \label{sec:class-lag-pres}

We start the discussion with the algebra in which the symbol of operators of $\lag (A, B)$ takes their values. The class $\lag (A,B)$ is defined in Subsection~\ref{sec:operators}  and its main properties are stated, the proof are postponed to Section~\ref{sec:schw-kern-oper}. 

\subsection{Symbol spaces} \label{sec:symbol-spaces}

\subsubsection{The algebra $\symb (\C^n)$}
Let $n$ be a positive integer and denote by $z_1,\, \ldots,\, z_n$ the linear coordinates of $\C^n$. Let $\D (\C^n)= \C [\con{z}_1, \,\ldots,\, \con{z}_n]$ be the space of antiholomorphic polynomial maps from $\C^n$ to $\C$. Introduce the scalar product
\begin{gather}\label{eq:scal_product_bargm}
\langle f, g \rangle = (2\pi)^{-n}\int_{\C^n} e^{-|z|^2} f (z)\, \con{g (z)} \; d \mu_n (z), \qquad f,\; g \in \D (\C^n)
\end{gather}
where $|z|^2 = \sum_{i=1}^n |z_i|^2$ and $ \mu_n $ is the measure $\prod_{i=1}^n dz_i d\con{z}_i$. The family $((\al !)^{-\frac{1}{2}} \con{z}^\al, \, \al \in
\N^n)$ is an orthonormal basis of $\D (\C^n)$. We will also need the decomposition into even and odd functions
\begin{gather}\label{eq:paire_impaire}
\D (\C^n) = \D^+ (\C^n) \oplus \D ^{-} (\C^n)
\end{gather}
where $ \D^+ (\C^n) $ is spanned by the $\con{z}^\al$ with $|\al| = \sum \al(i)$ even and $\D ^{-} (\C^n) $ by the $\con{z}^\al$ with $|\al|$ odd.

Let $\symb (\C^n)$ be the space of endomorphisms\footnote{
In the whole paper, the endomorphisms are vector space endomorphisms.
}
$s $ of $\D(\C^n)$ such that $s (\con{z}^\al) = 0$ except for a finite number of $\al \in
\N^{n}$. We claim that $\symb (\C^n)$ is closed under product and taking adjoint. To see that, simply observe that $\symb (\C^n)$ is the space of endomorphisms having a matrix in the basis $(\con{z}^\al)$ whose almost all entries are equal to zero. Notice as well that the family $(\Uu_{\al, \be}, \;
\al, \be \in \N^n)$ of $\symb (\C^n)$ defined by
\begin{gather}\label{eq:def_Ualphabeta}
\begin{split}
\Uu_{\al \be } \bigl((\be !)^{-\frac{1}{2}} \con{z}^\be \bigr) = (\al !)^{-\frac{1}{2}}
\con{z}^{\al}, \\ 
\Uu_{\al \be} (\con{z}^{\ga}) = 0, \quad \forall \ga \in \N^n
\setminus \{ \be \}
\end{split}
\end{gather}
is a vector space basis of $\symb (\C^{n})$. And we have
\begin{gather}\label{eq:rel_U}
\Uu_{\al \be} \circ \Uu_{\tilde \al \tilde \be } = \delta_{\be \tilde{\al}}
\Uu_{\al \tilde \be}, \qquad \Uu_{\al \be}^* =
\Uu_{\be \al}
\end{gather}
for all $\al$, $\be$, $\tilde \al$, $\tilde \be$ in $\N^n$.

Each element $s \in \symb (\C^n)$ can be written in a block matrix $s =
(\begin{smallmatrix} s_{++} & s_{-+}
\\ s _{+-} & s_{--}
\end{smallmatrix})$ in the decomposition~\eqref{eq:paire_impaire}, which leads to a decomposition into even/odd endomorphisms
\begin{gather}\label{eq:symb+_-}
\symb (\C^n) = \symb^{+} (\C^n) \oplus \symb^{-} (\C^n)
\end{gather}
where $s \in \symb^{+} (\C^n)$ iff $s_{-+} = s_{+-} =0$, and $s \in
\symb^{-} (\C^n)$ iff $s_{++} = s_{--} = 0 $. Observe that $\Uu_{\al\be}$ has the same parity as $|\al | + | \be|$. Furthermore
\begin{gather}
\symb^{\ep} (\C^n)\, \cdot \, \symb^{\ep'} (\C^n) \subset
\symb^{\ep \ep'} (\C^n)
\end{gather}
for any $\ep, \ep' \in \{ 1, -1 \}$.


\subsubsection{Extension to vector bundles}
In the previous definitions, we can replace $\C^n$ by any $n$-dimensional Hermitian vector space $\EE$. We denote by $\D (\EE)$ the space of antiholomorphic polynomial maps $\EE \rightarrow
\C$. Choosing an orthonormal basis $(e_i)$ of $\EE$, we can identify $\EE$ with $\C^n$ and then define the scalar product of $\D (\EE)$ by the formula~\eqref{eq:scal_product_bargm}. Since the weight $|z|^2$ and the measure $d
\mu_n$ are invariant by unitary change of coordinates, the resulting scalar product of $\D (\EE)$ is independent of $(e_i)$. Similarly, we define the subspace $\symb (\EE)$ of the space of endomorphisms of $\D (\EE)$ and associated to the basis $(e_i)$ of $\EE$, we have a basis $(\Uu_{\al,\be},\;
\al, \be \in \N^n)$ of $\symb (\EE)$. The decompositions into even/odd elements are defined and denoted as for $\C^n$ by
\begin{gather}
\D (\EE) = \D^+ (\EE) \oplus \D^{-} (\EE), \qquad \symb (\EE) = \symb ^{+} (\EE) \oplus
\symb^{-} (\EE).
\end{gather}

We can extend all these constructions to vector bundles. Let $\EE \rightarrow M$ be a Hermitian vector bundle with rank $n$. Define the infinite-dimensional vector bundles $\D (\EE)$ and $\symb (\EE)$ over $M$ with fibers $\D (\EE)_x = \D (\EE_x)$ and $\symb (\EE)_x = \symb (\EE_x)$. Later, we will choose for $\EE$ the complex tangent bundle of an almost-complex manifold, and we will construct operator whose symbols are smooth sections of $\symb (\EE) \otimes A$, where $A$ is an auxiliary vector bundle. Since the bundle $\symb(\EE)$ has infinite rank, let us make precise the definition of its smooth sections: a section $s \in \Ci (M,
\symb (\EE) \otimes A)$ is a family $(s(x) \in \symb (\EE_x) \otimes A_x, \; x \in M)$ such that for any orthonormal frame $(e_i)$ of $\EE$ and $(a_j)$ of $A$ over the same open set $U$ of $M$, if $(\Uu_{\al,
\be}(x))$ is the basis of $\symb (\EE_x)$ associated to the basis $(e_i(x))$, then
\[
s (x) = \sum \la_{\al,\be,j} (x) \, \Uu_{\al,\be} (x) \otimes a_j(x), \qquad
x \in U
\]
where the $\la_{\al, \be, j} $ are smooth functions on $U$, almost all equal to zero.


\subsection{Operators} \label{sec:operators}

\subsubsection{Schwartz kernel}

Consider a compact symplectic manifold $(M, \om)$ with a compatible almost-complex structure $j$ and a prequantum bundle $L \rightarrow M$. Assume we have two auxiliary Hermitian vector bundles $A$ and $B$ over $M$. The dimension of $M$ is $2n$ with $n \in \N$.


We will define a space $\lag (A, B)$ consisting of families of operators
\begin{gather}\label{eq:P_k_family}
\left(P_k : \Ci \left(M, L^k \otimes A\right) \rightarrow \Ci \left(M, L^k
\otimes B\right), \; k \in \N \right)
\end{gather}
having smooth Schwartz kernels satisfying some conditions. Let us first recall some standard definitions and notations.

 We denote by $\con{A}$ the conjugate bundle of $A$ and by $A \boxtimes B$ the external tensor product of $A$ and $B$.  The Schwartz kernel of $P_k$ is the section $K_k$ of $(L^k \otimes B) \boxtimes (\con{L}^k \otimes \con{A})$ such that
\[
(P_k f)(x) = \int_M K_k (x,y)\cdot f(y) \; \mu_M (y), \qquad \forall \;
f \in \Ci \left(M, L^k \otimes A\right)
\]
where the $\cdot$ stands for the scalar product $(\con{L}_y^k \otimes
\con{A}_y) \times (L_y^k \otimes A_y) \rightarrow \C$, and $\mu_M =
\om^n/n!$. We will denote the operator and its Schwartz kernel by the same letter, hoping it is not too confusing.

Since $L$, $A$ and $B$ are Hermitian bundles, the bundle $(L^k \otimes B) \boxtimes (\con{L}^k
\otimes \con{A})$ has a natural metric, so the pointwise norm $|P_k (x,y)|$ is well-defined. For any $N \in \N$, we will say that $(P_k)$ is in $\bigo (k^{-N})$ on an open set $U$ of $M^2$ if $|P_k (x,y)| =
\bigo (k^{-N})$ for $(x,y) \in U$ with a $\bigo$ uniform on any compact subsets of $U$. We say that $(P_k)$ is in $\bigo (k^{-\infty})$ on $U$ if $(P_k)$ is in $\bigo (k^{-N})$ on $U$ for any $N$.

We will also use the uniform norm $\| P_k \| = \sup \|P_k (f) \| / \| f\|$ with respect to the usual $L^2$ norms of section: $\|f\|^2 = \int_M |f(x)|^2
\; d\mu_M (x)$.


\subsubsection{Definition of $\lag (A,B)$}

By definition, a family $(P_k)$ as in~\eqref{eq:P_k_family} belongs to $\lag (A, B)$ if each $P_k$ has smooth Schwartz kernel satisfying for any $N$
\begin{gather}\label{eq:exp_Lag_bis}
P_k (x,y) = \Bigl(\frac{k}{2\pi} \Bigr)^n E^{k}(x,y) \sum_{\substack{\ell\, \in\Z, \\
\ell + m (\ell)\,\leqslant\,N}} k^{-\frac{\ell}{2}} b_{\ell}(x,y) + \bigo \left(k^{n - \frac{N+1}{2}}\right)
\end{gather}
where
\begin{enumerate}
\item $E$ is a section $L \boxtimes \con{L}$ such that $|E (x,y)| < 1$ for any $x\neq y$, and for any $y \in M$, the section~$E_y (x) = E(x,y)$ of $L
\boxtimes \con{L}_y$ satisfies $E_y (y) = u
\otimes \con{u}$ for any $u \in L_y$ with $|u|=1$, $(\nabla E_y)(y) =0$ and
\[
\left(\nabla_\xi \nabla_\eta E_y\right) (y) = -\left(\tfrac{i}{2} \om (\xi, \eta) +
\tfrac{1}{2} \om (\xi, j \eta)\right) E_y (y), \quad \forall \; \xi,\eta
\in T_yM.
\]
\item $m : \Z \rightarrow \N \cup \{ \infty \}$ is such that $\{ \ell; \; \ell + m (\ell) \leqslant N \}$ is finite for any $N$ and $\ell + m (\ell)
\geqslant 0$ for any $\ell$. Moreover, for any $\ell$, $b_{\ell} $ is a section of $B \boxtimes \con{A}$ vanishing to order $m(\ell)$ along the diagonal.
\end{enumerate}
As already mentioned in the introduction, the analytic meaning of the expansion~\eqref{eq:exp_Lag_bis} is not obvious, nevertheless we postpone the explanations to Section~\ref{sec:global-expansion}, cf. Lemma~\ref{lem:estimate} and Lemma~\ref{lem:estim_expansion}. The existence of $E$ will be proved in Section~\ref{sec:section-ee}, it is not unique,   any section~$E$ satisfying the stated conditions can be used for the expansion~\eqref{eq:exp_Lag_bis}, but the coefficients $b_{\ell}$ depend on the choice of $E$, cf. Lemma~\ref{lem:changement_E}.  

By rescaling the coordinates transverse to the diagonal by a factor $k^{\frac{1}{2}}$, we can write the expansion~\eqref{eq:exp_Lag_bis} in the following alternative way. To simplify the statement, we assume first that $A$ and $B$ are the trivial line bundle $\C_M := M \times \C$ and let $\lag (
\C) := \lag (\C_M, \C_M)$.

\begin{prop}\label{prop:un_de_plus}
Let $(P_k)$ be an operator family $(P_k)$ of the form~\eqref{eq:P_k_family} with smooth Schwartz kernels. Then $(P_k)$ belong to $\lag (\C)$ if and only if the Schwartz kernel family is in $\bigo (k^{-\infty})$ on $M^2
\setminus \op{diag} M$ and for any coordinate chart $U \subset M$ and unitary frame $t : U \rightarrow L$, for any $N \in \N$, we have over $U^2$ that
\begin{gather}\label{eq:expansion_kernel}
P_k (x + \xi, x) = \left(\frac{k}{2\pi} \right)^n e^{- k\varphi (x, \xi) }
\sum_{p =0 }^{N} k^{- \frac{p}{2}} a_{p} \left(x, k^{\frac{1}{2}} \xi \right) + \bigo \left(k^{n - \frac{N+1}{2}} \right)
\end{gather}
where we have identified $L_{x+\xi}^k \otimes \con{L}_{x}^k \simeq \C$ by using $t$ and
\begin{itemize}
\item[-] $\varphi(x, \xi) = - i (\sum_{i=1}^{2n} \al_i (x)\xi_i + \frac{1}{2}\sum_{i,j=1}^{2n} (\partial_{x_i}\al_j)(x) \xi_i \xi_j) + \frac{1}{4} |\xi|_x^2 $, with $\al =
\sum \al_i dx_i \linebreak\in \Om^1 (U, \R)$ the connection one-form defined by $\nabla t = \frac{1}{i} \al \otimes t$.
\item[-] $a_{p} (x,\xi) \in \C$ depends polynomially on $\xi$, meaning that for some $d(p) \in \N$, $a_p (x,\xi) =
\sum_{|\al |\leqslant d(p)} a_{p,\al} (x) \xi^\al$ with smooth coefficients $a_{p, \al}$.
\end{itemize}
\end{prop}

The proof is postponed to Section~\ref{sec:global-expansion}. Since the real part of $\varphi (x,\xi)$ is $\frac{1}{4} |\xi|_x^2$, we have for any $p$
\[
e^{-k \varphi (x, \xi)} a_p \left(x, k^\frac{1}{2} \xi \right) = \bigo
(1).
\]
So the $p^{\rm th}$ summand in~\eqref{eq:expansion_kernel} is in $\bigo (k^{n -\frac{p}{2}})$ and the expansion is meaningful. In the case where $A$ and $B$ are general vector bundles, we introduce frames of $A$ and $B$ on $U$, so that the Schwartz kernel $P_k$ on $U^2$ becomes a $\C^r$-valued functions with $r =(\op{rank} A)(\op{rank} B)$, and we have the same characterization with $\C^r$-valued coefficients $a_p$.


The advantage of the expansion~\eqref{eq:expansion_kernel} is that its analytical meaning is more transparent, the drawback is that it depends on local choices (coordinates, frames, rescaling $k^{\frac{1}{2}} \xi$) whereas the expansion~\eqref{eq:exp_Lag_bis} is global.

\subsubsection{Properties of $\lag (A,B)$}

$\lag (A, B)$ has a natural filtration defined as follows. For any $q \in
\N$, $\lag _q (A, B)$ is the subspace of $\lag (A,B)$ consisting of the operators such that the local expansions~\eqref{eq:expansion_kernel} hold with a sum starting at $p = q$, that is the coefficients $a_0,\, \ldots\,a_{q-1}$ are zero.

\begin{prop}\label{prop:lag}\leavevmode
\begin{enumerate}
\item\label{prop2.2.1} $\lag _q (A, B) = k^{-\frac{q}{2}}
\lag (A, B)$ and if $q \geqslant q'$, then $\lag_q(A,B) \subset \lag_{q'} (A,B)$.
\item\label{prop2.2.2} For any $(P_k)$ in $\lag (A, B)$,
\begin{align*}
&(P_k) \in \lag_q(A,B) \Leftrightarrow \| P_k \| = \bigo \left(k^{-\frac{q}{2}}\right) \\
&\Leftrightarrow \text{the Schwartz kernel family of $(P_k)$ belongs to }
\bigo \left(k^{n-\frac{q}{2}}\right)
\end{align*}
\item\label{prop2.2.3} $\lag_{\infty} (A,B) := \bigcap_q \lag_q(A,B)$ consists of the families~\eqref{eq:P_k_family} with a smooth Schwartz kernel in $ \bigo (k^{-\infty})$.
\item\label{prop2.2.4} for any sequence $(P_q)_{q\,\in\,\N}$ of $\lag (A,B)$ such that $P_q \in \lag_q(A, B)$ for any $q$, there exists $P \in \lag (A,B)$ satisfying $P = \sum_{p=0}^q P_p$ modulo $\lag_{q+1} (A,B)$ for any $q$.
\end{enumerate}
\end{prop}

We will now describe the quotients $\lag_q (A,B)/ \lag_{q+1} (A,B)$ by using the material introduced in Section~\ref{sec:symbol-spaces}. Since $M$ has an almost complex structure $j$ compatible with $\om$, the tangent bundle $TM$ is a complex vector bundle with a Hermitian metric, which defines our bundle $\symb (M) := \symb (TM)$.

\begin{theo}\label{theo:lag}
For any $q \in \N$, there exists a linear map
\[
\si_q: \lag_q (A,B) \rightarrow \Ci (M, \symb (M) \otimes \op{Hom}
(A,B))
\]
which is onto and has kernel $\lag_{q+1} (A,B)$. Furthermore, the
following holds for any $ P \in \lag_q (A,B)$:
\begin{enumerate}
\item\label{theo2.3.1} $\si _q (P) = \si_0 (k^{\frac{q}{2}} P)$.
\item\label{theo2.3.2} For any $f \in \Ci (M, \op{Hom} (B, C))$, $(f \circ P_k)$ belongs to $\lag_q (A,C)$ and $\si_q(f \circ P) = f \circ \si_q (P)$. For any $g \in \Ci (M, \op{Hom} (C, A))$, $(P_k
\circ g)$ belongs to $\lag_q (C,B)$ and $\si_q (P \circ g) = \si_q (P)
\circ g$.
\item\label{theo2.3.3} $P^* $ belongs to $\lag_q (B,A)$ and $\si (P^*) = \si (P)^*$.
\item\label{theo2.3.4} For any $P' \in \lag_{q'}(B,C)$, $P' \circ P$ belongs to $\lag_{q'+ q }(A,C)$ and
\[
\si_{q'+ q} (P' \circ P) = \si_{q'} (P') \circ \si_q (P).
\]
\item\label{theo2.3.5} The Schwartz kernel of $P_k$ on the diagonal satisfies
\[
P_k (x,x) = \frac{k^{n-\frac{q}{2}}}{ (2\pi)^n} \Bigl[\op{tr} (\si_q(P)(x))
+ \bigo \left(k^{-\frac{1}{2}}\right) \Bigr]
\]
where $\op{tr}$ is the map $\symb (T_xM) \otimes \op{Hom} (A_x,B_x) \rightarrow (L_x \otimes \con{L}_x)^k \otimes B_x \otimes
\con{A}_x \simeq \op{Hom} (A_x,B_x)$ sending $s \otimes f$ to $(\op{tr} s) f$.
\item\label{theo2.3.6}  The operator norm of $P_k$ satisfies
\[
\| P_k \| = k^{-\frac{q}{2}} \left(\sup_{x\,\in\,M} \left\|\si_q (P) (x) \right\| +
\op{o} (1) \right)
\]
where $\|\si_q (P) (x) \|$ is the operator norm for the norm of $\D (T_xM)$ corresponding to the scalar product~\eqref{eq:scal_product_bargm}. 
\end{enumerate}
\end{theo}



Let us explain how is defined the symbol map $\si_0$ for $A= B = \C_M$. Consider $P \in \lag (A,B)$ and the local expansion~\eqref{eq:expansion_kernel}. We view $(x,\xi)$ as a tangent vector of $M$, that is $
\xi \in T_xM$, so we consider $a_0 (x, \cdot)$ as a polynomial of $T_xM$. Then it is not obvious but nevertheless true that this polynomial does not depend on the choice of the coordinate chart $U$ and the unitary frame $t$. To compare, the coefficients $a_p$ in~\eqref{eq:expansion_kernel} with $p \geqslant 1$ do depend on the choice of the coordinates and the frame of $L$.

To pass from $a_0(x,\cdot)$ to the symbol of $P$ at $x$, we first choose a unitary frame $(e_i)$ of $T_xM$. So $T_xM \simeq \C^n$ by sending $\xi = \sum z_i e_i(x) $ to $z(\xi)= (z_i)$. We also have a basis $\Uu_{\al, \be}(x)$ of $\symb (T_xM) $ defined in~\eqref{eq:def_Ualphabeta}. Then
\[
\si_0(P) (x) = \sum f_{\al, \be} (x) \Uu_{\al,\be} (x) \Leftrightarrow
a_0 (x,\xi) = \sum f_{\al, \be} (x) \Ppp_{\al, \be} (z(\xi))
\]
where we use the polynomials $ \Ppp_{\al, \be} (z) = (\frac{1}{\al! \be !}) ^{1/2} (
\partial_z - \con{z})^{\al} z^\be.$ These polynomials form a basis of $\C [z,
\con{z}]$, cf. proof of Proposition~\ref{prop:tilde_symb_and_op}.

\begin{exam}\label{exemple:projecteur}
Choose a connection of $A$ and let $\Delta_k$ be the Laplacian
\begin{gather}\label{eq:laplacien_fibre_auxiliaire}
\Delta_k = \tfrac{1}{2} \left(\nabla^{L^k
\otimes A}\right)^* \nabla^{L^k \otimes A} : \Ci\left(L^k \otimes A\right) \rightarrow \Ci \left(L^k \otimes A\right).
\end{gather}
For any $m \in \N$, let $\Pi_{m,k}$ be the spectral projector
\begin{gather}\label{eq:projecteur_fibre_auxiliaire}
\Pi_{m,k} := 1_{\left[m - \frac{1}{2}, m +
\frac{1}{2}\right]} \left(k^{-1} \Delta_k\right).
\end{gather}
By~\cite{oim_copain}, the family $(\Pi_{m,k})$ belongs to $\lag(A,A)$, its $\si_0$-symbol at $x$ is $\pi_m(x) \otimes \op{id}_{A_x}$ where $\pi_m(x)$ is the projector of $\D (T_xM)$ onto the subspace $\D_m(T_xM)$ of homogeneous degree $m$ polynomials.

Since $\pi_m (x) = \sum_{|\al|=m} \rho_{\al, \al}(x)$, if the auxiliary bundle $A$ is trivial, the corresponding function $a_0$ is
\begin{gather}\label{eq:a_0laguerre}
a_0(x, \xi) = \sum_{|\al| =m} p_{\al,\al}(z(\xi)) = Q_m^{(n-1)} \left(|z(\xi)|^2\right)
\end{gather}
where $Q_m^{(p)} $ is the Laguerre polynomial $Q^{(p)}_m (x) = \frac{x^{-p}}{m!}(\frac{d}{dx} -1) ^m x ^{m+p}$. The second equality in~\eqref{eq:a_0laguerre} follows from $p_{m,m} (z) = Q_m^{(0)} (|z|^2)$ and the identity
\[
Q_m^{(n-1)} (x_1+ \ldots + x_n) =
\sum_{|\al| = m } Q_{\al(1)}^{(0)} (x_1) \ldots Q_{\al(n)}^{(0)} (x_n).
\]
Actually we won't use the expression in terms of Laguerre polynomials, what really matters is the fact that $\si_0 (\Pi_m) (x)$ is the orthogonal projector onto $\D _m(T_xM)$.
%\qed
\end{exam}

The definition of the symbol map $\si_0$ is motivated by Theorem~\ref{theo:lag} and its efficiency in the proofs of Section~\ref{sec:projectors-Toeplitz}. But this definition does not explain why it is natural to associate to $P \in \lag (A,B)$ an endomorphism of $\D (T_xM)$. A first explanation is provided by the following construction of peaked sections. A deeper reason will be provided later in Section~\ref{sec:landau-levels-cn}.

We still assume that $A = B = \C_M$ to simplify the exposition. Let $x \in M$ be a base point, with a coordinate chart $U$ at $x$ and a unitary frame $t: U \rightarrow L$. Let $\psi \in \Ci_0(U)$ be equal to $1$ on a neighborhood of $x$. To any $f \in \D (T_xM)$, we associate a family $\Phi_k^f \in \Ci (M, L^k)$ defined by
\[
\Phi_k^f (x+\xi) = \left(\frac{k}{2\pi} \right)^{\frac{n}{2}} e^{-k
\varphi(x,\xi)} \, f\left(k^{\frac{1}{2}} \xi\right) \, \psi (x+\xi) \, t
^k(x+\xi), \qquad k \in \N
\]
where $\varphi$ is the same function as in~\eqref{eq:expansion_kernel}.  Let $\| f \| = \sqrt{(f,f)}$ be the norm associated to the scalar product~\eqref{eq:scal_product_bargm}. 

\begin{prop}\label{prop:peaked-sections}
For any $f \in \D (T_xM)$,
\begin{gather}\label{eq:peaked_1}
\| \Phi^f_k \| = \| f \| + \bigo \left(k^{-1/2}\right),
\end{gather}
and for any $P \in \lag (\C_M, \C_M)$,
\begin{gather}\label{eq:peaked_2}
P_k \Phi^f_k = \Phi^g_k + \bigo \left(k^{-1/2}\right)
\end{gather}
where $g = \si_0(P)(x) \cdot f$.
\end{prop}

By~\eqref{eq:peaked_1} the map sending $f$ into $(\Phi_k^f)$ is injective, so~\eqref{eq:peaked_2} characterizes the symbol $\si_0 (P) (x)$. For a more general result with auxiliary bundles $A$, $B$ and the estimates of the scalar product of peaked sections, cf. Proposition~\ref{prop:peaked-sections_++}.

We say that an element $P$ in $\lag (A,B)$ is even (resp. odd) if the expansion~\eqref{eq:exp_Lag_bis} holds with $b_{\ell } = 0$ for any odd $\ell \in \Z$ (resp. even).

\begin{lemm}\label{lem:parite}
For any $P$ in $\lag (A,B)$, $(P_k)$ is even (resp. odd) if and only if in the local expansions~\eqref{eq:expansion_kernel}, every polynomial $a_p(x,\cdot)$ has the same (resp. the opposite) parity as $p$.
\end{lemm}


Denote by $\lag^+ (A, B)$ and $\lag^- (A, B)$ the subspaces of even and odd elements respectively.

\begin{theo}\label{theo:parity}
We have
\begin{enumerate}
\item\label{theo2.7.1} $ \lag (A, B) = \lag^+ (A, B) + \lag^- (A, B)$, $\lag^+ (A,B)
\cap \lag ^{-} (A,B) = \lag_{\infty} (A,B)$.
\item\label{theo2.7.2} $\lag ^{\ep} (A,B) \cdot \lag^{\ep'} (B,C) \subset \lag ^{\ep \ep'} (A, C) $ for any choice of signs $\ep$, $\ep'$.
\item\label{theo2.7.3} $\si_q(\lag_q(A,B) \cap \lag^{\ep} (A,B)) = \Ci (M,
\symb^{\ep (-1)^q}(M) \otimes \op{Hom} (A,B))$.
\end{enumerate}
\end{theo}


The proofs of Proposition~\ref{prop:lag}, Theorem~\ref{theo:lag} and Theorem~\ref{theo:parity} are postponed to Section~\ref{sec:proofs-results}. The proof of Lemma~\ref{lem:parite} is at the end of Section~\ref{sec:global-expansion}.


\subsection{Comparison with earlier works}

The expansions~\eqref{eq:exp_lag}, \eqref{eq:expansion_kernel} or similar versions appeared in the literature~\cite{MaMa,oim_op,ShZe} to describe Bergman kernels of ample line bundles and their symplectic generalizations as well as the associated Toeplitz operators. In a more general context, the Boutet de Monvel--Guillemin theory~\cite{BoGu} is built on two classes of operators: Hermite operators and Fourier integral operators respectively. The spaces $\lag^{+} (A,B)$ may be viewed as an intermediate choice in the semi-classical setting.


In~\cite{oim}, we considered a subalgebra of $\lag^{+} (A,A)$, denoted by $\clA(A)$, consisting of operators having an expansions~\eqref{eq:expansion_kernel} in which each $a_p(x,\cdot)$ has degree $\leqslant \frac{3}{2} p$. For our applications in this paper, it is necessary to consider the larger spaces $\lag (A,B)$, because our generalized projectors $\Pi_m$ and unitary equivalences do not belong to $\clA (A)$. More precisely, only the projector corresponding to the first Landau level belongs to $\clA (A)$.

Theorem~\ref{theo:lag} is a generalization of similar results for $\clA(A)$ established in~\cite{oim}, and surprisingly the proofs are somehow easier in this new generality. However, a crucial difference with~\cite{oim} relies in the symbols. Roughly, the symbols of the elements of $\clA(A)$ were defined directly as the polynomials $a_0 (x,\cdot)$. This had the advantage that it is easier to pass from the Schwartz kernel of the operator to the symbol. The drawback of this definition is that the product for these symbols is given by the mysterious formula
\begin{gather}\label{eq:old_prod}
(u \star v) (x, z, \con{z}) =\bigl[\exp (\Box) \left(u \left(x, - \zeta, \con{z} - \con {\zeta}\right) v \left(x, z+ \zeta, \con{\zeta}\right)\right)\bigr]_{\zeta = \con{\zeta} =0 }
\end{gather}
where $\Box = \sum \partial^2 / \partial \zeta^i \partial \con{\zeta}^i$. This was actually tractable for what we did in~\cite{oim} because our main interest was the projector $\Pi$ for the first Landau level, with symbol $\si_0(\Pi)=
\Uu_{00}$ and associated function symbol $u(x,\xi) =1$. But for our new projectors whose symbols are typically a sum of the $\Uu_{\al\al}$'s, it is essential to work with the symbol in $\symb (M)$. For instance, it is even not obvious how to recover the relations~\eqref{eq:rel_U} from the product~\eqref{eq:old_prod}.



\section{Projectors of \texorpdfstring{$\lag (A)$}{L(A)} and Toeplitz operators} \label{sec:projectors-Toeplitz}

In this section, we consider an auxiliary Hermitian vector bundle $A$ with arbitrary rank. We denote by $\lag (A): = \lag (A,A)$ the associated algebra and by $\lag^+(A)$ the subalgebra consisting of even elements. The symbols of operators of $\lag (A)$ are sections of $\symb (M) \otimes \op{End} A$. We will view $\symb (M)_x \otimes \op{End} A_x$ as a subspace of $\op{End} (\D (T_xM)
\otimes A_x)$.



Let $F$ be a subbundle of $\D_{\leqslant m}(TM)
\otimes A$ for some $m \in \N$, where $\D_{\leqslant\,m} (T_xM)$ is the subspace of $\D(T_xM)$ of polynomial with degree $\leqslant m$. We assume that $F$ has a definite parity, that is $F \subset \D ^\ep (TM) \otimes A$ with $\ep \in \{\pm 1\}$. Associated to $F$ is the section $\pi$ of $\symb (M) \otimes \op{End} A$ such that $\pi(x)$ is the orthogonal projector of $\D(T_xM) \otimes A_x$ onto $F_x$ at each point $x \in M$. The content of the following subsections is:
\begin{itemize}
\item[-] Subsection~\ref{sec:constr-proj}: we construct a selfadjoint projector $\Pi \in \lag ^+(A)$ with symbol $\pi$.
\item[-] Subsection~\ref{sec:toeplitz-algebra}: we study the Toeplitz algebra $ \clT = \{ \Pi P \Pi, \; P \in \lag^{+} (A) \} $
\item[-] Subsection~\ref{sec:unitary-equivalence}: we prove $(
\op{Im} \Pi_k)$ is isomorphic with any quantization of $(L,M)$ twisted by $F$, and deduce that the dimension of $ \op{Im} \Pi_k$ is the Riemann--Roch number of $L^k \otimes F$ when $k$ is sufficiently large.
\end{itemize}

A possible choice for $F$ is $F = \D_m (TM)\otimes A $ where $m\in \N$ and $\D_m (T_xM)$ is the subspace of $\D (T_xM)$ consisting of homogeneous polynomials with degree $m$. As explained in Example~\ref{exemple:projecteur}, the projector $\Pi_{m,k}$ onto the $m^{\rm th}$ Landau level $\Hilb_{m,k}$ belongs to $\lag^+ (A)$ and has symbol the projector onto $F$, so it can be used as the projector $\Pi$. Theorems~\ref{theo:dim_landau}, \ref{theo:Toeplitz_Landau} and~\ref{theo:ladder} will mainly follow from the results in Sections~\ref{sec:toeplitz-algebra} and~\ref{sec:unitary-equivalence}.

By~\cite{oim_copain}, the spectral projectors of Laplacians with a magnetic field, not necessarily constant but still satisfying some convenient assumptions, give other instances of projectors in $\lag ^+(A)$.

Another choice for $F$ is $F = \D_0 (TM) \otimes A$ where $\D_0 (T_xM) =\C$ is the subspace of $\D (T_xM)$ of constant polynomials. The corresponding quantum space and Toeplitz algebra is the quantization of $(M,L)$ twisted by $A$.

A last example is the Spin-c Dirac quantization twisted by an auxiliary bundle $B$. In this case, $A = S \otimes B$ where $S$ is the spinor bundle  $\bigoplus \bigwedge^k (T^*M)^{0,1} $ and $F = \D_0(TM)
\otimes \bigwedge^0 (T^*M)^{0,1} \otimes B$.  This example will be used to compute the dimension of our quantum spaces from the Atiyah--Singer theorem.



\subsection{Construction of the projector}\label{sec:constr-proj}
Let $\chi : \R \rightarrow \R$ be defined by $\chi (x) =1$ if $x \geqslant
\frac{1}{2}$ and $\chi (x) =0$ otherwise. If $P$ is a bounded self-adjoint operator of a Hilbert space $\Hilb$, then using the functional calculus for Borel bounded functions, we define a new bounded operator $\chi (P)$ of $\Hilb$, cf. as instance~\cite[Theorem~VII.2]{ReSi}. Since $\chi$ is real valued and $\chi^2 = \chi$, $\chi (P)$ is a self-adjoint projector.

\begin{theo}\label{theo:constr-proj}
Let $P \in \lag (A)$ be self-adjoint and having symbol $\si_0 (P) = \pi$. Then $\chi(P)$ belongs to $\lag (A)$ and $\si_0(\chi (P)) = \pi$. If furthermore $P \in \lag ^+ (A)$, then $\chi (P)$ is in $\lag ^+ (A)$.
\end{theo}

 An operator $P$ satisfying the assumptions exists by the surjectivity of $\si_0$, cf. Theorem~\ref{theo:lag}.  Theorem~\ref{theo:constr-proj} holds without the assumption that $F$ has a definite parity. When $F$ does have a definite parity, $\pi$ is even, so we can choose $P \in \lag^{+}(A)$ with symbol $\pi$.

\begin{proof}
To prove that the $\chi (P_k)$'s have smooth kernels, we will use the following basic fact: let $Q$, $Q'$ be two operators with smooth kernels acting on $\Ci (M,A)$ and $Q''$ be a bounded operator of $L^2(M, A)$. Then $Q Q''Q'$ has a smooth kernel. This follows from the Schwartz theorem saying that the operators with smooth kernel are the operators which can be continuously extended $\clC^{-\infty} \rightarrow \Ci $. We will also need the following pointwise norm estimates: consider families of operator $Q_k, Q'_k : \Ci (M,L^k \otimes A) \rightarrow \Ci (M, L^k\otimes A)$ and $Q''_k : L^2 (M, L^k
\otimes A) \rightarrow L^2 (M, L^k \otimes A)$. Then by~\cite[Section~4.3]{oim}, if the Schwartz kernel families of $(Q_k)$ and $(Q'_k)$ are respectively in $\bigo (k^{-N})$ and $\bigo (k^{-N'})$, and the operator norms of $Q_k^{''}$ are in $\bigo (1)$, then the Schwartz kernel family of $Q_k Q''_k Q_k'$ is in $\bigo (k^{-(N +N)'})$.

Back to our problem, we can write $\chi(P_k) = P_k \tilde{\chi} (P_k) P_k$ with $\tilde{\chi} (x) = \chi(x) / x^2$. Since $ \tilde{\chi}(P_k)$ is bounded, this shows that $\chi (P_k)$ has a smooth kernel. This also shows that the Schwartz kernel family of $\chi (P)$ is in $\bigo (k^{2n})$. To improve this, observe that $Q = P^2 - P$ is in $\lag (A)$ and $\si _0 (Q) = \pi^2 - \pi =0$, so $\| Q_k \| = \bigo (k^{-1/2})$, which implies easily that $\frac{1}{2}$ is not in the spectrum of $P_k$ when $k$ is sufficiently large, cf.~\cite[Proposition~4.2]{oim}.

Now for $x \in \R \setminus \{ \frac{1}{2}\}$, $y = x^2 -x > -1/4 $ and we have
\[
\chi (x) = x + (1-2x) f \left(x^2 -x\right) \quad \text{with } \quad f (y) = \tfrac{1}{2}
\left(1- (1+ 4y)^{-1/2}\right)
\]
 For any $m \in \N$, write the Taylor expansion of $f$ at $0$ at order $m $ as follows:  $f (y) = \sum_{\ell =0}^m a_\ell y^\ell + y^{m+1} f_m (y)$ with $f_m \in \clC^0(]-\frac{1}{4},
\infty[, \R)$. Then
\begin{gather}\label{eq:toutestla}
\chi (P) = P + \sum_{\ell = 0 }^m a_\ell (1-2P) Q^\ell + (1 - 2 P) Q^{m+1} f_m (Q).
\end{gather}
Now $\si_0(Q)=0$ implies that $Q^\ell$ and $PQ^{\ell}$ belong both to $\lag_{\ell} (A)$. Furthermore, $\|f_m(Q_k) \| = \bigo (1)$. Since $Q^{m+1} f_m (Q) = Q^m f(Q) Q$ and similarly $PQ^{m+1} f_m (Q) = PQ^m f(Q) Q$, it follows from the preliminary observation that the Schwartz kernel family of $(1 - 2 P) Q^{m+1} f_m (Q)$ is in $\biginf (k^{2n -m})$. We can now conclude easily the proof from~\eqref{eq:toutestla} by choosing at each step sufficiently large value of $m$: first the Schwartz kernel family of $\chi (P)$ is in $\bigo (k^{-\infty})$ outside the diagonal and second the local expansions~\eqref{eq:expansion_kernel} hold.
\end{proof}


\subsection{Toeplitz algebra} \label{sec:toeplitz-algebra}

Choose a self-adjoint projector $\Pi \in \lag ^+ (A)$ with symbol $\pi$, which exists by Theorem~\ref{theo:constr-proj}. For any $k\in \N$, let $\Hilb_k = \op{Im} \Pi_k \subset \Ci (M, L^k \otimes A)$. Computing the trace of $\Pi_k$ by integrating its Schwartz kernel over the diagonal, we deduce from the last assertion of Theorem~\ref{theo:lag} that $\Hilb_k$ is finite dimensional and 
\begin{gather}\label{eq:estim_dim}
\op{dim} \Hilb_k \sim \left(\frac{k}{2 \pi} \right)^n (\op{rank} F) \;
\op{vol} (M, \om)
\end{gather}
As we will see later, when $k$ is sufficiently large, this dimension depends polynomially on $k$ and is a Riemann--Roch number, cf. Theorem~\ref{theo:dim_general}. 

We will now work with families of operators $(T_k \in \op{End} \Hilb_k$, $k \in
\N)$. Equivalently, we can consider that each $T_k$ acts on the larger space $\Ci (M, L^k \otimes A)$ and satisfies $\Pi_k T_k \Pi_k = T_k$. Define the space
\begin{gather}\label{eq:Toeplitz_Pi}
\clT = \left\{ T \in \lag^{+} (A) / \, \Pi T \Pi = T \right\}.
\end{gather}
For any $q \in \N$, set $\clT_q := k^{-q} \clT =
\lag_{2q} (A) \cap \clT $ and $\clT_{\infty} = \bigcap_q \clT_q $. Clearly,
\[
\clT_\infty
\subset \clT_q
\subset \clT_p \subset \clT
\]
when $q \geqslant p$.

\begin{theo}\label{theo:toeplitz}\leavevmode
\begin{enumerate}
\item\label{theo3.2.1} For any $T \in \clT $, $T$ belongs to $\clT_q $ iff $\| T_k \| = \bigo (k^{-q})$. Furthermore, $\clT_{\infty}$ consists of the families $(T_k \in \op{End} (\Hilb_k), \, k \in \N^*)$ such that $\| T_k \| = \bigo(k^{-N})$ for any $N$.
\item\label{theo3.2.2} $ \clT $ is closed under composition and taking adjoint: $ \bigl(\clT_q \bigr)^* = \clT_q$ and $\clT_q \cdot \clT_p \subset \clT_{q+p} $ for any $q$ and $p$.
\item\label{theo3.2.3} For any $q$, there exists a linear map $\tau_q : \clT_q \rightarrow
\Ci (M, \op{End} F) $, which is onto, has kernel $\clT_{q+1} $ and is determined by $\si_{2q} (T) = \tau_q(T)
\pi$. Furthermore, if $P \in \clT_q $ and $Q \in
\clT_p $, then
\begin{equation}
\begin{split}  \tau_q(P) &= \tau_0(k^{q}P), \qquad \tau_q(P ^*) =
\tau_q(P)^*, \\ 
\tau_q (P) \tau_p(Q) &= \tau_{q+p} (PQ) \\
 \| P_k \| &= k^{-q} \bigl(\sup \left\{ \| \tau_q (P)_x \|, \; x \in M \right\} + \op{o} (1) \bigr),
\end{split}
\end{equation}
and the restriction to the diagonal of the Schwartz kernel of $P_k$ satisfies
\begin{gather}\label{eq:trace_diag}
P_k (x,x) = \frac{k^{n-q}}{(2\pi)^n} \Bigl[\op{tr} \left(\tau_q (P) (x)\right) +
\bigo \left(k^{-1}\right) \Bigr].
\end{gather}
\end{enumerate}
\end{theo}

Let us give more details on the equation $\si_{2q} (T) = \tau_q(T)
\pi$ defining the symbol map $\tau_q$. Recall that $\pi(x)$ is the orthogonal projector of $\D (T_xM) \otimes A_x$ onto $F_x$.

Then for an endomorphism $s$ of $F_x$, we define $s
\pi(x) \in \symb (T_xM) \otimes \op{End} A_x $ as the endomorphism of $\D (T_xM) \otimes A_x$ sending $\psi $ into $s (\pi(x) \psi)$.

\begin{proof}
\eqref{theo3.2.1} The first assertion follows from Part~\eqref{prop2.2.2} of Proposition~\ref{prop:lag}. To establish the second assertion, we deduce from the first part of the
proof of Theorem~\ref{theo:constr-proj} that if a family $(P_k \in \op{End} (\Hilb_k))$ satisfies $\| P_k \| = \bigo (k^{-\infty})$, then its Schwartz kernel is in $\bigo (k^{-\infty})$ because $ \Pi_k P_k \Pi_k = P_k$ and the Schwartz kernel of $\Pi_k$ is in $\bigo (k^{n})$.

Parts~\eqref{theo3.2.2} and~\eqref{theo3.2.3} follow from Theorem~\ref{theo:lag} and the fact that $\lag ^{+} (A)$ is a subalgebra of $\lag (A)$ by theorem~\ref{theo:parity}. To define $\tau_q (P)$, simply observe that $\Pi P \Pi = P$ implies that $\pi \si_{2q} (P) \pi = \si_{2q}(P)$, so we can write $\si_{2q} (P) = \tau_q(P) \pi$ with $\tau_q (P)$ a section of $ \op{End} F$. The map $\tau_q$ is onto because for any section $s$ of $\op{End} F$, there exists $P$ in $ \lag_{2q} (A)$ with $\si_{2q} (P) = s \pi$. Since $\pi (s \pi) \pi = s
\pi$, we have $\si_{2q} (\Pi P \Pi) = s \pi$ and clearly $\Pi P \Pi \in
\clT_q$.

The kernel of $\tau_q$ is $\clT_{q+1}$ because $\tau_q (P) =0$ implies that $\si_{2q} (P) =0$ so $P \in \lag^{+}_{2q+1} (A)$ so $\si_{2q+1} (P)$ is odd by Theorem~\ref{theo:parity}. But $\Pi P \Pi = P$ implies that $\si_{2q+1} (P) = \pi \si_{2 q+1} (P) \pi$. This implies that $\si_{2q+1} (P)$ is even because $F$ has a definite parity. Indeed, if for instance $F \subset \D^+(TM) \otimes A$, and $f = \pi g \pi$ with $g \in \symb (T_xM) \otimes \op{End} A_x$, then
in the decomposition
\[
\D \left(T_xM\right) \otimes A_x = \left(\D^+ (T_xM) \otimes A_x\right) \oplus \left(\D^{-} (T_xM) \otimes
A_x\right),
\]
$f$ has the form $
(\begin{smallmatrix} f_{++} & 0 \\ 0 & 0
\end{smallmatrix})$, so $f$ is even. Consequently, $\si_{2q+1} (P) = 0$ so $P \in \lag_{2q+2} (A)$. The formulas giving the symbol of products, adjoints, the operator norm and the Schwartz kernel on the diagonal follow directly from Theorem~\ref{theo:lag}.  Observe that the $\bigo (k^{-\frac{1}{2}})$ in Assertion 5 of Theorem~\ref{theo:lag} becomes a $\bigo (k^{-1})$ in~\eqref{eq:trace_diag} because $P$ being even, the restriction of the asymptotic expansion of its Schwartz kernel~\eqref{eq:exp_Lag_bis} to the diagonal only involves integral powers of $k^{-1}$. 
\end{proof}

\begin{rema} \label{rem:odd_toeplitz}
We can consider as well {\em odd Toeplitz operators}, that is $T \in \lag^{-} (A)$ such that $T =
\Pi T \Pi$. The space of these operators is $k^{-\frac{1}{2}} \clT$. Indeed, if $T $ is such an operator, then its symbol $ \si_0 (T) = \pi \si_0 (T) \pi$ is at the same time even and odd because $F$ has a definite parity, so $\si_0(T) =0$ so $T
\in \lag_1 (A)$, so $k^{\frac{1}{2}} T \in \lag (A)$ and is even, so $k^{\frac{1}{2}} T \in \clT$. \qed
\end{rema}


For any $f \in \Ci(M)$ and $k \in \N$, define the endomorphism $T_k(f)$ of $\Hilb_k$ such that
\[
\left\langle T_k (f) \psi, \psi' \right\rangle = \left\langle f \psi, \psi'\right\rangle,
\qquad \forall \psi, \psi' \in \Hilb_k.
\]
Viewed as an operator of $\Ci (M, L^k \otimes A)$, $T_k(f)$ is merely $\Pi_k f \Pi_k$. It follows from part~\eqref{theo2.3.2} of Theorem~\ref{theo:lag} that the family $(T_k(f))$ belongs to $\clT$ and has symbol $\tau_0 (T_k (f))= f
\op{id}_{F}$. By part~\eqref{theo3.2.3} of Theorem~\ref{theo:toeplitz}, we deduce that
\[
T_k (f) T_k (g) = T_k (fg) + \bigo \left(k^{-1}\right).
\]
A consequence of Theorem~\ref{theo:unitary_equivalence} will be that
\[
\left[T_k (f), T_k (g)\right] = i k^{-1} T_k (\{ f,g \}) + \bigo \left(k^{-2}\right)
\]
with $\{ f, g \}$ the Poisson bracket of $f$ and $g$ with respect to $\om$. This equality does not follow from Theorem~\ref{theo:toeplitz}. However,  the center of $\op{End} F_x$ consisting on the scalar multiple of the identity,  the following characterization of the Toeplitz operators having a scalar symbol follows from Theorem~\ref{theo:toeplitz}: for any $P \in
\clT $,
\[
P = T (f) +\bigo (k^{-1}) \text{ for some } f \in \Ci (M) \, \Leftrightarrow \;
\forall Q \in \clT, \, [P,Q] \in \clT_{1}.
\]
By Jacobi identity, this proves that $[T_k (f), T_k (g)] = k^{-1} T_k (h) +
\bigo (k^{-2})$ for some function $h \in \Ci(M)$.

\subsection{A unitary equivalence}\label{sec:unitary-equivalence}

Consider an auxiliary vector bundle $B$ with an arbitrary rank. Set $F' = \D_0 (T M) \otimes B $ with $\D_0 (T_xM) \subset \D (T_xM)$ the subspace of constant polynomials. Then by Theorem~\ref{sec:constr-proj}, there exists a projector $\Pi'$ in $\lag^{+}(B)$ having symbol
\[
\si_0(\Pi') = \Uu_{00} \otimes \op{id}_{B} \in \Ci (M,\symb(M) \otimes \op{End}
B),
\]
where $\Uu_{00} (x) \in \symb (T_xM)$ is the orthogonal projector of
$\D (T_xM)$ onto $\D_0 (T_xM)$, the notation being the same as in~\eqref{eq:def_Ualphabeta}.

Starting from $\Pi'$, we define $\Hilb_k' := \op{Im} \Pi'_k$ and the corresponding Toeplitz space $\clT' := \{ P
\in \lag^{+} (B)/ \; \Pi' P \Pi' = P \}$. Since $\D_0(TM) = \C$, $F' \simeq B$, so the symbols of the Toeplitz operators of $\clT'$ are sections of $\op{End} B$:
\[
0 \rightarrow \clT'_{q+1} \rightarrow \clT'_q
\xrightarrow{\tau'_q} \Ci (M, \op{End} B) \rightarrow 0.
\]
Our goal now is to establish an equivalence between the $(\Hilb_k, \clT)$ and $(\Hilb_k', \clT')$ when the bundle $B$ is $F$. The critical point is the existence of a convenient symbol. Recall our assumption that $F \subset
\D ^{\ep} (TM) \otimes A$ with $\ep \in \{ \pm 1 \}$.

\begin{lemm}\label{lem:unitary-symbole}
If $B = F$, then there is a canonical symbol $\rho \in \Ci (M, \symb^{\ep} (M) \otimes \op{Hom} (A, B))$ such that $ \rho^* \rho = \pi $ and $\rho \rho ^* = \Uu_{00} \otimes \op{id}_{B}.$
\end{lemm}

\begin{proof}
On the one hand, $\pi (x) $ is the
orthogonal projector of $\D (T_xM) \otimes A_x$ onto $F_x$. On the other hand, $\pi'(x) := \Uu_{00} (x) \otimes \op{id}_{B_x}$ is the orthogonal projector of $\D (T_xM) \otimes B_x$ onto $\C \otimes B_x$. Since $B = F$, the images of $\pi (x)$ and $\pi'(x)$ are isomorphic by the map $ \xi (x) : F_x \rightarrow \op{Im} \pi '(x)$ sending $f$ into $1
\otimes f$. We define $\rho (x)$ as the extension of $\xi(x)$
\begin{gather}\label{eq:def_rho}
\rho (x) : F_x \oplus F_x^{\perp} \rightarrow \left(\op{Im} \pi'(x)\right) \oplus \left(\op{Im} \pi'(x)\right)^{\perp}
\end{gather}
having the block decomposition $
(\begin{smallmatrix} \xi (x) & 0 \\ 0 & 0
\end{smallmatrix})$. So $\rho (x)$ is canonically defined. The equalities $ \rho(x)^* \rho(x) = \pi (x) $ and $\rho(x) \rho(x) ^* = \pi '(x)$ are easily verified by using that $\xi (x)$ is unitary. Writing $\rho $ in terms of a local frame of $F$, we see that $\rho (x)$ depends smoothly on $x$. Finally, $F_x \subset \D^{\ep} (T_xM)\otimes A_x$ and $\op{Im} \pi'(x) \subset \D^{+} (T_x M) \otimes B_x$, so $\rho (x) \in \symb^{\ep} (M)_x \otimes \op{Hom} (A_x, B_x)$.
\end{proof}

\begin{theo}\label{theo:unitary_equivalence}
Assume that $B = F$ and $\rho$ is the symbol defined above. Then there exists $U \in \lag ^{\ep}(A,B)$ with symbol $\si_0(U) = \rho$ and such that
\begin{gather}\label{eq:U}
U^*_k U_k = \Pi_k, \qquad U_k U_k^* = \Pi'_k
\end{gather}
when $k$ is sufficiently large. Modifying $\Pi_k'$ for a finite number of $k$, we can choose $U$ so that~\eqref{eq:U} holds for any $k$. In this case, the Toeplitz algebras $ \clT$ and $\clT'$ are isomorphic by the map sending $P$ into $U P U^* $. Furthermore, $P \in \clT_q$ if and only if $UPU^* \in \clT'_q$ and when this is satisfied
\begin{gather}\label{eq:symbol_iso}
\tau'_q (U P U^*) = \tau_q (P).
\end{gather}
\end{theo}

\begin{proof}
Choose $W \in \lag^{\ep} (A,B)$ with symbol $\rho$ and set $V := \Pi' W \Pi$. Then $V
\in \lag^{\ep} (A,B)$ with $ \si_0(V) = \pi' \rho \pi = \rho$ and since $\rho^* \rho =
\pi$ and $\rho \rho^* = \pi'$, we have
\begin{gather}\label{eq:first_Step}
V^*_k V_k = \Pi_k + \bigo \left(k^{-\frac{1}{2}}\right), \qquad V_k V^*_k = \Pi_k' + \bigo \left(k^{-\frac{1}{2}}\right)
\end{gather}
So $V_k$, viewed as an operator from $\Hilb_k$ to $\Hilb_k'$, is invertible when $k$ is sufficiently large.

Observe also that $V^*V$ is a Toeplitz operator of $\clT$ with symbol $\op{id}_{F}$. So $V^*V = \Pi + Q$ with $Q \in
\clT_1$. Since $\| Q_k\| = \bigo (k^{-1})$, the spectrum of $Q_k$ is contained in $[-\frac{1}{2}, \frac{1}{2}]$ when $k$ is sufficiently large. Modifying $Q_k$ for a finite number of $k$, we can assume this holds for any $k$, and when $k$ is sufficiently large, we still have $V_k^*V_k = \Pi_k + Q_k$. Let $P_k$ be the endomorphism of $\Ci (M, L^k \otimes A$) which is zero on $\Hilb_k^{\perp}$ and equal to $(\op{Id}_{\Hilb_k} + Q_k) ^{-1/2}$ on $\Hilb_k$. We claim that $(P_k)$ belongs to $\clT$ and has symbol $\op{id}_{F}$.

Assuming this temporarily, it follows that $U_k := V_k P_k $ belongs to $\lag ^{\ep} (A,B)$, has symbol $\rho$ and satisfies when $k$ is sufficiently large $U_k^* U_k = P_k V_k^* V_k P_k = \Pi_k$. Since $U_k$, viewed as an operator from $\Hilb_k$ to $\Hilb'_k$ is invertible, this also implies that $U_kU_k^* = \Pi'_k$.

To prove the claim above, we write the Taylor expansion $(1+x)^{-1/2} = 1 +
\sum_{\ell = 0 }^m a_\ell x^\ell + x^{m+1} f_m (x)$ with $f_m$ a continuous function $[-\frac{1}{2},\frac{1}{2}] \rightarrow \R$. Then
\begin{gather}\label{eq:racine}
P_k = \Pi_k + \sum_{\ell = 0 }^m a_\ell Q_k^{\ell} + Q_k^{m+1} f_m (Q_k).
\end{gather}
Then we show that $P$ belongs to $\lag ^+ (A)$ by arguing as in the proof of Theorem~\ref{theo:constr-proj}: $Q^\ell \in \clT_\ell $ and $\| f_m (Q_k) \| = \bigo (1)$, so the Schwartz kernel family of $Q_k^{m+1} f_m (Q_k) = Q_k^m f_m (Q_k) Q_m$ is in $\bigo (k^{2n - m-1 })$. Choosing $m$ sufficiently large at each step, we then deduce from~\eqref{eq:racine} that the Schwartz kernel family of $P_k $ is $\bigo (k^{-\infty})$ outside the diagonal and that the local expansions~\eqref{eq:expansion_kernel} hold.

So we have proved the existence of $U \in \lag^{\ep}(A,B)$ with $\si_0(U) = \rho$ and satisfying~\eqref{eq:U} for any $k$ except a finite set. For the missing $k$'s, we modify $\Pi_k'$ by choosing any subspace $\Hilb'_k$ of $\Ci (M, L^k \otimes B)$ having the same dimension as $\Hilb_k$, define $\Pi_k'$ as the orthogonal projector onto $\Hilb_k'$ and $U_k$ as any isometry $\Hilb_k \rightarrow \Hilb'_k$ extended to zero on $\Hilb_k^{\perp}$. Then $\Pi_k'$ and $U_k$ have a smooth Schwartz kernel, so the new families $\Pi '$ and $U$ are still in $\lag ^{+} (B)$ and $\lag ^{\ep} (A,B)$ respectively.

It is now easy to prove the last assertion: if $P \in \lag ^{+} (A)$, then $U P U^* \in \lag ^{+}(B)$ because $U \in \lag ^{\ep} (A,B)$ and $U^* \in
\lag^{\ep} (B,A)$. If $ \Pi P \Pi = P$, then $\Pi' (U P U^*) \Pi' = U P U^*$ by~\eqref{eq:U}. So $P \in \clT$ implies that $U P U^* \in
\clT'$, which defines an isomorphism from $\clT$ into $\clT'$ because we can invert it by sending $Q$ into $U^* Q U$. Furthermore, $\si_0 (UPU^*) = \rho^* \si_0 (P)
\rho$ which leads to~\eqref{eq:symbol_iso}.
\end{proof}

A first corollary is the computation of the symbols of commutators in terms of Poisson bracket. Recall the Toeplitz operators $T_k (f) : \Hilb_k
\rightarrow \Hilb_k$ associated to $f \in \Ci (M)$. Define similarly $T'_k (f) : \Hilb'_k \rightarrow \Hilb_k'$.

\begin{coro}
$ [T_k (f), T_k (g)] = i k^{-1} T_k (\{ f,g \}) + \bigo (k^{-2})$ for any $f, g \in \Ci (M)$.
\end{coro}
Another proof will be provided in Proposition~\ref{prop:subsymbolic}.

\begin{proof}
This amounts to show that for any two Toeplitz operators $T$, $S$ of
$\clT$ with symbol $\tau_0(T) = f \op{id}_F$, $\tau_0(S) = g
\op{id}_F$, we have $\tau_1 ([T,S]) = i \{ f, g\} \op{id}_F$. By Theorem~\ref{theo:unitary_equivalence}, this holds for $\clT$ if and only if this holds for $\clT'$.  The results for $\clT'$ has been proved in~\cite[Theorem~1.4]{oim}, when the projector is chosen as in~\cite[Theorem~1.1]{oim}.
\end{proof}

The operators $T'_k(f)$ are defined not only for $f \in \Ci (M)$ but also for $f \in \Ci (M, \op{End} B)$. Since $\tau_q'(k^{-q} T_k' (f)) = f$, it follows that we can define the Toeplitz operators of $\clT'$ as the families $(T_k)$ such that for any $N$,

\[
T_k = \sum_{\ell =0}^{N} k^{-\ell} T_{k} (f_\ell) + \bigo \left(k^{-(N+1)}\right)
\]
for a sequence $(f_\ell)$ of $\Ci (M, \op{End} B)$. This provides a definition of $\clT'$ without any reference to the algebra $\lag ^{+} (B)$. Observe also that the coefficients $f_{\ell}$ are uniquely determined by $T$ and the map $\clT' \rightarrow \Ci (M, \op{End} B) [[\hbar]]$ sending $T$ into $\sum \hbar^\ell f_{\ell}$ is a full symbol map, meaning that it is onto and its kernel is $\clT' \cap \bigo (k^{-\infty})$. This full symbol map can also be used to get uniform control of the product of Toeplitz operators, cf.~\cite{oim}. But unfortunately, this does not hold for $\clT$, except in the particular case where $F$ has rank one, so that $\op{End} F \simeq \C$. This happens in particular for higher Landau level in dimension $n=1$.

A second consequence of Theorem~\ref{theo:unitary_equivalence} is the computation of the dimension of our quantum spaces.  Here the parity assumption is not necessary.


% Official author native source lines 1103--1110
\begin{theo}\label{theo:dim_general}
Let $\Pi \in \lag (A)$ be a projector whose symbol $\pi = \si_0 (\Pi)$ has a constant rank. Then the dimension of $\Hilb_k= \op{Im}(\Pi_k)$ is
\[
\op{dim} \Hilb_k = \int_M \op{ch} \left(L^k \otimes F\right) \; \op{Td}
(M)
\]
when $k$ is sufficiently large, where $F$ is the subbundle of $\D (TM) \otimes A$ given by $F_x = \op{Im} \pi (x)$ for any $x \in M$.
\end{theo}

% Official author native source lines 1111--2324



\begin{proof}
We introduce a new family $(\Hilb''_k:= \op{Ker} D_k)$ where $D_k$ is the spin-c Dirac operator acting on $\Ci (M, L^k \otimes B \otimes S)$ with  $S :=\bigwedge (T^*M)^{0,1} $  the spinor bundle. By the Atiyah--Singer theorem and a vanishing theorem~\cite{BoUr96,MaMa02}, the dimension of $\Hilb''_k$ is given by the Riemann--Roch number of $L^k \otimes B$ when $k$ is sufficiently large. We claim that the projector $\Pi''_k$ of $\Ci (M, L^k \otimes B \otimes S)$ onto $\Hilb_k''$ belongs to $\lag ^+ (B \otimes S)$ and has symbol $\Uu_{00}
\otimes p_B$ where $p_B $ is the section of $\op{End} (B \otimes S)$ equal at each $x\in M$ to the projector of $B_x \otimes S_x$ onto $B_x \otimes
\C$. This is actually a reformulation of results by Ma and Marinescu~\cite{MaMa}, as is explained in ~\cite[Appendix A]{oim}.  Alternatively this follows from the companion paper~\cite{oim_copain}.

Now the image of the symbol $\rho_{00} \otimes p_B$ is isomorphic with $B$, so by Theorem~\ref{theo:unitary_equivalence}, when $k$ is sufficiently large, $\Hilb''_k$ has the same dimension as $\Hilb_k' = \op{Im} \Pi'_k$, where $\Pi_k'$ is any self-adjoint projector of $\lag^+(B)$ with symbol $ \rho_{00} \otimes \op{id}_B$.

To conclude, when $\Pi$ is even, by another application of Theorem~\ref{theo:unitary_equivalence}, for $B=F$, $\Hilb'_k$ and $\Hilb_k$ have the same dimension when $k$ is sufficiently large.  The same proof works for $\Pi$ not being necessarily even. Actually, the existence of $V$ satisfying~\eqref{eq:first_Step} already implies that the dimensions of $\Hilb_k$ and $\Hilb_k'$ are the same when $k$ is large. 
\end{proof}



\section{Landau Hamiltonian algebra} \label{sec:landau-levels-cn}

In this section, we come back to the algebra $\symb (\C^n)$ introduced in Section~\ref{sec:symbol-spaces}. We extend the action of the elements of $\symb (\C^n)$ on $\D (\C^n)$ to the complete polynomial space and we compute the corresponding Schwartz kernel. This will be used in the sequel to give an intrinsic definition of the symbol maps $\si_q$, cf. Definition~\ref{def:symbol}, and to understand the composition properties of the class $\lag (A,B)$.


Let $\pol (\C^n)$ be the space of polynomials map from $\C^n$ to $\C$, so any $f
\in \pol (\C^n)$ has the form $ f = \sum a_{\al \be} z^ \al \con{z}^\be$ where the sum is finite and the $a_{\al \be}$ are complex numbers. The space $\D (\C^n)$ introduced in Section~\ref{sec:symbol-spaces} is the subspace of $\pol (\C^n)$ of antiholomorphic maps. We endow $\pol (
\C^n)$ with the same scalar product
\begin{gather}\label{eq:scal_prod_2}
\langle f, g \rangle = (2\pi)^{-n} \int_{\C^n} e^{-|z|^2} f (z) \con{g (z)} \; d \mu_n (z)
\end{gather}
as in~\eqref{eq:scal_product_bargm} for $\D (
\C^n)$. The family $(((\al + \be)!)^{-\frac{1}{2}} z^{\al} \con{z}^\be, \;
\al, \be \in \N^n)$ is an orthonormal basis of $\pol (\C^n)$.

For any $i =1, \ldots, n $, introduce the endomorphism $a_i = \partial_{\con{z}_i}$ and its adjoint $a_i^{*} = \con{z}_i - \partial_{z_i}$. They satisfy the bosonic commutation relations
\[
\left[a_i,a_j\right] = \left[a_i^{*}, a_j^{*}\right] = 0, \qquad \left[a_i, a_j^{*}\right] =
\delta_{ij}.
\]
So the $a_i^{*} a_i$'s are mutually commuting Hermitian endomorphisms. Their eigenspaces are the Landau levels of $\C^n$. In the sequel, we use the notation $a^\al := a_1^{\al(1)} \ldots a_n^{\al (n)}$ and $(a^{*})^{\al} := (a_1^*)^{\al(1)} \ldots (a_n^*)^{\al (n)}$.

\begin{prop}\label{prop:linearlandau}\leavevmode
\begin{enumerate}
\item\label{prop4.1.1} For $i=1,\, \ldots,\, n$, $a_i^* a_i$ is diagonalisable with spectrum $\N$. So we have a decomposition into mutually orthogonal joint eigenspaces $ \pol (\C^n) = \bigoplus_{\al\,\in\,\N^n} \lag_\al$ with $\lag_{\al} = \bigcap_{i=1}^{n} \op{ker} (a_i^* a_i - \al (i)).$
\item\label{prop4.1.2} $\lag_0 = \C [z_1,\,\ldots,\,z_n]$ and for any $\al \in \N^n$, $
\lag_{\al} = (a^*)^\al \lag_0 $.
\item\label{prop4.1.3} For any $\al, \be \in \N^n$, let $\Vvv_{\al \be} := (\al ! \be !)^{-\frac{1}{2}} (a^{*})^{\al} \Vvv_{00} a ^{\be}$ with $\Vvv_{00} $ the orthogonal projector of $\pol (\C^n)$ onto $\lag_0$. Then
\begin{enumerate}
\item\label{prop4.1.3.a} $\Vvv_{\al\be}$ is zero on the $\lag_{\ga}$'s with $\ga \neq
\be$ and restricts to a unitary isomorphism from $\lag_{\be}$ to $\lag_{\al}$
\item\label{prop4.1.3.b} $\Vvv_{\al \al}$ is the orthogonal projector onto $\lag_\al$.
\item\label{prop4.1.3.c} $ \Vvv_{\al \be} \circ \Vvv_{\tilde{\al} \tilde{\be}} =
\delta_{\be \tilde{\al}} \Vvv_{\al \tilde \be} $ and $\Vvv_{\al \be} ^{*} = \Vvv_{\be \al} $
\end{enumerate}
\end{enumerate}
\end{prop}



\begin{proof}
The result is certainly standard in condensed matter theory. For the convenience of the reader, we explain briefly the proof
for $n=1$. The extension in higher dimension is straightforward. We write $a:= a_1$, recall the commutation relation $ [a, a^{*}]= 1 $ and set $\lag_m := (a^{*})^m (\C[z])$ for any $m \in \N$.

We check by induction that $\lag_m = \ker (a^* a - m)$. First, writing $\langle a^{*} a f, f \rangle = \| a f \|^2 $, it comes that $\ker a^* a = \ker a =
\lag_0$. Assume now that $\lag_m = \ker (a^* a - m)$. By the commutation relation, $f \in \lag_m$ implies that $a^{*} a a ^{*} f = (m+1) a^* f$, so $\lan_{m+1} \subset \ker (a^* a - (m+1))$. Conversely, by the commutation relation again, $a^* a f = (m+1) f $ implies that $ (a^* a) a f = m a f$ so $ af \in \lan_m$ and $f = (m+1)^{-1} a^* (af) \in
\lan_{m+1}$.

To conclude that $a^* a $ is diagonalizable with eigenvalues in $\N$, it suffices to prove that $\pol (\C)$ is spanned by the $\lan_m$. Introduce the filtration $\Fam_m := \oplus _{\ell=0}^{m} \con{z}^\ell
\C[z]$, $m \in \N$. If $f \in \C [z]$, then $\con{z}^m f = (a^{*})^m f \mod \Fam_{m-1}$. So $\Fam_m = \lan_m +
\Fam_{m-1} = \ldots = \lan_m + \lan_{m-1} + \ldots + \lan_0$ by reiterating.

So we have proved that $\pol (\C) = \bigoplus \lan_m$ with $\lan_m = \ker (a^* a - m)$, which shows the first and second assertions of the proposition. By the commutation relation, $a a^{*} = (m+1)$ on $\lan_m$. So $a^{*} : \lan_m \rightarrow \lan_{m+1}$ is invertible with inverse $(m+1)^{-1} a: \lan_{m+1} \rightarrow \lan_{m}$. We conclude that
\begin{itemize}
\item if $p \leqslant m$, then $a^p$ restricts to an isomorphism from $\lan_m$ to $\lan_{m-p}$, whose inverse is the restriction of $\frac{(m-p)!}{m!} (a^{*})^m$ to $\lan_{m-p}$.
\item if $p > m$, then $a ^p (\lan_m) = \{0\}$.
\end{itemize}
With these two facts, we easily check the third assertion.
\end{proof}

By the last assertion of Proposition~\ref{prop:linearlandau}, the space $\widetilde{\symb} (\C^n)$ of endomorphisms of $\pol (
\C^n)$ generated by the $\Vvv_{\al \be}$'s is closed under composition, so it is an algebra. By the following proposition, $\widetilde{\symb} (\C^n)$ is isomorphic with the algebra $\symb (\C^n)$ introduced in Section~\ref{sec:symbol-spaces}, through the map sending $\Vvv_{\al\be}$ into $\Uu_{\al\be}$.

\begin{prop}
The elements of $\widetilde{\symb} (\C^n)$ preserve the
subspace $\D (\C^n)$ of $\pol (\C^{n})$. Furthermore, the restriction map $\op{res} : \widetilde{\symb} (\C^n) \rightarrow \op{End} (\D (\C^{n}))$ is injective, with image $\symb (\C^n)$ and $\op{res} (\Vvv_{\al \be}) = \Uu_{\al
\be}$.
\end{prop}

Recall the decomposition~\eqref{eq:symb+_-} of $\symb (\C^n)$ into the subspaces of even and odd elements. Since $\widetilde \symb (\C^n) \simeq \symb (\C^n)$, this gives us a new decomposition
\[
\widetilde{\symb}
(\C^n) = \widetilde{\symb}^+ (\C^n) \oplus \widetilde{\symb}^-
(\C^n).
\]

\begin{proof}
Observe first that the operators $a_i$, $a_i^{*}$ and the projector $\Vvv_{00}$
preserves $\D (\C^n)$. Furthermore, for any $f \in \D (\C^n)$, $a_i f =
\partial_{\con{z}_i} f$, $a_i^{*} f = \con{z}_i f$ and $\Vvv_{00} f = f (0)$. Consequently, the operators $\Vvv_{\al \be}$ preserve $\D (\C^n)$ and an easy computation shows that
\[
\Vvv_{\al\be} \left((\be !)^{-\frac{1}{2}} \con{z}^\be \right) = (\al !)^{-\frac{1}{2}}
\con{z}^{\al}, \qquad \Vvv_{\al \be} (\con{z}^{\ga}) = 0, \quad \forall \ga \in \N^n
\setminus \{ \be \}.
\]
This means that the restriction of $\Vvv_{\al \be}$ to $\D (\C^n)$ is exactly the endomorphism $\Uu_{\al \be}$ introduced in Section~\ref{sec:symbol-spaces}, cf. Equation~\eqref{eq:def_Ualphabeta}. So the restriction map $\op{res}$ is well-defined, its image is $\symb (\C^n)$, and the $\Uu_{\al\be}$'s being linearly independent, it is injective.
\end{proof}

Let us compute the Schwartz kernel of each $\Vvv_{\al\be}$.

\begin{lemm}\label{lem_kenrel_u_alpha_beta}
For any $f \in \pol (\C^n)$, we have
\[
(\Vvv_{\al\be} f)(u) = (2 \pi)^{-n} \int_{\C^{n} } e^{ u\cdot \con{v} - |v|^2 }
\Ppp_{\al\be} (u-v) \, f(v) \; d\mu_n (v)
\]
where $u \cdot \con v = \sum u_i \con{v}_i$, $\Ppp_{\al,\be} (z) = \bigl(\al! \be ! \bigr) ^{-\frac{1}{2}} \bigl(
\partial_z - \con{z})^{\al} z^\be $.
\end{lemm}

In particular, the orthogonal projector $\sum_{|\al| = m } \tilde{\rho} _{\al
\al}$ onto $\bigoplus_{|\al| = m} \lag_{\al}$ has the Schwartz kernel
\begin{gather}\label{eq:noyau_niveau_m}
(2 \pi)^{-n} e^{u \cdot \con{v} - | v |^2 } Q_m ^{(n-1)} \left(| u - v |^2\right) \; d \mu_n (v)
\end{gather}
where $Q_m^{(n-1)}$ is the Laguerre polynomial and we have used~\eqref{eq:a_0laguerre}.
\begin{proof}
For $\al = \be =0$, this is the well-known formula for the Schwartz kernel $K (u,v) = (2\pi)^{-n} e^{u \cdot \con v - |v|^2} $ of the projector onto the Bargmann space, which in our setting is the $L^2$-completion of $\lag_0$. So the Schwartz kernel of $(a^{*})^{\al} \Vvv_{00} a^{\be}$ is $ K_{\al,\be} (u,v) = (\con{u} -\partial_u)^{\al} (-\partial_{\con{v}})^{\be} K(u,v)$. To compute this, we use that for a polynomial $g (z,\con{z})$
\begin{align*}
(\con{u}_i - \partial_{u_i}) (K(u,v) g(u-v)) & = K(u,v) (a_i^{*} g)(u-v), \\
(\partial_{\con{v}_i}) (K(u,v) g(u-v)) & = K(u,v) (b_i^{*} g)(u-v)
\end{align*}
where $b_i^{*} := z_i - \partial_{\con{z}_i}$. So $K_{\al,\be} (u,v) = K(u,v) p(u-v) $ with $p (z,\con z) = (a^{*})^{\al} (b^{*})^\beta 1= \bigl(
\partial_z - \con{z})^{\al} z^\be $, which ends the proof.
\end{proof}

So on one hand, the elements of $\widetilde{\symb} (\C^n)$ act on $\pol (\C^n)$, on the other hand, the Schwartz kernel of $\Vvv_{\al\be}$ is given by a polynomial $\Ppp_{\al\be} \in \pol (\C^n)$.

\begin{prop}\label{prop:tilde_symb_and_op}
$\widetilde{\symb} (\C^n)$ consists of the endomorphisms $V$ having the form
\begin{gather}\label{eq:sch_ker}
(V f)(u) = (2\pi)^{-n} \int_{\C^n} e^{u\cdot \con v - |v|^2} q (u-v) f(v) \; d\mu_n (v)
\end{gather}
with $q \in \pol (\C^n)$. Furthermore, the map $\op{Op} : \pol (\C^n) \rightarrow
\widetilde{\symb} (\C^n)$, sending $q$ into $V$, is an isomorphism which preserves the parity  and
\begin{equation} \label{eq:rel_bizare}
\begin{split}
\op{tr} \left(\op{Op} (q)\middle|_{\D (\C^n)} \right) &= q (0) \\
 \op{Op} (q) \circ \op{Op} (f) &= \op{Op} (\op{Op} (q) f) \\
 \langle q, f \rangle &= \left(\op{Op}(f)^* q\right)(0)
\end{split}
\end{equation}
for any $q, f \in \pol (\C^n)$. 
\end{prop}
\begin{proof}
Lemma~\ref{lem_kenrel_u_alpha_beta} says that $\op{Op} (\Ppp_{\al\be}) =
\Vvv_{\al \be}$. The family $(\Ppp_{\al\be})$ is a basis of $\pol (\C^n)$ because $\Ppp_{\al\be} (z) = (\al ! \be !) ^{-\frac{1}{2}}(- z)^{\al} \con{z}^{\beta} + $ a linear combination of $z^{\al'} \con{z}^{\be'}$ with $\al' < \al$ and $\be' < \be$. Since $(\Uu_{\al\be})$ is a basis of $\symb (\C^n)$, $(\Vvv_{\al \be})$ is a basis of $\widetilde { \symb } (\C^n)$ and it follows that $\op{Op}$ is an isomorphism. This isomorphism preserves the parity, because $\Vvv_{\al
\be} $ and $\Ppp_{\al\be}$ have both the same parity as $|\al| + |\be|$. For the first equation of~\eqref{eq:rel_bizare}, it suffices to prove it for $q = \Ppp_{\al,\be}$, and in this case, it follows from $\op{tr}
\Uu_{\al\be} =\delta_{\al\be} = \Ppp_{\al\be} (0)$.

To prove the second equation of~\eqref{eq:rel_bizare}, observe that we recover $q$ from the Schwartz kernel of $\op{Op} (q)$  by multiplying by $(2\pi)^{n}$ and setting $v = 0 $, that is $ f(u) = (2 \pi)^n \op{Op} (f)\linebreak(u,0)$. Let $g \in \pol (\C^n)$ be the function such that $\op{Op} (g) = \op{Op} (q) \circ \op{Op} (f)$. Then using the previous observation for $f$ and for $g$, we have
\begin{align*}
(\op{Op} (q) f) (u) &=  (2\pi)^n \int_{\C^n} \op{Op}(q) (u,v) \op{Op} (f) (v,0) \; d\mu_n (v) \\
&=  (2 \pi)^n \op{Op}(g) (u,0) \\
&=  g (u).
\end{align*}
The proof of the third equation of~\eqref{eq:rel_bizare} is similar by using that $ \op{Op}(f)^* (0,u) = \con{\op{Op} (f) (u,0)}= (2 \pi)^{-n} \con{f (u)}$.
\end{proof}



As a last remark, we can replace in the previous definitions $\C^n$ with any $n$-dimensional Hermitian space $\EE$ as we did in section~\ref{sec:symbol-spaces}. So we denote by $\pol (\EE)$ the space of polynomial maps $\EE \rightarrow \C$ and by $\widetilde{\symb} (\EE)$ the space of endomorphisms of $\pol (\EE)$ having the form~\eqref{eq:sch_ker}, where we interpret $u \cdot \con{v}$ as the scalar product of the vectors $u$, $v$ of $\EE$ and $|v|$ as the norm of $v$. Observe as well that the map $\op{Op} : \pol (\EE) \rightarrow
\widetilde{\symb} (\EE)$ is well-defined. Furthermore the restriction from $\pol (\EE)$ to its subspace $\D (\EE)$ induces an isomorphism $\widetilde{\symb} (\EE) \simeq \symb (\EE)$.


\section{The Schwartz kernels of operators of \texorpdfstring{$\lag (A, B)$}{L(A,B)} } \label{sec:schw-kern-oper}



\subsection{The section \texorpdfstring{$E$}{E}} \label{sec:section-ee}

An important ingredient in the global Schwartz kernel description of operators of $\lag (A,B)$ is a section~$E$ of $L \boxtimes \con{L}$ satisfying the following conditions. For any $y\in M$, denote by $E_y $ the section of $L \otimes \con{L}_y$ given by $E_y(x) = E(x,y)$ for any $x \in M$. Then we will assume that for any $y\in M$
\begin{equation}\label{eq:hypotheseE}
\begin{split}
 E_y(y) &= u \otimes \con{u},   \qquad\forall u \in L_y \text{ with } |u| =1, \\
 (\nabla E_y)(y) &=0,  \\
 (\nabla_{\xi} \nabla_{\eta} E_y)(y) &= - \left(\tfrac{i}{2} \om (\xi, \eta) + \tfrac{1}{2}
\om (\xi, j \eta) \right) E_y (y),  \qquad\forall \; \xi, \eta \in T_yM
\end{split}
\end{equation}
Such a section appeared already in the expansion~\eqref{eq:expansion_kernel} as follows. Choose a unitary frame $t$ of $L$ and a coordinate system on the same open set, then the section
\begin{gather}\label{eq:exempleE}
E (y+ \xi, y) := e^{-\varphi (y, \xi)} t(y+ \xi) \otimes \con{t}(y),
\end{gather}
with $\varphi$ defined as in~\eqref{eq:expansion_kernel}, satisfies~\eqref{eq:hypotheseE}. From this local construction, we easily obtain a global section~$E$ by using a partition of unity.

The conditions~\eqref{eq:hypotheseE} determine the second-order Taylor expansion of $E$ at $(y,y)$ in the directions tangent to the first factor of $M^2$. Since any tangent vector of $M^2$ at $(y,y)$ is the sum of a vector tangent to the diagonal and a vector tangent to the first factor, we deduce that $E$ is uniquely determined modulo a section vanishing to third order along the diagonal.

The function $\psi_y (x)= - 2 \ln |E_y (x)|$ vanishes to second order at $y$ and for any $\xi, \eta \in T_yM$, $(\xi.\eta. \psi_y)(y) = \om (\xi, j \eta)$, so $\psi_y(x)>0$ when $x\neq y$ is sufficiently close to $y$. So modifying $E$ outside the diagonal, we can assume that it satisfies as well
\begin{gather}\label{eq:normE}
|E (x,y) | <1, \qquad \forall (x,y) \in M^2 \text{ such that } x \neq y
\end{gather}
Another important property of $E$ is the symmetry:
\begin{gather}\label{eq:symE}
\con{E (x,y)} = E(y,x) +
\bigo \left(|x-y|^3\right)
\end{gather}
For a longer discussion, the reader is referred to~\cite{oim}.

In the sequel we will need the following expression of $E$ in terms of complex coordinates and a frame of $L$, both normal at a point $p_0 \in M$.  We say that a function or a section on $M$ (resp. $M^2$) is in $\bigo_{p_0} (m)$ (resp. $\bigo_{p_0, p_0} (m)$) if it vanishes to order $m$ at $p_0$ (resp. $(p_0,p_0)$).  Let $(\partial_i)_{i =1}^{n}$ be an orthonormal basis of $T^{1,0}_{p_0} M$, i.e. $\frac{1}{i}\om_{p_0} (\partial_i, \con{\partial}_j) = \delta_{ij}$. Choose complex valued functions $z_i$ on a neighborhood of $p_0$ such that
\begin{gather}\label{eq:coordonnees_z}
z_i (p_0) =0, \qquad dz_i (\partial_j) =\delta_{ij}, \quad dz_i (\con \partial_j) = 0 \quad
\text{ at }p_0.
\end{gather}
Then $(\op{Re} z_i, \op{Im}z_i)_{i=1}^n$ is a coordinate system on a neighborhood of $p_0$ and $\om|_{p_0} = i \sum dz_i \wedge d\con{z}_i$. The curvature of $L$ being $\frac{1}{i} \om$, there exists a unitary frame $t$ of $L$ at $p_0$ such that
\begin{gather}\label{eq:normal_frame}
\nabla t = \tfrac{1}{2} \textstyle{\sum} (z_i d \con{z}_i - \con{z}_i d z_i) \otimes t + \bigo_{p_0} (2).
\end{gather}
To construct such a $t$, multiply any unitary local section by $\exp (i (f_1 + f_2))$ where $f_1$ and $f_2$ are real valued functions respectively linear and quadratic in $z_i$, $\con{z}_i$. With $f_1$ conveniently chosen, $\nabla t = \bigo_{p_0} (1)$ and with the right $f_2$, we get~\eqref{eq:normal_frame}. 


Then
\begin{equation}\label{eq:E_linearise}
E(x,y) = e^{ z(x) \cdot \con{z} (y) -\frac{1}{2} \left(|z(x)|^2 + |z (y)|^2\right)} t(x) \otimes \con{t(y)} + \bigo_{\left(p_0, p_0\right)} (3).
\end{equation}
A similar expression already appeared in the description of the Schwartz kernels of the operators of $\widetilde{\symb} (\C^n)$. Indeed, let $L_{\C^n}= \C^n \times \C$ be the trivial holomorphic line bundle equipped with the metric such that the frame $s(z) = (z,1)$ has a pointwise norm $|s(z)|^2 = e^{-|z|^2}$. Then it is natural to interpret the elements of $\pol (\C^n)$ as sections of $L_{\C^n}$, because the scalar product~\eqref{eq:scal_prod_2} is the integral of the pointwise scalar product $(f (z) s(z), g(z) s(z)) = f(z) \con{g} (z) e^{-|z|^2}$. Furthermore, in the integral~\eqref{eq:sch_ker}, we can interpret $e^{-|v|^2} f(v)$ as the pointwise scalar product $(f(v) s(v), s(v))$. In other words, the Schwartz kernel of $V$ is
\begin{equation}\label{eq:2pin-e_cn-u}
(2\pi)^{-n} E_{\C^n} (u,v) q(u-v) \quad \text{ with }
\quad E_{\C^n}(u,v) =e^{u \cdot \con{v}} s(u) \otimes \con{s}(v).
\end{equation}
Now equip $L_{\C^n}$ with its Chern connection, that the unique connection compatible with both the holomorphic and Hermitian structures. Then $\nabla s = - \sum \con{z}_i dz_i \otimes s$. So the curvature is $\frac{1}{i} \om_{\C^n} $ with $\om_{\C^n} = i \sum dz_i
\wedge d \con{z}_i$. And if $t $ is the unitary frame $t(z) = e^{|z|^2/2} s (z)$, we have $\nabla t = \frac{1}{2} \sum (z_i d \con{z}_i - \con{z}_i d z_i) \otimes t$ and
\[
E_{\C^n}(u,v) = e^{ u \cdot \con v - \frac{1}{2} \left(
|u|^2+ |v|^2\right)} t(u) \otimes \con{t(v)},
\]
the same formula as~\eqref{eq:E_linearise}.

\subsection{Schwartz kernel expansion} \label{sec:global-expansion}


We consider operator families $(P_k : \Ci(M, L^k \otimes A) \rightarrow \Ci(M, L^k \otimes B), \; k \in \N)$ having smooth Schwartz kernels. Recall the notations introduced in the beginning of Section~\ref{sec:operators}. In particular, $\| P_k \|$ is the operator norm whereas $|P_k|$ is the function of $M^2$ sending $(x,y)$ into $|P_k(x,y)|$.

Let $E$ be a section of $L \boxtimes \con L $ satisfying~\eqref{eq:hypotheseE} and~\eqref{eq:normE} and $b \in \Ci (M^2, B \boxtimes \con A)$. Then, viewing $(L^k \otimes B) \boxtimes (\con{L}^k \otimes \con{A})$ as $(L \boxtimes \con L)^k \otimes (B \boxtimes
\con{A})$, we introduce the operator family $(P_k)$ with Schwartz kernels
\begin{gather}\label{eq:prel}
P_k (x,y) = \left(\frac{k}{2\pi} \right)^n E^k(x,y) b (x,y)
\end{gather}
The pointwise norms of $P_k$ depend in an essential way on the vanishing order of $b$ along the diagonal. If $m \in \N$, we write $b = \bigo (m)$ to say that all the derivatives of $b$ of order $\leqslant m-1$ are zero at each point of the diagonal. Recall that
\[
\psi (x,y) = - 2 \ln |E(x,y)|
\]
is positive outside the
diagonal, and vanishes to second order along the diagonal with a Hessian non-degenerate in the transverse direction.

\begin{lemm}\label{lem:estimate}
If $P_k$ is given by~\eqref{eq:prel} with $b = \bigo (m)$, then  $|P_k | = \bigo (k^{n-\frac{m}{2}} e^{-k
\frac{\psi}{4}})$  and $\| P_k \| = \bigo (k^{- \frac{m}{2}})$
\end{lemm}

\begin{proof}
Since $b = \bigo (m)$, $ | b | = \bigo (\psi ^{\frac{m}{2}}
)$. So 
\[
|P_k | \leqslant C k^n e^{-k \frac{\psi}{2}} \psi ^{\frac{m}{2}}
= C k^{n- \frac{m}{2}} e^{-k \frac{\psi}{4}} e^{-k \frac{\psi}{4}} (k\psi) ^{\frac{m}{2}}
\leqslant C' k^{n- \frac{m}{2}} e^{-k \frac{\psi}{4}}
\]
because $t \rightarrow e^{-\frac{t^2}{2} } t^m$ is bounded on $\R_{\geqslant 0}$. This proves the first estimate and can be written locally in a coordinate system as
\[
|P_k (x,y) | \leqslant C k^{n - \frac{m}{2}} e^{-k|x-y|^2 /C}.
\]
So $\int_M |P_k(x,y)| \, d\mu_M (y)$ and $\int_M |P_k (x,y)|\, d\mu_M (x)$ are both $\leqslant C k^{-\frac{m}{2}}$ and the operator norm estimate follows from Schur test.
\end{proof}


Consider now $(P_k) \in \lag (A,B)$. Recall that by definition~\eqref{eq:exp_Lag_bis} we have for any $N \in \N$, 
\begin{gather}\label{eq:expansion}
P_k (x,y) = \left(\frac{k}{2\pi} \right)^n E^k(x,y) \sum_{\substack{\ell\,\in\,\Z, \\
\ell + m (\ell)\,\leqslant\,N }} k^{- \frac{\ell}{2}} b_{\ell} (x,y) + R_{N,k} (x,y)
\end{gather}
where
\begin{enumerate}\romanenumi
\item $m : \Z \rightarrow \N \cup \{\infty \} $ is such that for any $N$, $\{ \ell/\; \ell + m(\ell) \leqslant N \}$ is finite, and $\ell + m (
\ell) \geqslant 0$ for any $\ell$.
\item $(b_{\ell})_{\ell\,\in\,\Z}$ is a family of $\Ci (M^2, B \boxtimes \con{A})$ such that $b _{\ell} = \bigo (m(
\ell))$ for any $\ell$.
\item $|R_{N,k} (x,y) | = \bigo (k^{n - \frac{N+1}{2}})$ uniformly on $M^2$.
\end{enumerate}
By Lemma~\ref{lem:estimate}, $| E^k(x,y) k^{-\frac{\ell}{ 2}} b_{\ell}(x,y) | \in \bigo (k^{ - \frac{1}{2} (\ell + m (\ell))})$. So the expansion~\eqref{eq:expansion} is consistent in the sense that passing from $N$ to $N+1$, we add new terms $k^{-\frac{\ell}{2}} b_\ell$ such that $\ell + m (\ell) = N+1$, which contribute to $P_k (x,y)$ with a $\bigo (k^{n -\frac{1}{2}(N+1)})$.

\begin{lemm}\label{lem:estim_expansion}
If the expansion~\eqref{eq:expansion} holds, then for any $q>0$
\[
|P_k | = \bigo \left(k^{n} e^{- k
\frac{\psi }{4}} \right) + \bigo \left(k^{-q}\right),  \qquad \| P_k \| = \bigo (1)
\]
Similarly the remainders $R_{N,k}$'s satisfy for any $q>0$,
\[
 |R_{N,k} | = \bigo \left(k^{n- \frac{N+1}{2}} e^{- k
\frac{\psi }{4}} \right) + \bigo\left(k^{-q}\right),  \qquad \| R_{N,k} \| = \bigo
\left(k^{-\frac{N+1}{2}}\right).
\]
\end{lemm}

\begin{proof}
To prove the first estimate, we use~\eqref{eq:expansion} with $N$ sufficiently large so that $| R_{N,k} | = \bigo (k^{-q})$, and the result follows from Lemma~\ref{lem:estimate} because $\ell + m (\ell) \geqslant 0$. The operator norm estimate is proved similarly by choosing $N$ so that $|R_{N,k} | =\bigo (1)$ which implies that $\| R_{N,k} \| = \bigo (1)$. The proof for the $R_{N,k}$ is essentially the same.
\end{proof}

We next show that in the expansion~\eqref{eq:expansion}, we can choose any section~$E$ satisfying the assumptions given in Section~\ref{sec:section-ee}.

\begin{lemm}\label{lem:changement_E}
Assume~\eqref{eq:expansion} holds and let $E'$ be a section satisfying~\eqref{eq:hypotheseE} and~\eqref{eq:normE}. Then there exists a family $(b'_{\ell})$ of $\Ci (M^2, B \boxtimes \con A)$ such that~\eqref{eq:expansion} holds with $E'$ and $b'_{\ell}$ instead of $E$ and $b_{\ell}$.
\end{lemm}

\begin{proof}
Observe first that~\eqref{eq:expansion} holds outside the diagonal if and only $|P_k (x,y)|$ is in $\bigo (k^{-\infty})$ outside the diagonal, and this condition is clearly independent of the choice of $E$ and the $b_{\ell}$'s. On a neighborhood of the diagonal, we have $E = e^{g} E'$. Since $g \in \bigo (3)$, we can assume that  $|g| \leqslant \psi/8$.  Let us write
\[
P_k = \left(\frac{k}{2 \pi}\right)^n E^k
b_{N,k} + \bigo \left(k^{n- \frac{N+1}{2}}\right) \quad \text{with } \quad b_{N,k} = \sum_{\ell + m (\ell)
\,\leqslant\,N} k^{-\frac{\ell}{2}}b_{\ell}.
\]
By Lemma~\ref{lem:estimate},  $|E|^k b_{N,k} = \bigo (e^{-k \frac{\psi}{4}})$.  Using that
\[
\exp z = \sum_{p = 0 }^N \frac{z^p }{p!} + r_N (z) \quad \text{with} \quad |r_{N} (z) | \leqslant \frac{|z|^{N+1} }{ (N+1)!} e^{|\op{Re} z |},
\]
we deduce from $E^k = e^{kg } (E')^k$ that
\begin{gather}\label{eq:fin}
E^k b_{N,k} = (E')^k b_{N,k} \sum_{p = 0 }^N k^p \frac{g^p}{p!} + R_{N,k}
\end{gather}
where 
\[
| R_{N,k} | \leqslant C_N e^{-k \frac{\psi}{4} } |kg|^{N+1} e^{ k
|\op{Re} g |} \leqslant C_N e^{-k \frac{\psi}{8} }
|kg|^{N+1} = \bigo \left(k^{- \frac{N+1}{2}}\right)
\]
because $|\op{Re} g | \leqslant \psi/8$.  To conclude now, it suffices to define the $b'_{\ell}$ so that
\begin{gather}\label{eq:last_goal}
(E')^k \left[\sum_{\ell + m (\ell)
\,\leqslant\,N} k^{-\frac{\ell}{2}}b_{\ell} \right] \left[\sum_{p = 0 }^N k^p \frac{g^p}{p!} \right]= (E')^k \sum_{\ell + m ' (\ell)\,\leqslant\,N} k^{- \frac{\ell}{2}} b'_{\ell} + \bigo \left(k^{-\frac{N+1}{2}}\right)
\end{gather}
holds for any $N$. This suggests that each $b'_{\ell}$ should be equal to the infinite sum
\[
b_{\ell}
+ b_{\ell +2 } \, g + b_{\ell +4 } \frac{g^2}{2} + b_{\ell+ 6} \frac{g^3}{6}
+\ldots
\]
But by Lemma~\ref{lem:estimate}, the equality~\eqref{eq:last_goal} depends only on the class of $b'_{\ell}$ modulo $\bigo (N- \ell)$, so we can interpret these infinite sums as sums of Taylor expansions along the diagonal. Since $b_{\ell +2p} \, g^{p} = \bigo (m(\ell +2 p) + 3p) = \bigo (3p)$, by Borel lemma, there exists $b'_{\ell}$ such that for any $M$
\begin{gather}\label{eq:changement_E}
b'_{\ell} = \sum_{p=0}^M b_{\ell+ 2p } \frac{g^p}{p!} + \bigo (3(M+1))
\end{gather}
So $b'_{\ell} = \bigo (m '(\ell))$ with $ m' (\ell) := \op{min} \{ m(\ell + 2p) + 3p, \; p \in \N \}.$ We easily check that $m'$ satisfies the same condition as $m$. We finally deduce~\eqref{eq:last_goal} by removing with Lemma~\ref{lem:estimate} all the coefficients leading to a $\bigo (k^{-\frac{N+1}{2}})$.
\end{proof}

Suppose now we have an open set $U$ of $M$, and functions $u_i \in \Ci (U^2)$, $i =1, \,\ldots,\, 2n$ vanishing along the diagonal and such that for any $y \in U$, $(u_i(
\cdot, y))$ is a coordinate system on a neighborhood of $y$. Then we can write the Taylor expansions along the diagonal as follows: any $ f \in \Ci (U^2)$ has a decomposition
\begin{gather}\label{eq:taylor_expansion}
f(x,y) = \sum_{m=0}^M f_m (y, u(x,y)) + \bigo (M+1)
\end{gather}
where each $f_m (y,\xi)$ is homogeneous polynomial in $\xi$ with degree $m$. This can be done also for sections of $B \boxtimes \con{A}$ by introducing frames of $A$ and $B$ on $U$, so that $\Ci (U^2, B \boxtimes \con{A}) \simeq \Ci (U^2, \C^r)$.

\begin{lemm}\label{lem:new_expansion}
The expansion~\eqref{eq:expansion} holds on $U^2$ if and only if there exists a sequence $(a_p)$ of $\Ci (U \times \R^{2n}, \C^r)$, each $a_p (x,\xi)$ being polynomial in $\xi$, such that for any $N$
\begin{equation}\label{eq:new_expansion}
P_k (x,y) = \kpi E^k (x,y) \sum_{p=0}^{N} k^{-\frac{p}{2}} a_p \left(y, k^{\frac{1}{2}} u (x,y)\right) + \bigo \left(k^{n- \frac{N+1}{2}}\right).
\end{equation}
\end{lemm}

The remainders in~\eqref{eq:new_expansion} satisfy the same pointwise estimates as the $R_{N,k}$ given in Lemma~\ref{lem:estim_expansion}, the proof is identical.

\begin{proof}
If~\eqref{eq:expansion} holds on $U^2$, writing the Taylor
expansion of each $b_{\ell}$ as in~\eqref{eq:taylor_expansion}, we have by Lemma~\ref{lem:estimate}
\begin{align*}
E^k (x,y) k^{-\frac{\ell}{2}} b_{\ell} (x,y) &= E^k (x,y) k^{-
\frac{\ell}{2}} \sum_{m = m (\ell) }^{N- \ell } b_{\ell,m} (y, u(x,y)) +
\bigo\left(k^{-\frac{N+1}{2}}\right)\\
&=E^k (x,y) k^{-\frac{\ell +m }{2}} \sum_{m = m (\ell) }^{N- \ell } b_{\ell,m} \left(y,k^{\frac{1}{2}} u(x,y)\right) +
\bigo\left(k^{-\frac{N+1}{2}}\right)
\end{align*}
So we obtain~\eqref{eq:new_expansion} with $ a_p = \sum _{\ell + m (\ell)
\leqslant p } b_{\ell, p - \ell} $, this sum being finite because of the assumption satisfied by $m(\ell) $.

Conversely, starting from the $a_{p}$'s, for each $\ell \in \Z$, we construct by Borel summation a function $b_{\ell}$ such that $ b_{\ell} (x, y)=
\sum _{m =0 }^{M} a _{ \ell+m, m}(y,u (x,y)) + \bigo (M+1)$ for all $M$, where by convention $a_{p}=0$ for $p<0$, and $a_{m+\ell, m}$ is the degree $m$ homogeneous component of $a_{m+\ell}$. We readily deduce the expansion~\eqref{eq:expansion} from~\eqref{eq:new_expansion} by using Lemma~\ref{lem:estimate} again.

Observe that $b_{\ell} =
\bigo(m(\ell))$ with $m(\ell)$ the smallest $m$ such that $a_{\ell + m, m } \neq 0$. Since $a_{p} = 0$ for $p<0$, we have $\ell + m(\ell)
\geqslant 0$. Furthermore, $\ell + m (\ell) \leqslant N$ happens only if there exists $m \leqslant N - \ell$ such that $a_{\ell + m, m } \neq 0$, that is if there exists $p \leqslant N$ such that $ a _{p, p-\ell } \neq 0$, so necessarily $ p - \ell \leqslant d (p)$ where $d(p)$ is the degree of $a_p$. So $\ell + m (\ell) \leqslant N$ implies that $\ell \geqslant \min \{ p - d(p) / \; p =0,\,\ldots,\, N \}$. So $\ell + m (\ell) \leqslant N$ only for a finite number of $\ell$.
\end{proof}
 We have essentially proved Proposition~\ref{prop:un_de_plus}. Here are the details.

\begin{proof}[Proof of Proposition~\ref{prop:un_de_plus}]
Identify $U$ with an open convex set of $\R^{2n}$, then the functions $u_i(x,y)= x_i -y_i$ satisfy the above conditions. And for $(x,y) = (x' + \xi', x')$, we have $a_p (y, k^{\frac{1}{2}} u(x,y)) = a_p (x', k^{\frac{1}{2}} \xi')$, so the expansions~\eqref{eq:new_expansion} and~\eqref{eq:expansion_kernel} are the same when $E=e^{-\varphi}$. Now Proposition~\ref{prop:un_de_plus} follows from Lemma~\ref{lem:new_expansion}, the local version of Lemma~\ref{lem:changement_E} and the fact that $E=e^{-\varphi}$ in~\eqref{eq:expansion_kernel} satisfies the conditions~\eqref{eq:hypotheseE}.
\end{proof}

It is the good place to prove Lemma~\ref{lem:parite} on the characterization of the parity in terms of local expansions.
\begin{proof}[Proof of Lemma~\ref{lem:parite}]
This follows from the relation between the coefficients $a_{p}$ and the coefficients $b_{\ell}$ given in the proof of Lemma~\ref{lem:new_expansion}. For instance, if $b_{\ell} =0$ for any odd integer $\ell$, then $b_{\ell, p-\ell} \neq 0$ only for even $\ell$ and in this case it has the same parity as $p$, so $ a_p = \sum _{\ell + m (\ell)
\,\leqslant\,p } b_{\ell, p - \ell} $ has the same parity as $p$. Conversely, if $a_{p}$ has the same parity as $p$ for any $p$, then $a_{ \ell+m, m } =0$ for odd $\ell$, so $b_{\ell} $ vanishes to infinite order on the diagonal for odd $\ell$, so we can assume that $b_{\ell} =0$. The proof for odd elements is the same.
\end{proof}



\subsection{Filtration and symbol}
For an operator $(P_k) \in \lag (A,B)$, we have two different ways of writing the expansion of its Schwartz kernel: a global one~\eqref{eq:expansion} with coefficients $b_\ell$ and a local one~\eqref{eq:new_expansion} with coefficients $a_p$. We now discuss the uniqueness of these coefficients. Recall that $(P_k) \in \lag_q (A,B)$ if in all the local expansions~\eqref{eq:new_expansion}, the coefficients $a_{p} $ are zero for $p < q$.

\begin{prop}\label{prop:filtration-symbol}\leavevmode
%\hspace{1em}
\begin{enumerate}
\item\label{prop5.5.1} In the local expansions~\eqref{eq:new_expansion}, the coefficients $a_{p}$ are uniquely determined by the section~$E$, the functions $(u_i)$ and the frames of $A$ and $B$.
\item\label{prop5.5.2} In the global expansion~\eqref{eq:expansion}, the Taylor expansions of the coefficients $b_{\ell}$ along the diagonal are uniquely determined by the section~$E$.
\item\label{prop5.5.3} $ (P_k) \in \lag_q (A,B)$ iff $\bigl(\forall \, \ell \in\Z, \; b_{\ell} = \bigo (q - \ell) \bigr)$ iff $| P_k | = \bigo (k^{n-\frac{q}{2}}) $.
\item\label{prop5.5.4} If $(P_k) \in \lag_q (A,B)$, then the coefficient $a_q$ of the local expansion~\eqref{eq:new_expansion}, viewed as a section of $ \pol (TM) \otimes B \otimes \con{A} \rightarrow U$ does neither depend on $E$ nor on the functions $(u_i)$. Furthermore,
\begin{gather}\label{eq:a_n_fonction_b_ell}
a_q = \sum_{\ell + m (\ell) = q} b_{\ell, q - \ell}
\end{gather}
where the $b_{\ell}$ are the coefficients of the global expansion~\eqref{eq:expansion} and $b_{\ell, q-\ell}$ is defined as in~\eqref{eq:taylor_expansion}.
\end{enumerate}
\end{prop}

\begin{proof}
Assertions 1, 2 and 3 follow from the following facts: Let $f_0$, \ldots, $f_q$ in $\Ci (M^2)$. Let $\psi = -2
\ln |E|$. Then
\begin{gather}\label{eq:esti_iff}
e^{-k \psi} \sum_{\ell =0}^q k^{-\frac{\ell}{2}} f_{\ell} = \bigo \left(k^{-\frac{q}{2}}\right) \quad \Leftrightarrow \quad f_0 \in \bigo (q),\, \ldots,\, f_q \in \bigo (0).
\end{gather}
Indeed, recall that $\psi \geqslant 0 $, is in $\bigo (2)$ and its Hessian is non-degenerate in the direction transverse to the diagonal. The converse of~\eqref{eq:esti_iff} follows from the same proof as Lemma~\ref{lem:estimate}. The direct sense of~\eqref{eq:esti_iff} follows from~\cite[Proposition~2.4 and Remark~2.5]{oim}.


From this, we deduce that $|P_k | = \bigo (k^{n -\frac{q}{2}})$ iff $b_{\ell} = \bigo (q-\ell)$ for any $\ell$. Since $a_{p} = \sum_{\ell + m (\ell) \leqslant p} b_{\ell, p -\ell}$ by the proof of Lemma~\ref{lem:new_expansion}, $(b_\ell = \bigo (q-\ell)$ for any $\ell$) iff $(a_p = 0,$ for any $p < q$). This last condition is the definition of $\lag _q (A,B)$. We have just proved Assertion 3. This implies that $|P_k | =\bigo (k^{-\infty})$ iff ($a_{p} = 0$ for any $p$) iff $(b_{\ell} = \bigo(
\infty)$ for any $\ell$), which proves Assertions 1 and 2.

For the fourth assertion, since $P \in \lag_q(A,B)$, we have $b_{\ell} = \bigo (q-\ell)$ for any $\ell$, so we can assume that $\ell + m (\ell) \geqslant q$, so
\[
a_q =
\sum_{\ell + m (\ell) \,\leqslant q } b_{\ell, q- \ell} =\sum_{\ell + m (\ell
) = q } b_{\ell, q- \ell}.
\]
Since $b_{\ell} = \bigo (q-\ell)$, the $(q-\ell)^{\rm th}$ order term in the Taylor expansion of $b_{\ell}$ is intrinsically defined as a function
\[
\xi \in T_xM \rightarrow b_{\ell, q- \ell}
(x,\xi) \in B_x \otimes \con{A}_x,
\]
so we can view $a_q(x,\cdot)$ as an element
of $\pol (T_xM) \otimes B_x \otimes \con{A}_x$.

It remains to prove that for $\ell + m (\ell) =q$, $b_{\ell, q-\ell}$ does not depend on the choice of $E$. With the notation of the proof of Lemma~\ref{lem:changement_E}, this amounts to prove that $b'_{\ell} = b_{\ell} +
\bigo (q-\ell +1)$. This follows from~\eqref{eq:changement_E}, because $b_{\ell + 2p } g^p = \bigo (m (\ell+ 2 p) +3 p)$ and $m(\ell +2p) + 3p \geqslant q - (\ell+ 2p) + 3p = q + p - \ell \geqslant q + 1 -\ell$ when $p
\geqslant 1$.
\end{proof}


We are now ready to define the symbol map
\[
\si_q : \lag_q (A,B) \rightarrow \Ci (
\symb(M) \otimes \op{Hom} (A,B)).
\]
First, for any $x \in M$, $T_xM$ is a Hermitian space, so it has an associated algebra $\widetilde{\symb} (T_xM)$ with a map $\op{Op}:\pol (T_xM) \rightarrow \widetilde{\symb} (T_xM)$ as in Section~\ref{sec:landau-levels-cn}.

For any $(P_k)$ in $\lag_q (A,B)$, by the fourth assertion of Proposition~\ref{prop:filtration-symbol}, $a_q(x,
\cdot) \in \pol (T_xM) \otimes B_x \otimes \con{A}_x$. Identifying $B_x
\otimes \con{A}_x$ with $ \op{Hom} (A_x, B_x)$, we set
\begin{gather}\label{eq:def_si_tilde}
\widetilde{ \si}_q (P) (x) := \op{Op} (a_q(x, \cdot)) \in
\widetilde{\symb} (T_xM) \otimes \op{Hom}(A_x, B_x)
\end{gather}
Recall that we have an isomorphism $\widetilde{\symb} (T_xM) \simeq \symb (T_xM)$ defined by restriction from $\pol (T_xM)$ to $\D (T_xM)$.

\begin{defi}\label{def:symbol}
$\si _q (P) (x) \in \symb (T_xM) \otimes \op{Hom}(A_x, B_x)$ is defined as the restriction of $\widetilde{\si}_q (P) (x)$.
\end{defi}

\subsection{Proofs of the results of Section~\ref{sec:operators}} \label{sec:proofs-results}
We now give the proof of Proposition~\ref{prop:lag}, Theorem~\ref{theo:lag} and Theorem~\ref{theo:parity}.

\begin{proof}[Proof of Proposition~\ref{prop:lag}]
The first assertion is an easy consequence of the definition of $\lag_q (A,B)$ by the local expansions. In the second assertion, the characterisation in terms of pointwise norm is the third assertion of Proposition~\ref{prop:filtration-symbol}. By Lemma~\ref{lem:estimate} or Lemma~\ref{lem:estim_expansion}, every $(P_k) \in \lag_q(A,B)$ satisfies $\| P_k \| = \bigo (k^{-\frac{q}{2}})$. For the converse, it suffices to show that if $\si_0 (P) \neq 0$, then $\| P _k \| \geqslant c >0$. This is a consequence of Corollary~\ref{cor:norm_estim}. The third assertion is straightforward. The fourth assertion is a variation on Borel Lemma, cf. for instance~\cite[Proposition~2.1]{oim}.
\end{proof}

\begin{proof}[Proof of Theorem~\ref{theo:lag}]
In Definition~\ref{def:symbol}, we have defined a map
\[
\si_q : \lag_q(A,B) \rightarrow \Ci (M, \symb (M) \otimes \op{Hom}
(A,B)).
\]
having kernel $\lag_{q+1} (A,B)$ by  the injectivity of $\op{Op}$, cf. Proposition~\ref{prop:tilde_symb_and_op}.  To prove that it is surjective, we show that for any $c \in \Ci (M, \pol (TM) \otimes B \otimes \con{A})$, there exists $P \in \lag_q (A,B)$ such that in the local expansions~\eqref{eq:new_expansion}, $a_q =c$. To do this, let $d \in \N$  be an upper bound of the degree of $c(x, \cdot)$ for any $x\in M$. For any  $m =0, \,\ldots,\, d$,  let $c_m(x,\cdot)$ be the homogeneous component with degree $m$ of $c(x, \cdot)$. Choose a section $b_{q-m}$ of $B \boxtimes
\con{A}$ vanishing to order $m$ along the diagonal and satisfying $b_{q-m} (x+ \xi, x) = c_m (x,\xi) + \bigo (m+1)$. Then we set 
\[
P_k (x,y) := \left(\frac{k}{2\pi} \right)^n E^k(x,y) \sum_{\ell = q-
d }^q k^{-
\frac{\ell}{2}} b_{\ell} (x,y).
\]

Since $b_\ell = \bigo (q-\ell)$ for any $\ell$, $(P_k) \in \lag_q (A,B)$ by Assertion~\eqref{prop5.5.3} of Proposition~\ref{prop:filtration-symbol}.  By~\eqref{eq:a_n_fonction_b_ell}, we have $a_q =  \sum_{\ell = q-d}^q b_{\ell, q- \ell} = \sum_{m=0}^{d} c_{m}  = c$, as was to be proved.

Let us prove the remaining assertions of Theorem~\ref{theo:lag}. Let $P \in \lag_q (A,B)$. Assertion~\eqref{theo2.3.1}, that is $\si _q (P) = \si_0 (k^{q/2} P)$, follows directly from the local expansions~\eqref{eq:new_expansion}. Let us prove Assertion~\eqref{theo2.3.2}. If $f \in \Ci (M,
\op{Hom} (B,C))$, then the Schwartz kernel of $ P_k'= f \circ P_k$ is $P_k' (x,y) = f(x) (P_k (x,y))$, so $P_k'$ has the same expansion~\eqref{eq:expansion} as $P_k$ with $b'_{\ell} (x,y) = f(x)(b_{\ell}(x,y))$ instead of $b_{\ell}$, which implies that $P'_k$ belongs to $\lag_q (A,B)$ with the same function $\ell \mapsto m(\ell) $. Furthermore, with the notation~\eqref{eq:taylor_expansion}, $b'_{\ell, m(\ell)} (x, \cdot) = f(x) b_{\ell, m(\ell)} (x, \cdot)$, which implies by~\eqref{eq:a_n_fonction_b_ell} that $\si_q(P')(x) = f(x) \circ \si_q(P)(x)$.

Let us prove Assertion~\eqref{theo2.3.3}. Since $P_k^* (x,y) = \con{P_k (y,x)}$, the Schwartz kernel of $P_k^*$ has the expansion~\eqref{eq:expansion} with $E' (x,y) =
\con{E (y,x)}$ instead of $E$ and $b'_{\ell} (x,y) = \con{b_{\ell} (y,x)}$ instead of $b_{\ell}$. By~\eqref{eq:symE}, we deduce that $(P_k^*) \in \lag_q (B,A)$. Furthermore, $b'_{\ell, m(\ell)} (x, \xi) = \con{b}_{\ell,m(\ell)}(x, - \xi)$ so $a_q' (x,\xi) = \con{a}_q (x, -\xi)$. By~\eqref{eq:sch_ker}, $\op{Op} (q)^* = \op{Op} (r)$ with $r(\xi) = \con{q} (- \xi)$, so $\si_q(P^*) = \si_q(P)^*$.

Let us prove Assertion~\eqref{theo2.3.5}. By~\eqref{eq:new_expansion},
\[
P_k(x,x) = \frac{k^{n-q/2}}{ (2\pi)^n} \left(a_q(x,0) + \bigo \left(
k^{-\frac{1}{2}}\right) \right)
\]
and by the first equation of~\eqref{eq:rel_bizare}, $a_q(x,0) = \op{tr} (\si_q(P) (x))$.

 Let us prove half of Assertion~\eqref{theo2.3.6}. More precisely, we will deduce from Assertion~\eqref{theo2.3.4} that for any $P \in \lag (A,B)$, we have
\begin{gather}\label{eq:tbp}
\limsup_{k\,\rightarrow\,\infty} \| P_k \| \leqslant \sup_{x\,\in\,M} \|
\si_0 (P)(x) \|
\end{gather}
With the lower bound provided by Corollary~\ref{cor:norm_estim}, this will show Assertion 6. Let $f := \si_0(P)^* \si_0 (P)$. Let $m \in \N$ be sufficiently large so that $\pi f \pi = f $ where for any $x \in M$, $\pi (x)
\in \symb (T_xM)
\otimes \op{End} A_x $ is the selfadjoint projector onto $\D_{\leqslant\,m } (T_xM) \otimes A_x$. Then for any $C > \sup_{x\,\in\,M} \|
\si_0 (P)(x) \|$, there exists a symbol $g \in \Ci (M, \symb (M) \otimes
\op{End} A)$ such that $g ^* = g$, $\pi g \pi = g $, $g^2 = C^2 \pi - f$. Indeed, $g(x)$ is zero on $\D_p (T_xM) \otimes A_x$ for any $ p > m$, and $g(x) $ is the positive square root of $C - f$ on $\D_{\leqslant\,m} (T_xM) \otimes A_x$. Let $\Pi$ and $Q$ in $\lag (A)$ be self-adjoint and having symbol $\pi$ and $g$ respectively. Then by Assertion 4,
\[
C^2 \Pi^2 - P^* P = Q^2 + R
\]
with $R \in \lag_1 (A,B)$. So
\begin{align*}
\| P_k \Psi \| ^2 = \left\langle P_k^* P_k \Psi, \Psi \right\rangle = & C^2 \| \Pi_k
\Psi \|^2 - \| Q_k \Psi \|^2 - \left\langle R_k \Psi, \Psi \right\rangle \\
\leqslant & \left(C^2 \| \Pi_k \|^2 + C' k^{-\frac{1}{2}}\right) \| \Psi \|^2.
\end{align*}
by Assertion~\eqref{prop2.2.2} of Proposition~\ref{prop:lag}. Since $\pi^2 = \pi$, by Assertion~\eqref{prop2.2.4} and Assertion~\eqref{prop2.2.2} of Proposition~\ref{prop:lag} again, $\Pi_k^2 -
\Pi_k = \bigo (k^{-\frac{1}{2}})$ so $\| \Pi_k \| \leqslant 1 + \bigo (k^{-\frac{1}{2}})$. Consequently, $\| P_k \| \leqslant C (1 + \bigo (k^{-\frac{1}{2}}))$ which implies~\eqref{eq:tbp}.




It remains to prove Assertion~\eqref{theo2.3.4}. Let $(P'_k) \in
\lag_{q'} (B,C)$. We will prove that $Q_k = P'_k \circ P_k$ belong to $\lag _{q'' } (A,C)$ with $q'' = q + q'$ and compute its symbol. Since the composition of operators with kernels in $\bigo (k^{-p})$ and $\bigo (k^{-\ell})$ respectively, has a kernel in $\bigo (k^{-(p+\ell)})$, we can consider each summand of the expansions~\eqref{eq:expansion} for $P_k$ and $P'_k$ separately. In other words, we can assume that $P_k = (\frac{k}{2\pi} )^n E^k f$ and $P'_k = (\frac{k}{2\pi})^n E^k f'$ with $f = \bigo (q)$ and $f' = \bigo (q')$. So
\begin{gather}\label{eq:R_kernel}
Q_k (x,z) = \left(\frac{k}{2\pi} \right)^{2n} \int_M (E(x,y) \cdot E(y,z))^k g(x,y,z) d \mu_M (y)
\end{gather}
with $g(x,y,z) = f(x,y) f' (y,z)$. Observe that $g$ vanishes to order $q''$ along $\Si = \{ (x,y,z) \in M^3, \; x= y = z \}$.



By~\eqref{eq:normE}, $|E(x,y) \cdot E(y,z) | < 1 $ if $(x,y,z) \notin \Si$. This implies first that the Schwartz kernel of $(Q_k)$ is in $\bigo (k^{-\infty})$ outside the diagonal. Furthermore, to compute $Q_k$ on a neighborhood of $(p,p)$ up to a $\bigo (k^{-\infty})$, we can reduce the integral~\eqref{eq:R_kernel} to a neighborhood of $p$. So we can work locally.

Introduce a local orthonormal frame $(\partial_i, i =1,\, \ldots,\, n)$ of $T^{1,0}M$ on an open neighborhood $U$ of $p$ in $M$. Let $\si_i \in \Ci (U^2)$, $i =1,\,\ldots n$ be such that
\begin{gather}\label{eq:dsi_i--partial_j}
d\si_i (\partial_j, 0) = \delta_{ij} + \bigo (1), \qquad d\si_i \left(\con{\partial}_j, 0\right) = \bigo (1)
\end{gather}
for any $i$ and $j$. Observe that if the $z_i$ are coordinates as in~\eqref{eq:coordonnees_z}, then
\begin{gather}\label{eq:si_z}
\si_i (x,y) = z_i(x) - z_i (y) + \bigo_{(p_0, p_0)} (2).
\end{gather}
So we can use the functions $u_i = \op{Re} \si _i$ and $u_{i+n} = \op{Im} \si_i$ when we write the Taylor expansion~\eqref{eq:taylor_expansion} and the local expansion~\eqref{eq:new_expansion}.



Restricting $U$ if necessary, we can assume that for any $z$, the map $ y\in U \rightarrow (\si_i (y,z)) \in \C^n$ is a diffeomorphism onto its image. Let $\mu_z$ be the pull-back of the volume $\mu_n $ by this map. By~\eqref{eq:si_z}, we have $\mu_M (y) = \rho (y,z) \mu_z (y)$ with $\rho \in \Ci (U^2)$ satisfying $\rho (y,y) = 1$.

Now using the expressions~\eqref{eq:E_linearise} and~\eqref{eq:si_z}, we readily prove that
\[
E(x,y) \cdot E(y,z) = e^{ \varphi (x,y,z) + r(x,y,z) } E
(x,z),
\]
where $r(x,y,z) = \bigo_{\Si} (3)$ and $ \varphi(x,y,z) = (\si (x,z) - \si (y,z))\cdot \con{\si} (y,z) $. Arguing as in the proof of Lemma~\ref{lem:changement_E}, it comes that
\[
(E(x,y) \cdot E(y,z))^k = E^k(x,z) e^{k \varphi (x,y,z)} \sum_{\ell =0 }^{N}
\frac{k^{\ell}}{\ell !} (r (x,y,z))^{\ell} + \bigo \bigl(k^{-\frac{1}{2} (N+1) } \bigr)
\]
so the integrand of~\eqref{eq:R_kernel} is equal to
\[
E^k(x,z) e^{k \varphi (x,y,z)} \sum_{\ell =0 }^{N}
k^{\ell} g_{\ell} (x,y,z) \; d\mu_z (y) + \bigo \left(k^{-\frac{1}{2} (q'' + N+1)} \right)
\]
with $g_{\ell} (x,y,z) = \rho(y,z) g (x,y,z) (r (x,y,z)) ^\ell / (\ell !) = \bigo_{\Si} (q''+ 3\ell)$.

For any $z \in U$, we write the Taylor expansion of $(x,y) \rightarrow g_{\ell}(x, y,z)$ at $(z,z)$ with the coordinates system
\[
(x,y) \rightarrow \op{Re} \si_i (x,z), \op{Im} \si_i (x,z),
\op{Re} \si_i (y,z), \op{Im} \si_i (y,z).
\]
We obtain
\[
g_\ell(x,y,z) = \sum_{m= q'' + 3 \ell}^{p} h_{\ell,m} (z, \si (x,z), \si (y,z))
+ \bigo_{\Si} (p+1)
\]
with $h_{\ell,m} (z, \xi, \eta)$ homogeneous polynomial in $\xi, \eta$ with degree $m$. Arguing as in Lemma~\ref{lem:new_expansion}, we obtain that
\begin{gather}\label{eq:dev_loc_R}
Q_k (x,z) = \left(\frac{k}{2\pi} \right)^n E^k(x,z) \sum_{\ell= 0}^{N}
\sum_{m= q''+ 3 \ell}^{q''+ 2 \ell +N} k^{\ell - \frac{m}{2}} I_{\ell,m} (x,z) + \bigo \left(k^{-\frac{1}{2} \left(q'' + N+1\right)} \right)
\end{gather}
with
\[
I_{\ell,m} (x,z) = \left(\frac{k}{2\pi} \right)^n \int_U e^{k
\varphi (x,y,z) } h_{\ell,m} \left(z, k^{\frac{1}{2}} \si (x,z), k^{\frac{1}{2}}
\si (y,z)\right) \; d\mu_z (y)
\]
Set $u_i = \si_i (x,z)$ and let us use the coordinates $v_i = \si_i (y,z)$ for the integration so that $\varphi (x,y,z) = u \cdot \con v - |v|^2$ and $d\mu_z (y) = |dv d \con{v}|$. It comes that $ I_{\ell, m} (x,z) = J_{\ell, m} (z, k^{\frac{1}{2}} \si (x,z))$ with
\begin{gather}\label{eq:j_ell_m}
J_{\ell, m } (z, u) = \left(\frac{k}{2\pi} \right)^n \int e^{ k^{\frac{1}{2}} u \cdot \con v - k |v|^2} h_{\ell,m} \left(z, u, k^{\frac{1}{2} } v\right) \; d \mu_n (v)
\end{gather}
where we integrate on a neighborhood of the origin in $\C^n$. We can actually integrate on $\C^n$ because this will modify $E^k(x,z) I_{\ell,m} (x,z)$ by a $\bigo (e^{-k/C})$. Indeed, $| E(x,z) | = e^{-\frac{1}{2} |u|^2} + \bigo (|u|^3)$ so $|E(x,z) | = \bigo (e^{-\frac{1}{3} |u|^2})$ so
\[
\left|E(x,z) e^{u
\cdot \con v - |v|^2}\right| = \bigo \left(e^{ - \frac{1}{3} |u|^2 + |uv| - |v|^2 }\right) = \bigo \left(
e^{ -\frac{1}{4} |v| ^2}\right)
\]
and we conclude by using that $ \int_{|v|\geqslant \epsilon} e^{- \frac{k}{4} |v|^2} |v|^m \; |dv d
\con{v}| = \bigo (e^{-k/C})$ for any $\ep >0$ and $m \in \N$.

Taking the integral~\eqref{eq:j_ell_m} over $\C^n$, it comes that
\begin{gather}\label{eq:j_ell-m}
J_{\ell,m} (z,u) = (2 \pi)^{-n} \int_{\C^n} e^{u \cdot \con v - | v| ^2 } h_{\ell, m } (z, u, v) \; d\mu_n (v)
\end{gather}
So $J_{\ell,m}$ does not depend on $k$. Furthermore it is polynomial in $u$. To see this, it suffices to view $h_{\ell, m } (z,u, v)$ as a polynomial in the variables $u-v$, $v$ and to compare with the formula~\eqref{eq:sch_ker}. So $Q_k (x,z)$ has the local expansion~\eqref{eq:new_expansion}, so $(Q_k)$ belongs to $\lag _{q''} (A, C)$. Its symbol is given by the leading order term in~\eqref{eq:dev_loc_R} which corresponds to $\ell =0$ and $m = q''$, that is
\[
\widetilde{\si}_{q''} (Q) (x) = \op{Op} (J_{0, q''} (x,
\cdot)).
\]
We can compute it in terms of the symbols of $P$ and $P'$ as follows: by~\eqref{eq:a_n_fonction_b_ell}, $\widetilde \si_{q} (P)(x) = \op{Op} (a_q(x, \cdot))$ where $\xi \rightarrow a_q(x,\xi)$ is the homogeneous polynomial of degree $q$ such that $f(x,y) = a_q (y, \si (x,y)) + \bigo (q+1)$. Similarly, $\widetilde \si_{q'} (P')(x) = \op{Op} (a'_{q'}(x, \cdot))$ with $f'(x,y) = a'_{q'} (y, \si (x,y)) + \bigo (q'+1)$. Now by~\eqref{eq:si_z}, $\si (x,y) = \si (x,z) - \si (y,z) + \bigo _{\Si} (2)$ and it comes that
\begin{align*}
g(x,y,z) & = a_q (y, \si (x,z) - \si (y,z)) a'_{q'} (z, \si (y,z)) +
\bigo_{\Si} (q'') \\
& = a_q (z, \si (x,z) - \si (y,z)) a'_{q'} (z, \si (y,z)) +
\bigo_{\Si} (q'')
\end{align*}
leading to $ h_{0, q''} (z, u,v) = a_q (z, u - v) a'_{q'} (z,v) $ and using first~\eqref{eq:sch_ker} and then~\eqref{eq:rel_bizare}, we have that
\begin{align*}
\widetilde{\si}_{q''} (Q)(x) & = \op{Op} (J_{0, q''} (x,
\cdot)) = \op{Op} \bigl(\op{Op} (a_q (x,\cdot)) a'_{q'} (x,\cdot) \bigr) \\
& = \op{Op} (a_q (x,\cdot)) \circ \op{Op} (a'_{q'} (x,\cdot)) = \widetilde \si_{q} (P) (x)
\circ \widetilde \si_{q'} (P')(x)
\end{align*}
as was to be proved.
\end{proof}

\begin{proof}[Proof of Theorem~\ref{theo:parity}]
The first assertion follows from the definition of the parity and Proposition~\ref{prop:filtration-symbol}. For the composition, it suffices to consider the case treated in the previous proof: we start from $(P_k)$ and $(P'_k)$ both even. Since $h_{\ell,m}(x,\cdot) $ has degree $m$, $J_{\ell,m}(x, \cdot)$ given by~\eqref{eq:j_ell-m} has the same parity as $m$, so by~\eqref{eq:dev_loc_R}, $(P'_k \circ P_k)$ is even. Last assertion is simply the fact that $\op{Op} (a_q(x,\cdot))$ has the same parity as $a_q (x,\cdot)$ by Proposition~\ref{prop:tilde_symb_and_op}.
\end{proof}

\subsection{Peaked sections}
In this section, we state and prove a generalisation of Proposition~\ref{prop:peaked-sections}. Consider an auxiliary bundle $A$. Let us choose a base point $x \in M$, with a coordinate chart $U$ centered at $x$, and a trivialisation $A|_U \simeq U \times A_x$. To any $f \in \pol (T_xM) \otimes A_x$, we associate the section of $L^k \otimes A$
\begin{gather}\label{eq:phi_peaked_section_+}
\Phi_k^f (x+\xi) = \left(\frac{k}{2\pi} \right)^{\frac{n}{2}} E^k (x+ \xi, x) \, f\left(k^{\frac{1}{2}} \xi\right) \, \psi (x+ \xi)
\end{gather}
where $E$ is chosen as in Section~\ref{sec:section-ee}, and $\psi \in \Ci_0(U)$ is equal to $1$ on a neighborhood of $x$.

 The space $\pol (T_xM) \otimes A_x$ has a natural scalar product obtained by tensoring the scalar product~\eqref{eq:scal_prod_2} of $\pol (T_xM)$ with the Hermitian metric of $A$.

\begin{prop}\label{prop:peaked-sections_++}\leavevmode
\begin{enumerate}
\item\label{prop5.7.1} For any $f$, $g \in \pol (T_x M) \otimes A_x$, $\bigl\langle
\Phi_k^f, \Phi_k^g \bigr\rangle = \langle f, g \rangle + \bigo (k^{-\frac{1}{2}})$.
\item\label{prop5.7.2} For any $f \in \pol (T_x M) \otimes A_x$ and $Q \in \lag (A,B)$, $ Q_k \Phi_k^f = \Phi_k^h + \bigo (k^{-\frac{1}{2}})$ where $h = \tilde{\si}_0 (Q)(x) \cdot f \in \pol (T_xM) \otimes B_x$.
\end{enumerate}
\end{prop}

In the second part, we used the symbol $\tilde{\si}_0 (P)$ defined in~\eqref{eq:def_si_tilde}, and $\Phi_k^{h}$ is defined as $\Phi_k^f$ with a trivialisation of $B$.

\begin{proof}
Consider the operator $P^f \in \lag (\C, A)$ with Schwartz
kernel 
\[
P^f_k (y+ \xi, y) = \left(\frac{k}{2\pi} \right)^n E^k (y+ \xi, y) \, f 
\left(k^{\frac{1}{2}} \xi\right)\, \psi (y+ \xi) \, \psi (y).
\]

On one hand, $(\frac{k}{2\pi})^{\frac{n}{2}} \Phi^f_k = P_k^f (\cdot,x)$. On the other hand, $\tilde{\si}_0(P^f)(x) = \op{Op} (f)$. So we can compute the scalar product of $\Phi^f_k$ and $\Phi^g_k$ as a composition of Schwartz kernels
\begin{align*}
\left (\frac{k}{2\pi} \right)^n \left\langle \Phi_k^f, \Phi_k^g \right\rangle & = \left(\left(P_k^g\right)^* P_k^f\right)(x,x) \\
& = \left(\frac{k}{2\pi} \right)^n \op{tr}_{\D (T_xM)} (\op{Op} (g)^*
\op{Op} (f)) + \bigo \left(k^{n-\frac{1}{2}}\right)
\end{align*}
by the last part of Theorem~\ref{theo:lag}.  To conclude, we have by the second equation of~\eqref{eq:rel_bizare} with $\op{Op} (q) = \op{Op} (g)^*$ that $ \op{Op} (g)^* \op{Op} (f) = \op{Op} (\op{Op} (g)^* f) $ and then by the first and third equations of~\eqref{eq:rel_bizare}
\[
\op{tr}_{\D \left(T_xM\right)} (\op{Op} \left(\op{Op} (g)^* f)\right) = \left(\op{Op} (g)^* f\right)(0)
= \langle f, g \rangle.
\]

The proof of the second part is similar, we have
\[
\left(\frac{k}{2\pi} \right)^{\frac{n}{2}} Q_k \Phi^f_k = \left(Q_k
P_k^f\right)(\cdot, x)
\]
By Theorem~\ref{theo:lag}, $(Q_k P_k^f) \in \lag (\C,B)$ with symbol at $x$
equal to
\[
\tilde{\si}_0 (Q) (x) \circ \op{Op} (f)
= \op{Op} \big(\tilde{\si}_0 (Q) (x) f\big) = \op{Op} (h)
\]
by~\eqref{eq:rel_bizare}.  So by the local expansion~\eqref{eq:expansion_kernel},
\[
\left(Q_k P_k\right)(\cdot,x) = \left(\frac{k}{2\pi} \right)^{\frac{n}{2}} \Phi_k^h + r_k
\]
with $r_k =R_{k} (\cdot, x)$ where $(R_{k}) \in \lag_ 1 (\C,B)$. Finally, $\| r_k \|^2 = (R_k^* R_k) (x,x) = \bigo (k^{n-1})$ because $(R_k^*R_k) \in
\lag_2 (\C)$ by Theorem~\ref{theo:lag}. 
\end{proof}

We deduce the following lower bound for the operator norm of operators of $\lag (A,B)$. If $\rho \in \symb_x (M) \otimes \op{Hom} (A_x, B_x)$, then we denote by $\| \rho \|$ the norm
\[
\| \rho \| = \sup \bigl\{ \| \rho f \| / \| f \|,\; f \in \D (T_xM)
\otimes A_x, \; f \neq 0 \bigr\}.
\]


\begin{coro} \label{cor:norm_estim}
For any $P \in \lag (A,B)$, we have
\[
\liminf_{k\,\rightarrow\,\infty} \| P_k \| \geqslant \sup_{x\,\in\,M} \|
\si_0 (P)(x) \|
\]
\end{coro}

\begin{proof}
By Proposition~\ref{prop:peaked-sections_++}, for any $f \in \D (T_xM) \otimes A_x$ non zero,
\[
\frac{\left\| P_k \Phi_k^f\right\|}{\left\| \Phi_k^f\right\|} = \frac{ \| \si_0 (P) (x) f
\|}{\| f \| } + \bigo \left(k^{-\frac{1}{2}}\right).
\]
So $ \liminf_{k\,\rightarrow\,\infty} \| P_k \| \geqslant \| \si_0 (P) (x) f
\| / \| f \|$.
\end{proof}

\section{Derivatives} \label{sec:derivatives}

The class $\lag (A, B)$ has been defined without any control on the derivatives of the Schwartz kernels. The reason was merely to simplify the exposition but in the applications it is natural and necessary to understand the composition of operators of $\lag(A, B)$ with covariant derivatives. We start from general considerations, then we define a subclass $\lag ^{\infty}(A,B)$ where the asymptotic expansion of the Schwartz kernels hold with respect to a convenient $C^{\infty}$ topology. Finally we apply this to complete the proofs of the theorems stated in the introduction.
\subsection{The class \texorpdfstring{$\biginf (k^{-N})$}{O-infinity}}

Consider as before a Hermitian line bundle $L \rightarrow M$ and an auxiliary Hermitian vector bundle $A
\rightarrow M$. Let $\Fam$ be the space of families
\[
s= \left(s_k \in \Ci \left(M, L^k \otimes A\right)
,\, k \in \N\right).
\]
Recall that $s \in \bigo (k^{-N})$ if for any $x \in M$, $|s_k (x)| = \bigo (k^{-N})$ with a $\bigo $ uniform on any compact subsets of $M$. Here we do not assume that $M$ is compact.


The definition of $\biginf$ involves the derivatives. If $L$ and $A$ are trivial bundles so that $(s_k)$ is a sequence of $\Ci (M, \C^r)$ with $r$ the rank of $A$, then we say that $(s_k) \in \biginf (k^{-N})$ if for any $m \in
\N$, the derivatives of order $m$ of $(s_k)$ are in $\bigo (k^{-N+m})$. More precisely, for any vector fields $X_1$, \ldots, $X_m$ of $M$, we require that
\begin{gather}\label{eq:derk}
X_1 \ldots X_m s_k = \bigo \left(k^{-N+m}\right).
\end{gather}
So we loose one power of $k$ for each derivative. Because of this, the class $\biginf (k^{-N})$ is invariant by multiplication by $e^{ik h}$, where $h$ is any real-valued function of $M$.


For actual vector bundles $L$ and $A$, we introduce unitary frames $u$ and $(v_j)_{j=1}^{r}$ of $L$ and $A$ over the same open set $U$ of $M$ and write $s_k = \sum f_{k,j} u^k
\otimes v_j$ with $f_{k} \in \Ci (U, \C^r)$. Then we say that $(s_k)$ belongs to $\biginf (k^{-N})$ if for all choices of unitary frames of $L$ and $A$, the corresponding local representative sequence $(f_k)$ is in $\biginf(k^{-N})$. Observe that changing the frame $u$ of $L$ amounts to multiply $f_k$ by $e^{ikh}$, so the condition that $f_k \in \biginf (k^{-N})$ does not depend on the frame choice when these frames are defined on the same open set.

The typical example of a family in $\biginf (k^{-N})$ is an oscillating sequence
\[
s_k (x) = k^{-N}
e^{-k \varphi (x) } a (x)
\]
with $ \varphi \in \Ci (M)$ having a non negative real part and $a \in \Ci (M, \C^r)$. More generally, for actual bundles, we can set
\[
s_k (x)
= k^{-N} E^k(x) a(x)
\]
where $E \in \Ci (M, L)$ is such that $|E|\leqslant 1$ and $a
\in \Ci (M, A)$.


Obviously, if $N' \geqslant N$, $\biginf (k^{-N'}) \subset \biginf (k^{-N})$. Define $\biginf (k^{-\infty}) := \cap_N \biginf (k^{-N})$. We will need the following result.


\begin{lemm}\label{lem:dev_der_bor}
Let $(s_\ell)$ be a sequence of $\Fam$ such that for any $\ell$, $s_{\ell} \in
\biginf (k^{-p(\ell)}) $ where $(p(\ell))$ is an increasing real sequence, and $p(\ell)
\rightarrow \infty$ as $\ell \rightarrow \infty$. Then
\begin{enumerate}
\item There exists $s \in \Fam \cap \biginf (k^{-p(0)})$ unique modulo $\biginf (k^{-\infty})$ such that
\begin{gather}\label{eq:dev_Der}
s_k = \sum_{\ell =0}^{N-1} s_{\ell,k} + \biginf \left(k^{-p(N)}\right), \qquad
\forall N
\end{gather}
\item Let $s \in \Fam$ such that $s \in \biginf (k^{p})$ for some $p$ and $s_k =
\sum_{\ell =0}^{N-1} s_{\ell,k} + \bigo (k^{-p(N)})$ for any $N$. Then $s
\in \biginf (k^{-p(0)})$ and~\eqref{eq:dev_Der} holds.
\end{enumerate}
\end{lemm}

The first part is a variation of Borel Lemma, the second part follows from interpolation inequalities, cf. as instance~\cite[Lemma~32.]{Shubin}

In the sequel, we will apply this material to Schwartz kernels. So instead of $M$, $L$, $A$, we will have $M^2$, $L \boxtimes \con{L}$ and $B
\boxtimes \con{A}$. 

\subsection{Application to \texorpdfstring{$\lag (A,B)$}{L(A,B)}}

Choose a section~$E$ as in Section~\ref{sec:section-ee} and let $b \in \Ci (M^2, B \boxtimes \con{A})$ vanishing to order $m$ along the diagonal. Then by the same proof as Lemma~\ref{lem:estimate}, the family $ (E^k b)$ is in $\biginf (k^{-\frac{m}{2}})$. Actually, we even have a better result if instead of using any derivatives, we only consider covariant derivatives for the connection of $(L \boxtimes \con{L})^k \otimes (B \boxtimes \con{A})$ induced by the connection of $L$ and any connections of $A$ and $B$.

\begin{lemm}\label{lem:comp_der_lag}
For any $\ell \in \N$, any vector fields $X_1$, \ldots, $X_{\ell}$ of $M^2$, we have $ \nabla _{X_1} \ldots \nabla_{X_\ell} (E^k b) $ is in $\bigo (k^{-\frac{1}{2} (m - \ell)})$.
\end{lemm}

The improvement is that we only loose a half power of $k$ for each derivative.

\begin{proof}
The main observation is that $\nabla E $ vanishes on the
diagonal. Indeed, $\nabla_{X} E =0$ on the diagonal when $X$ is tangent to the first factor because of the second equation in~\eqref{eq:hypotheseE}, but also when $X$ is tangent to the diagonal by the first equation in~\eqref{eq:hypotheseE}. So on a neighborhood of the diagonal, we have  $\nabla_X E = f E$  with $f \in \bigo (1)$. By Leibniz rule, $ \nabla_X (E^k b) = E^k (k f b + \nabla_X b) $. Using this repeatedly, we obtain
\[
\nabla _{X_1} \ldots \nabla_{X_\ell} \left(E^k b\right) = E^k \left(k^{\ell} b_{\ell} +
k^{\ell -1} b_{\ell-1} + \ldots + b_0\right)
\]
where
$ b_{\ell} = \bigo (m+ \ell)$, $b_{\ell-1 } \in \bigo (m+ \ell - 2)$,
\ldots, $b_0 \in \bigo (m - \ell)$. And we conclude as in the proof of Lemma~\ref{lem:estimate}.
\end{proof}

Recall that the Schwartz kernel family of an operator $P \in \lag (A,B)$ has by definition an expansion of the form
\begin{gather}\label{eq:exp_der_part}
P_k (x,y) = \left(\frac{k}{2\pi} \right)^n E^k(x,y) \sum_{ \ell + m (\ell)\,\leqslant\,N} k^{- \frac{\ell}{2}} b_{\ell} (x,y) + R_{N,k} (x,y)
\end{gather}
with $R_{N,k} \in \bigo (k^{n- \frac{N+1}{2}})$.  Let $\lag^{\infty} (A,B)$ (resp. $\lag_q^{\infty} (A,B)$) be the subspace of $\lag (A,B)$ (resp. $\lag_q (A,B)$) consisting of the operator families having a Schwartz kernel in $\biginf (k^n)$.  Identifying operators and their kernels,
\[
\lag^{\infty} (A,B) = \lag (A,B) \cap \biginf (k^n), \quad \lag_q^{\infty} (A,B) = \lag_q (A,B) \cap \biginf (k^n)
\]
By the following proposition, these new classes have the same properties than the $\lag (A,B)$ and this follows directly from Lemma~\ref{lem:dev_der_bor}.

\begin{prop}\label{prop:versionlinfty}\leavevmode
\begin{enumerate}
\item \label{item:1}
If $P \in
\lag^{\infty} (A,B)$ and the expansion~\eqref{eq:exp_der_part} holds with $R_{N,k} \in \bigo (k^{n- \frac{N+1}{2}})$ for any $N$, then $R_{N,k} \in\biginf (k^{n- \frac{N+1}{2}})$ for any $N$.
\item \label{item:2}
For any $P \in \lag (A,B)$ there exist $Q \in
\lag^{\infty} (A,B)$ unique modulo $\biginf (k^{-\infty})$ such that $Q = P +\bigo (k^{-\infty})$.
\item \label{item:3}
 For any $ P \in \lag^{\infty}(A,B)$,
\begin{enumerate}
\item \label{item:a}
the adjoint of $P$ belongs to $\lag^{\infty} (B,A)$,
\item \label{item:b}
$(P_kQ_k) \in \lag ^{\infty} (A,C)$ for any $Q \in \lag ^{\infty} (B,C)$,
\item \label{item:c}
$(f \circ P_k)$ belongs to $\lag^{\infty} (A,C)$ for any $f \in \Ci (M,
\op{Hom} (B,C))$; $(P_k \circ g) $ belongs to $\lag ^{\infty} (C,B)$ for any $g \in \Ci (M, \op{Hom} (C,A))$
\item \label{item:d}
for any vector field $X$ of $M$ and connections on $A$ and $B$, $(k^{-\frac{1}{2}} P_k \circ \nabla^{L^k \otimes A}_X) $ and $(k^{-\frac{1}{2}} \nabla^{L^k \otimes B}_X \circ P_k)$ belong to $\lag^{\infty} (A,B)$.

Furthermore, if $(P_k)$ is even (resp. odd), these two operators are odd (resp. even).
\end{enumerate}

\item \label{item:4}
$\lag_q^{\infty} (A,B) = \lag (A,B) \cap \biginf (k^{n -
\frac{q}{2}})$, the restriction of $\si_q$ to $\lag_q^{\infty} (A,B)$ is onto and has kernel $\lag_{q+1} ^{\infty} (A,B)$.
\end{enumerate}
\end{prop}

\begin{proof}
Assertion~\ref{item:1} follows from the preliminary
observation on $E^kb$ and the second part of Lemma~\ref{lem:dev_der_bor}. Assertion~\ref{item:2} follows from the first part of Lemma~\ref{lem:dev_der_bor}. Claim~\ref{item:b} follows from the fact that the composition of two kernels in $\biginf (k^{n})$ is in $\biginf (k^{2n})$, and $\biginf (k^{2n}) \cap \lag (A,C) = \lag^{\infty} (A,C)$ by the second part of Lemma~\ref{lem:dev_der_bor}. Claims~\ref{item:a} and~\ref{item:c} are straightforward. Ones proves claim~\ref{item:d} by arguing as in the proof of Lemma~\ref{lem:comp_der_lag}. Part~\ref{item:4} follows from the second part of Lemma~\ref{lem:dev_der_bor}.
\end{proof}

\begin{rema} \label{sec:rem_linfty}
We can adapt Theorems~\ref{theo:constr-proj} and~\ref{theo:unitary_equivalence} to the spaces $\lag^{\infty}$:
\begin{enumerate}
\item in Theorem~\ref{theo:constr-proj}, if we start with $P \in \lag^{\infty} (A)$, then $\chi (P) \in \lag ^{\infty} (A)$.
\item in Theorem~\ref{theo:unitary_equivalence}, if $\Pi$ and $\Pi'$ are in $\lag^{\infty} (A)$ and $\lag^{\infty} (B)$ respectively, then we can choose $U \in \lag^{\infty} (A,B)$.
\end{enumerate}
In both cases, the only change in the proof is the fact that for any families of operators $Q_k, Q'_k : \Ci (M,L^k \otimes A) \rightarrow \Ci (M, L^k\otimes A)$ and $Q''_k : L^2 (M, L^k
\otimes A) \rightarrow L^2 (M, L^k \otimes A)$, by~\cite[Section~4.3]{oim}, if the Schwartz kernel families of $(Q_k)$ and $(Q'_k)$ are respectively in $\biginf (k^{-N})$ and $\biginf (k^{-N'})$, and the operator norms of $Q_k^{''}$ are in $\bigo (1)$, then the Schwartz kernel family of $Q_k Q''_k Q_k'$ is in $\biginf (k^{-(N +N')})$.
\end{rema}

By Theorem~\ref{theo:lag}, we already know how to compute the symbols of $P^*$, $PQ$, $f P$ or $gP$ in terms of the symbols of $P$ and $Q$. To complete this, we compute the symbol of the compositions of $P$ with the covariant derivatives $\nabla_X^{L^k \otimes A}$ and $\nabla_X^{L^k \otimes B}$. Recall that for any $Y \in T_xM$, we defined in the introduction some endomorphisms $\rho (Y) \in \op{End} (\D (T_xM))$ in~\eqref{eq:rho}. If $Y = U + \con V$ with $U, V \in T_x ^{1,0} M$, then $\rho (Y)=
\rho (U) + \rho (\con V)$ where $\rho(U)$ is the multiplication by $i \om (U, \cdot)$ and $\rho (\con{V})$ is the derivation with respect to $\con{V}$.


\begin{lemm}\label{lem:symb_der}
For any $P \in \lag^{\infty}(A,B)$ and vector field $X$ of $M$, we have
\begin{gather*}
\si_0 \left(k^{-\frac{1}{2}} P_k \circ \nabla^{L^k
\otimes A}_X\right) (x) = \si_0 (P_k) (x) \circ \rho (X(x)) \\
 \si_0 \left(k^{-\frac{1}{2}} \nabla^{L^k
\otimes B}_X \circ P_k\right) (x) = \rho (X(x)) \circ \si_0 (P_k) (x). 
\end{gather*}
\end{lemm}

\begin{proof}
We deduce one formula from the other by taking adjoint. To
prove the first one, it suffices by Proposition~\ref{prop:peaked-sections_++} to show that if $(\Phi_k^f)$ is the peaked section associated to $ f\in \D(T_xM) \otimes A_x$, then
\[
k^{-\frac{1}{2}}
\nabla_X \Phi_k^f = \Phi_k^g + \bigo \left(k^{-\frac{1}{2}}\right)
\]
where $g = \rho (X(x)) f $. This is easily checked if we use the normal coordinates as in~\eqref{eq:coordonnees_z} centered at $p_0 = x$. We have to differentiate~\eqref{eq:phi_peaked_section_+}. We have first that $k^{-\frac{1}{2}} \nabla_X (E^k) = k^{\frac{1}{2}} E^k (\nabla_X
E)E^{-1}$ and by~\eqref{eq:E_linearise},
\[
(\nabla_{\partial_i} E)E^{-1} =
- \con{z}_i + \bigo (2), \qquad (\nabla_{\con{\partial}_i} E)E^{-1} =
\bigo (2).
\]
Second we have $ k^{-\frac{1}{2}} \partial_i f (k^{\frac{1}{2}} \xi) = 0$ since $f \in \C [\con{z}_1, \,\ldots,\, \con{z}_n]$ and $ k^{-\frac{1}{2}} \con{\partial}_i f (k^{\frac{1}{2}} \xi) = (\partial f / \partial \con{z}_i) (k^{\frac{1}{2}} \xi)$. To conclude recall that $\rho (\partial_i) $ is the multiplication by $-\con{z}_i$ whereas $\rho (\con{\partial}_i)$ is the derivation with respect to $\con{z}_i$.
\end{proof}

\subsection{Kostant--Souriau operators and subprincipal estimates}

In our context, the Kostant--Souriau operators are the operators of the form
\[
f + \tfrac{i}{k} \nabla_X^{L^k \otimes A} : \Ci \left(M, L^k \otimes A\right)
\rightarrow \Ci \left(M, L^k \otimes A\right)
\]
where $ f \in \Ci (M)$ and $X$ is its
Hamiltonian vector field, that is $\om (X, \cdot) + df =0$.

\begin{lemm}\label{lem:ledernier}
For any $f \in \Ci (M, \R)$ with Hamiltonian vector field $X$, for any $P \in \lag^{\infty}_0 (A)$ we have
\begin{itemize}
\item $ [f, P]$ belongs to $\lag_{1}^{\infty} (A)$
\item $[f,P] \equiv (ik)^{-1} [\nabla_X^{L^k
\otimes A}, P]$ modulo $\lag_{2}^{\infty} (A)$.
\end{itemize}
So the commutator $[f +\frac{i}{k} \nabla_X^{L^k \otimes A}, P]$ belongs to $\lag_{2}^{\infty}(A)$.
\end{lemm}

\begin{proof}
The Schwartz kernel of $[f, P]$ is the product of $g(x,y)= f(x) -
f(y)$ by the Schwartz kernel of $P$. Since $g$ vanishes along the diagonal, this implies that $[f, P] \in \lag_1 (A)$. Furthermore, it follows from the definition~\eqref{eq:def_si_tilde} of the symbol that if $\wt{\si}_0 (P)(x) = \op{Op} (b)$ with $b \in \pol (T_xM)$, then $\wt{\si}_1 ([f, P])(x) = \op{Op} (\ell b)$ where $\ell = d_x f \in T_x^*M$. Now a computation from~\eqref{eq:sch_ker} shows that
\[
\left[\op{Op} (b), a_i\right] = \op{Op} (z_i b), \qquad \left[\op{Op} (b), a_i^*\right] = \op{Op} (\con{z}_i b)
\]
where as in Section~\ref{sec:landau-levels-cn} we use the annihilation and creation operators $a_i = \partial_{\con{z}_i}$, $a_i^* = \con{z}_i -
\partial_{z_i}$.  So working with normal coordinates at $p_0 =x$ as in~\eqref{eq:coordonnees_z},
\[
\wt{\si}_1 ([f, P])(x) = \sum_i (\partial_{z_i} f)(x) [\op{Op} (b), a_i] + (\partial_{\con{z}_i} f)(x) [\op{Op} (b), a_i^*]
\]
Moreover $ i X = (\partial_i f) \con{\partial}_i - (\con{\partial}_if)
\partial_i $ at $x$. Since $ \rho (\con{\partial}_i)$ and $\rho (\partial_i)$ are the restrictions of $a_i$ and $-a_i^*$ to $\C[\con{z}_1,\,
\ldots,\,\con{z}_n]$ respectively, we deduce that
\[
\si_1 ([f, P]) = i [\si_0 (P), \rho(X)].
\]
On the other hand, by Lemma~\ref{lem:symb_der}, $\si_0 (\bigl[k^{-\frac{1}{2}}
\nabla_X^{L^k \otimes A}, P]) = [\rho(X), \si_0 (P)]$. So
\[
\si_1 \left(\left[f + \tfrac{i}{k} \nabla_X^{L^k \otimes A}, P\right]\right) =0
\]
and the result follows. 
\end{proof}

Let us apply this result to the Toeplitz algebra associated to a self-adjoint projector $\Pi \in \lag^{\infty}(A) \cap \lag^{+}(A)$, whose symbol $\pi$ is the projector on a subbundle $F$ of $\D (TM) \otimes A$ having a definite parity. Introduce the operators associated to $f \in \Ci (M)$, $X, Y \in
\Ci (M, TM)$,
\[
T_k (f) = \Pi_k f \Pi_k, \qquad T_k (X, Y) = k^{-1} \Pi_k \nabla^{L^k
\otimes A}_X \nabla^{L^k
\otimes A}_Y \Pi_k.
\]
By Proposition~\ref{prop:versionlinfty}, $T_k (X, Y)$ belongs to $\lag ^+(A)$ so it belongs to the Toeplitz algebra~\eqref{eq:Toeplitz_Pi} associated to $\Pi$. By Lemma~\ref{lem:symb_der}, its symbol is $\pi \rho (X) \rho (Y)
\pi$.

\begin{prop}\label{prop:subsymbolic}
For any $f,g \in \Ci (M)$,
\begin{gather}\label{eq:lastproofeq}
T_{k} (f) T_k (g) = T_k (fg) + k^{-1} T_k (X,Y) + \bigo \left(k^{-2}\right)
\end{gather}
where $X$ and $Y$ are the Hamiltonian vector fields of $f$ and $g$ respectively. Consequently $ i [T_k (f), T_k (g)] = k^{-1} T_k (\om (X,Y)) + \bigo (k^{-2})$.
\end{prop}
\begin{proof}
By a straightforward computation, we have
\[
\Pi f \Pi g \Pi = \Pi fg \Pi + \Pi [f, \Pi]\,[g, \Pi] \Pi
\]
By lemma~\ref{lem:ledernier}, $\Pi [f, \Pi] [g, \Pi] \Pi$ belongs to $\lag_2 (A)$ and its symbol is
\begin{gather*}
\si_2 \big(\Pi [f, \Pi]\,[g, \Pi] \Pi\big) = - \pi [\rho (X), \pi]\,[\rho (Y),
\pi] \pi
\end{gather*}
Observe that $\rho (X) |_x$ is an odd operator of $\D (T_xM)$. Since $F_x$ has a definite parity, every endomorphism of $F_x$ is even. So $\pi \rho (X)
\pi |_x$ is at the same time odd and even, so $\pi \rho (X) \pi =0$.  Similarly $\pi \rho (Y) \pi =0$, so
\begin{gather*}
\si_2 \big(\Pi [f, \Pi]\,[g, \Pi] \Pi\big) = \pi \rho (X) \rho (Y) \pi = \si_2 \left(k^{-2} \Pi \nabla_X^{L^k \otimes A} \nabla_Y^{L^k \otimes A} \Pi\right)
\end{gather*}
which proves~\eqref{eq:lastproofeq}. Consequently, the rescaled commutator of $T_k(f)$, $T_k(g)$ satisfies
\[
k \left[T_k (f), T_k (g)\right] = T_k (X, Y) - T_k (Y,X) + \bigo \left(k^{-1}\right),
\]
so its symbol is $\pi ([\rho (X), \rho (Y)] \pi$. To conclude the proof, we simply use that
\begin{gather}\label{eq:commutator_rho}
\left[\rho (X), \rho (Y)\right] = \tfrac{1}{i} \om (X,Y),
\end{gather}
as follows easily by using a basis $U_i$ of $T_x^{1,0}M$ such that $\frac{1}{i} \om (U_i, \con{U}_j) = \delta_{ij}$, $\rho (U_i) = - a_i^*$ and $\rho (\con{U}_i) = a_i$.
\end{proof}


\subsection{Proofs of Theorems~\ref{theo:dim_landau}, \ref{theo:Toeplitz_Landau} and~\ref{theo:ladder}} \label{sec:ledernier}

In this last section, we complete the proof of the theorems stated in the introduction, a generalization actually since we will consider more general projectors.

Let $\Pi \in \lag^{\infty}(A) \cap \lag^{+} (A)$ be a self-adjoint projector with symbol $\pi= \pi_m \otimes \op{id}_A$, where $\pi_m$ is the projector of $\D (TM)$ onto $\D_m (TM)$. Such an operator exists by Theorem~\ref{theo:constr-proj} and Remark~\ref{sec:rem_linfty}. Alternatively, the projector $\Pi = (\Pi_{m,k})$ onto the $m^{\rm th}$ Landau level defined in~\eqref{eq:projecteur_fibre_auxiliaire} has the expected properties~\cite[Theorems~5.2, 5.3]{oim_copain}.

By Theorem~\ref{theo:dim_general}, the dimension of $\Hilb_k= \op{Im}(\Pi_k)$ is
\[
\op{dim} \Hilb_k = \int_M \op{ch} \left(L^k \otimes A \otimes \D_m(TM)\right) \; \op{Td}
(M)
\]
when $k$ is sufficiently large, which implies Theorem~\ref{theo:dim_landau}.

Define the Toeplitz algebra
\[
\clT^{\infty} = \left\{ P \in \lag^{\infty}(A) \cap \lag^{+}
(A)/\; \Pi P \Pi = P \right\}.
\]
Clearly $\clT^{\infty}$ is contained in the Toeplitz algebra $\clT$ defined in~\eqref{eq:Toeplitz_Pi}, and by assertion 2 of Proposition~\ref{prop:versionlinfty}, the difference is rather small: for every $ P \in
\clT$, there exists $P' \in \clT^{\infty}$ unique modulo $\biginf (k^{-\infty})$ such that $P' = P + \bigo (k^{-\infty})$. Recall the symbol map $\tau_0$ introduced in Theorem~\ref{theo:toeplitz} and denote by $\tau$ its restriction to $\clT^{\infty}$
\[
\tau : \clT^{\infty} \rightarrow \Ci (M, \op{End} (\D_m (TM) \otimes A))
.
\]
We could as well consider the maps $\tau_q$ with $q \geqslant 1$ but we will limit ourselves to $\tau_0$. By Theorem~\ref{theo:toeplitz} and Assertions 2, 4 of Proposition~\ref{prop:versionlinfty}, $\tau$ is onto and its kernel is $k^{-1}
\clT^{\infty}$. It follows as well from Theorem~\ref{theo:toeplitz} that for any $P,Q \in \clT^{\infty}$
\begin{gather*}
\tau (PQ) = \tau (P) \tau (Q), \quad \| P_k \| = \sup_{x\, \in\,M} \|
\tau(P)_x\| + \op{o} (1) \\
P_k (x,x) = \left(\frac{k}{2\pi} \right)^n \op{tr} (\tau (P)_x) + \bigo \left(k^{-1}\right)
\end{gather*}
Choose any connection on $A$. By Proposition~\ref{sec:rem_linfty}, for any $ f \in \Ci (M, \op{End} A)$, $p \in \N$ and vector fields $X_1$, \ldots, $X_{2p}$ of $M$, the operator
\begin{gather}\label{eq:toeplitz_mult_der}
T_k \left(f, X_1,\, \ldots,\, X_{2p}\right) = k^{-p} \Pi_k f \nabla_{X_1}^{L^k \otimes A}
\ldots \nabla_{X_{2p}}^{L^k \otimes A}
\Pi_k
\end{gather}
belong to $\clT^{\infty}$ and  by Lemma~\ref{lem:symb_der}  its $\tau$-symbol is $ (\pi_m \rho(X_1) \ldots
\rho(X_{2p})\pi_m) \otimes f $. By Lemma~\ref{lem:algebraic}, these symbols span $\Ci (M,
\op{End} (\D_m (TM) \otimes A))$ as a vector space, we deduce that any $P \in \clT^{\infty}$ is of the form
\[
P_k = \sum_{\ell =0 }^N k^{-\ell} P_{\ell,k} + \bigo \left(k^{-(N+1)}\right), \qquad
\forall N \in \N
\]
where for any $\ell$, $(P_{\ell,k})_k$ is a finite sum of operators such as~\eqref{eq:toeplitz_mult_der}. So in the case where the auxiliary bundle $A$ is trivial, $\clT^{\infty}$ is the space $\clT^{\op{sc}}_m$ defined in the introduction. Last assertion of Theorem~\ref{theo:Toeplitz_Landau} follows from Proposition~\ref{prop:subsymbolic}.

\begin{lemm}\label{lem:algebraic}\leavevmode
\begin{enumerate}
\item\label{lemm6.8.1} $\op{End} \D_m(T_xM)$ is spanned as a vector space by the
\[
\pi_m (x) \rho (X_1) \ldots \rho (X_{2p}) \pi_m (x)
\]
where $p \in \N$ and $X_1$,
\ldots, $X_{2p} \in T_xM$.
\item\label{lemm6.8.2} The linear map $\Psi_x : T^{(1,0)} _xM \otimes T^{(0,1)} _xM \rightarrow \op{End} \D_m (T_xM)$ such that $\Psi (U \otimes
\con{V})= \pi_m (x) \rho (U)
\rho (\con{V}) \pi_m (x)$, is injective when $m \geqslant 1$.
\end{enumerate}
\end{lemm}

\begin{proof}
Introduce a
basis $(U_i)$ of $T^{1,0}_xM$ such that $\frac{1}{i} \om (U_i, \con {U}_j) = \delta_{ij}$ and let $z_i = \frac{1}{i} \om (\cdot,\con{U}_i)$. So $ \D (T_x M) = \C [\con{z}_1, \ldots, \con{z}_n]$, $\rho (U_i) P = - \con{z}_i P = - a_i^* P$ and $ \rho (\con{U}_i) P = \partial P / \partial
\con{z}_i = a_i P$. For any $\al $, $\be$ in $\N^n$, $f_{\al, \be}= (\al!)^{-1} (a^*)^{\be} a^\al$ satisfies $f_{\al,\be} (\con{z}^\al) =
\con{z}^\be$ and for any $\ga \in \N^n$ such that $\ga \neq \al$ and $| \ga | = |
\al |$, $f_{\al,\be} (\con{z}^\ga) = 0$. So the family $\pi_m(x) f_{\al, \be}
\pi_m (x)$, $| \al | = | \be | =m$ is a basis of $\op{End} F_x$. This proves the first assertion. The second one is the fact that the restrictions to $\D_m (T_xM)$ of the $ \con{z}_i \con{\partial}_j$, $i, j =1,\, \ldots,\, n$ are linearly independent.
\end{proof}


Equation~\eqref{eq:bm_1-f-g} and Part~\eqref{rema1.4.2} of Remark~\ref{rem:Toeplitz} are consequences on the following Proposition.

\begin{prop}\label{prop:fonction}
For any $f \in \Ci (M)$, with Hamiltonian vector field $X$, we have
\begin{itemize}
\item if $n =1$ or $m=0$, $ T_k (f)^2 - T_k (f^2) = -k^{-1} (\frac{1}{2} + m)T_k (|X|^2) + \bigo (k^{-2}).$
\item if $n \geqslant 2$ and $m \geqslant 1$, $ X \neq 0$ implies that there exists no function $h \in \Ci (M)$ such that $T_k (f)^2 - T_k (f^2) = k^{-1} T_k (h) + \bigo (k^{-2})$.
\end{itemize}
\end{prop}

\begin{proof}
By Proposition~\ref{prop:subsymbolic}, $ T_k (f)^2 = T_k (f^2) + k^{-1} S_k $ where $S_k$ is a Toeplitz operator with symbol $\pi \rho(X)^2 \pi$. Writing $X = U + \con {U}$ with $U \in T^{1,0}_xM$, we have
\[
\pi_m \rho (X)^2 \pi_m = \pi_m \bigl(\rho (U) \rho (\con{U}) + \rho
(\con {U}) \rho (U) \bigr) \pi_m = 2 \Psi (U \otimes \con{U}) - \tfrac{1}{2} g
(X,X)
\]
where $\Psi$ is the bundle map introduced in Lemma~\ref{lem:algebraic}, and we have used first that $\pi_m \rho (U) ^2 \pi_m =0 =\pi_m
\rho(\con{U})^2 \pi_m $ and then that $[\rho(\con{U}), \rho (U)] = i \om (U, \con {U}) = -\frac{1}{2} g (X,X)$ by~\eqref{eq:commutator_rho}. Now if $m=0$, then $\Psi =0$. If $n=1$, then $2 \Psi (U \otimes \con{U}) = - m g (X,X)$. This proves the first assertion. For the second one, it suffices to prove that if $n \geqslant 2$, $m \geqslant 1$ and $X(x) \neq 0$, then $\Psi (U \otimes \con{U})_x \in \op{End}(\D_m (T_xM))$ is not scalar. This follows from the fact that $\Psi_x$ is injective, so that $\Psi_x (\al) $ is scalar only when $ \al$ is a multiple of $\sum_{i=1}^{n} U_i \otimes
\con{U}_i$, which never happens for $\al = U \otimes \con{U}$ when $n
\geqslant 2$ and $U \neq 0$.
\end{proof}


Let us prove now Theorem~\ref{theo:ladder}. Introduce a quantization $(\Hilb_{F,k})$ of $(M,L)$ twisted by $F = \D_m (TM) \otimes A$. We can adapt the definition~\eqref{eq:def_w_k} of $W_k$ with the auxiliary bundle $A$ by setting
\begin{equation}\label{eq:def_w_k_A}
\begin{split}
&W_k : \Ci \left(M, L^k \otimes A\right) \rightarrow \Ci \left(M, L^k \otimes F\right), \qquad k \in \N\\
&W_k = R_m D_{G^{\otimes (m-1)} \otimes A,k} \circ D_{G^{\otimes (m-2)}\otimes A,k} \circ \ldots \circ D_{G \otimes A, k} \circ D_{A,k}
\end{split}
\end{equation}

\begin{lemm}
The operator $(V_k = \frac{1}{m!} k^{-\frac{m}{2}} \Pi_{F,k} W_k \Pi_{k}$, $k \in \N$) belongs to $ \lag ^{\infty} (A, F)$, has the same parity as $m$ and its symbol $\si_0(V)$ viewed as a morphism from $\D(TM) \otimes A$ to $\D(TM) \otimes F$ is given by
\[
\forall \, f \in \D_p(TM), \; \forall \, a \in A, \; \si_0(V) (f \otimes a)
=
\begin{cases} 1 \otimes f \otimes a \text{ if $p = m $} \\ 0 \text{
otherwise}
\end{cases}
\]
\end{lemm}

\begin{proof}
We claim that for any even (resp. odd) operator 
$P \in \lag^{\infty} (B, A)$ , \newline $(k^{-\frac{1}{2}}D_{A,k} \circ P)$ belongs to $\lag^{\infty} (B, A
\otimes G) $, is odd (resp. even) and its symbol is $\varphi_A \circ \si_0 (P)$ where
\[
\varphi_A = \sum_ i a_i \otimes \con{z}_i \otimes
\op{id}_A \in \symb (TM) \otimes G \otimes \op{End} A.
\]
Here we have introduced an orthonormal frame $(
\partial_i)$ of $T^{1,0}M$, $(z_i)$ is the dual frame of $(T^{1,0}M)^*$ and $a_i =\partial_{\con{z}_i}$ is the annihilation operator.  This follows from Lemma~\ref{lem:symb_der} by writing $D_{A,k} s = \sum_i \con{z}_i \otimes
\nabla_{\con{\partial}_i} s$. 

Consequently, $k^{-\frac{m}{2}}R_m D_{A \otimes G^{m-1}} \circ \ldots \circ D_{A,k} \circ \Pi_k $ belong to $\lag (A, F)$ with symbol $\varphi_{A}^m \circ \pi_m$ where $\varphi_A^m $ is the morphism from $\D (TM) \otimes A$ to $\D (TM) \otimes \D_m (TM) \otimes A$ given by
\[
\varphi_A^m = \sum_{i_1, \,\ldots,\, i_m =1}^{n} a_{i_1} \ldots a_{i_m} \otimes
(\con{z}_{i_1} \ldots \con{z}_{i_m}) \otimes \op{id}_A
\]
Computing $\varphi_A^m (\con{z}^{\be})$ for any multi-index $\be$, we show that for any $f \in \D_p(TM)$ and $ a\in A$, $ \varphi_A^m (f
\otimes a) = m ! (1 \otimes f \otimes a)$ when $p=m$ and $0$ otherwise.
\end{proof}

The symbol $\sigma_0 (V)$ is exactly the symbol $\rho$ introduced in Lemma~\ref{lem:unitary-symbole}. Now Theorem~\ref{theo:ladder} and the remark on the unitarization of $V_k$ follows by the same proof as Theorem~\ref{theo:unitary_equivalence}.

\appendix
\section{} \label{sec:appendix}

In this appendix we discuss the known results for surfaces with constant curvature. Two important features appear for the negatively curved surfaces: only the lower part of the spectrum consists of Landau levels, and moreover, there is an isomorphism between the $m^{\rm th}$ Landau level and the first level of a Laplacian twisted by the $m^{\rm th}$ power of the complex determinant bundle.

These results appeared in the physics literature, cf. in particular~\cite{IeLi94} for the case with surface with genus $\geqslant 2$. A more recent mathematical reference is~\cite{Te06}.

\subsection*{The plane~\cite{La30}}
Consider a quantum particle confined in a two dimensional plane $(x,y)$ and subject to a constant magnetic field perpendicular to this plane. Its Hamiltonian is the operator
\begin{gather}\label{eq:laplacian_landau}
H = -\tfrac{1}{2} \left(\nabla_x^2 + \nabla_y^2\right) \quad \text{ with } \quad \nabla_x = \tfrac{\partial}{ \partial x} +
\tfrac{i}{2} B y,\; \nabla_y = \tfrac{\partial}{\partial y} - \tfrac{i}{2} Bx.
\end{gather}
$B$ is a positive constant representing the strength of the magnetic field. The spectrum of $H$ is $B(\frac{1}{2} + \N)$ and the Landau levels $\Hilb_m =
\op{ker} (H - B(\frac{1}{2} + m))$ are given in terms of the ladder operators $\nabla_z = \nabla_x - i \nabla_y$, $\nabla_{\con{z}} = \nabla_x + i
\nabla_y$ by
\begin{gather}\label{eq:landau_level}
\Hilb_0 = \op{ker} (\nabla_{\con z}), \qquad \Hilb_m = (\nabla_z)^m
\Hilb_0, \quad m \geqslant 1.
\end{gather}

\subsection*{Surface with constant curvature~\cite{IeLi94, Te06}}
Let $M$ be a compact orientable surface with a Riemannian metric having a constant Gauss curvature $S$. Introduce a Hermitian line bundle $ L
\rightarrow M$ with a connection $\nabla : \Ci (M, L) \rightarrow \Om^1 (M,L)$. Assume that the curvature satisfies
\begin{gather}\label{eq:constant_magnetic}
i \op{curv} (\nabla) = B
\op{vol}_g
\end{gather}
where $B$ is a non-zero constant and $\op{vol}_g$ is the Riemannian volume. Choosing the convenient orientation for $M$, we can assume that $B$ is positive. The quantum Hamiltonian is the Laplacian $ \Delta := \frac{1}{2} \nabla^* \nabla $ acting on sections of $L$. Then denoting its eigenvalue by $0 \leqslant \la_0 < \la_1 <
\ldots $, it is known that
\begin{gather}\label{eq:spec_magentic_surface}
\la_m = B \left(\tfrac{1}{2} + m\right) + S \tfrac{ m (m+1)}{2} \qquad \text{ if } B + m S >0
\end{gather}
For a sphere or a torus, $S \geqslant 0$, and these formulas describe the whole spectrum. If the genus of $M$ is larger than $2$, then $S<0$ and the condition $B + m S >0$ is satisfied only for a finite number of $m$. In this case, it is not reasonable to expect an explicit formula for the other eigenvalues. Indeed, if $S=-1$ and $L = K^r$ with $K$ the canonical bundle of $M$ and $r$ a positive integer, then~\eqref{eq:spec_magentic_surface} gives the first $(r+1)$ eigenvalues, and for any $n \in \N$, $\la_{r+n} = \la_r + \tfrac{1}{2} \mu_n$, where $\{\mu_n, \; n
\in \N\}$ is the spectrum of the Laplace--Beltrami operator of $M$. This latter spectrum depends in an essential way on the metric. Indeed, by Huber's theorem~\cite[Theorem~9.2.9]{Buser}, $\{ \mu_n\}$ determines the length spectrum of $M$.

The multiplicity of the first eigenvalues is equal to:
\begin{gather}\label{eq:mult}
\op{mult} (\la_m) = B \frac{ \op{vol}_g (M)}{2\pi} + \left(\tfrac{1}{2} + m\right)
\chi (M) \qquad \text{ if } B + (m+1)S >0.
\end{gather}
Here $\chi (M)$ is the Euler characteristic of $M$ and observe that $B
\op{vol}_g (M) / (2\pi)$ is the degree of $L$, so it is an integer. \eqref{eq:mult} follows from a description of the corresponding eigenspace $\Hilb_m = \ker (\Delta -
\la_m)$ similar to~\eqref{eq:landau_level}, that we explain below.

\subsection*{Proof of formulas~\eqref{eq:spec_magentic_surface} and~\eqref{eq:mult}}

To start with, we do not assume that the Gauss curvature $S$ and the function $B$ defined in~\eqref{eq:constant_magnetic} are constant. We choose any orientation on $M$. Let $j$ be the complex structure of $M$ compatible with $g$, i.e. $g(jX, jY) = g(X,Y)$ for any tangent vectors $X$, $Y \in T_pM$ and $ X
\wedge jX >0$ if $X \neq 0$. Since we are in real dimension~2, $j$ is integrable. Furthermore the associated volume form $\op{vol}_g$ is the symplectic form $\om (X, Y) = g(jX, Y)$.

$L$ has a natural holomorphic structure such that its $\con{\partial}$-operator is $\nabla^{0,1}$. We denote it by $\con{\partial}_L: \Ci (L) \rightarrow \Ci (L \otimes \con{K})$ with $K = (T^*M)^{1,0}$ the canonical bundle. The curvature of $\nabla$ has the form $\frac{1}{i} B \om$ with $B \in \Ci (M, \R)$. The canonical bundle has a natural metric induced by $g$. Its Chern connection, that is its connection compatible with both its metric and holomorphic structure, has curvature $i S \om$ where $S$ is the Gauss curvature.

\begin{theo}\label{theo:appendix}
The following identities holds:
\begin{enumerate}
\item\label{theo7.1.1} Weitzenb\"ock formula: $\Delta_L = \con{\partial}^*_L \con{\partial}_L + \frac{1}{2} B$.
\item\label{theo7.1.2} Bosonic commutation relation: $ \con{\partial}_L \con{\partial}_L^* = \con{\partial}_{L\otimes K^{-1}} ^* \con{\partial}_{L\otimes K^{-1}} + (B + S)$.
\end{enumerate}
\end{theo}

The Weitzenb\"ock formula is a classical relation, it holds more generally on K\"ahler manifolds. We call the second formula the bosonic commutation relation because it replaces the canonical commutation relation satisfied by the creation/annihilation operators $[a, a^{*}] =1$. In this formula, we identify $\con{K}$ with $K^{-1}$ through the metric so that the operators $ \con{\partial}_L \con{\partial}_L^*$ and $\con{\partial}_{L\otimes K^{-1}} ^* \con{\partial}_{L\otimes K^{-1}}$ act on the same space $\Ci (L\otimes \con{K}) = \Ci (L
\otimes K^{-1})$. A similar formula were obtained in~\cite[Proposition~9]{Te06} for the same purpose of computing the spectrum of $\Delta_L$.


\begin{proof}[Proof of the bosonic identity]
Introduce a local holomorphic frame $s$ of $L$ and a complex coordinate $z$ on $M$. We have first $\con{\partial}_L (f s) = f_{\con{z}} \; s \otimes d \con{z}$. To compute the adjoint, recall that the scalar products of $\Ci (L)$ and $\Ci(L \otimes \con{K})$ are defined by integrating the pointwise scalar products against the volume form. Write $\om = i h dz \wedge d \con{z}$ and $|s|^2 = e^{-
\varphi}$ with $h$ and $\varphi$ real valued functions. Then $|d \con z |^2 = h^{-1}$ and a direct computation leads to
\[
\con{\partial}_L^* (f s \otimes d\con{z}) = h^{-1} (-
f_z + f \varphi_z) s.
\]
With the identification $\con{K} \simeq K^{-1}$, we have $ d\con{z} = h^{-1} (dz)^{-1}$. We deduce that
\[
\con{\partial}_L \con {\partial}_L^* \left(f s \otimes (dz)^{-1}\right) = h^{-1} \Big(-
f_{z \con{z} } + f_{\con{z}} \left(\varphi_z - \tfrac{h_z}{h}\right) + f \left(\varphi_{z
\con{z}} - \partial_{\con{z}} \left(\tfrac{h_z}{h}\right)\right) \Big) s \otimes (dz)^{-1}
\]
Similar computations by using $|s \otimes dz^{-1} |^2 = h e^{-\varphi}$ leads to
\[
\con{\partial}_{L\otimes
K^{-1}} ^* \con{\partial}_{L\otimes K^{-1}} \left(f s \otimes (dz)^{-1}\right) = h^{-1} \left(- f_{z \con{z}} +
f_{\con{z}} \left(\varphi_z - \tfrac{h_z}{h}\right) \right) s \otimes (dz)^{-1}.
\]
To conclude, observe that $B = h^{-1} \varphi_{z \con z}$ and $S = - h^{-1}
\partial_z \partial_{\con z} \, \ln h $.
\end{proof}

From now on, we will assume that $B$ and $S$ are constant. With the Weitzenb\"ock formula, we pass directly from the spectrum of $\Delta _L$ to the one of $\con{\partial}^*_L \con{\partial}_L$. We can use the Bosonic relation exactly as it is usually done with the Landau Hamiltonian, cf. proof of Proposition~\ref{prop:linearlandau}. We deduce that for any $\la \neq 0$, $\la $ is an eigenvalue of $\con{\partial}^*_L
\con{\partial}_L$ if and only if $\la - (B+S)$ is an eigenvalue of $\con{\partial}^*_{L\otimes K^{-1}} \con{\partial}_{L\otimes K^{-1}}$. Moreover the eigenspaces have the same dimension. Indeed $\con{\partial}_L$ restricts to an isomorphism
\[
\op{Ker} \left(\la - \con{\partial}^*_L
\con{\partial}_L\right) \rightarrow \op{ker} \left(\la - (B+S) -
\con{\partial}^*_{L\otimes K^{-1}} \con{\partial}_{L\otimes K^{-1}}\right)
\]
with inverse the restriction of $\la^{-1} \con {\partial}_L^*$. Besides this, $\ker (\con{\partial}^*_L
\con{\partial}_L)$ is the space $H^0 (L)$ of holomorphic sections of $L$. By Riemann--Roch theorem, $H^0 (L)$ has dimension $d + \frac{1}{2} \chi (M)$ if the degree $d = B \op{Vol}(M)/(2\pi)$ of $L$ is larger than $- \chi (M)$.

To summarize, when $B$ is sufficiently large, $0 $ is an eigenvalue of $\con{\partial}^*_L \con{\partial}_L$ with multiplicity equal to $ B
\op{Vol}(M) / (2 \pi) + \frac{1}{2} \chi(M)$, and the remainder of the spectrum is identical with the spectrum of $(B+S) + \con{\partial}^*_{L\otimes K^{-1}}
\con{\partial}_{L\otimes K^{-1}}$, multiplicities included. We can iterate this argument and deduce by induction the formulas~\eqref{eq:spec_magentic_surface}, \eqref{eq:mult} giving the first eigenvalues of $\Delta_{L}$ with their multiplicity. Since $\op{deg} (L \otimes K^{-1}) = \op{deg} (L) + \chi (M)$, we can repeat ad infinitum this argument when $\chi (M) \geqslant 0$ and obtain the whole spectrum of $\Delta_L$; whereas for $\chi (M) <0$, only a finite number of iterations is possible.

